% Latin-script Interslavic source component: Work 44 (Deuring lecture notes).
% Audited editorial provenance: ../00_SESSION_INDEPENDENT_STATE_v019.json and ../NORMALIZATION_DECISIONS_v019.jsonl.
\documentclass[11pt]{article}
\usepackage{fontspec}
\IfFontExistsTF{Times New Roman}{\setmainfont{Times New Roman}}{\setmainfont{TeX Gyre Termes}}
\usepackage{amsmath,amssymb}
\usepackage{geometry}
\usepackage{microtype}
\usepackage{hyperref}
\usepackage{mathrsfs}
\newcommand{\srcfn}[2]{\footnote{#2}}
\providecommand{\tightlist}{\setlength{\itemsep}{0pt}\setlength{\parskip}{0pt}}
\allowdisplaybreaks

\geometry{a4paper,margin=2.1cm}
\setlength{\parindent}{1.1em}
\setlength{\parskip}{0pt}
\emergencystretch=10em
\sloppy
\hbadness=10000
\raggedbottom
\hypersetup{hidelinks}
\newcommand{\foreign}[1]{#1}
\newcommand{\tocline}[2]{\noindent #1\dotfill #2\par}
\newcommand{\tocsec}[3]{\noindent\hspace*{1.2em}\S\,#1.\quad #2\dotfill #3\par}

\begin{document}
\thispagestyle{empty}

\begin{center}
{\LARGE\bfseries Algebra hiperkompleksnyh veličin\par}
\vspace{3.5em}
{\large Lekcija prof. E. \foreign{Noether}\par}
{\large Zimny semestr 1929/30\par}
\vspace{1.5em}
{\large Obrabotal prof. M. \foreign{Deuring}\par}
\vspace{4em}
{\Large\bfseries Sodržanje\par}
\end{center}
\vspace{1em}

\begingroup
\small
\setlength{\parindent}{0pt}
\tocline{\textbf{Uvod}}{3}
\medskip
\tocline{\textbf{Glava I: Prědstavjenja}}{3}
\tocsec{1}{Definicija direktnogo prědstavjenja --- definicija recipročnogo prědstavjenja}{3}
\tocsec{2}{Klasy prědstavjenij}{4}
\tocsec{3}{Moduly prědstavjenij}{4}
\tocsec{4}{Svęz medžu modulami prědstavjenij i prědstavjenjami}{5}
\medskip
\tocline{\textbf{Glava II: Galoisova teorija komutativnyh polj}}{7}
\tocsec{5}{Razširjenje koeficientnogo polja hiperkompleksnyh sistemov}{7}
\tocsec{6}{Ireducibilne prědstavjenja komutativnyh sistemov}{8}
\tocsec{7}{Izomorfizmy polja}{8}
\tocsec{8}{Razpadno polje}{9}
\tocsec{9}{Izomorfizmy od $Z_1$ na $Z_i$}{9}
\tocsec{10}{Galoisova grupa}{9}
\tocsec{11}{Glavny teorem Galoisovoj teorije}{9}
\tocsec{12}{Formalno značenje komponentov $e_i$}{11}
\tocsec{13}{Obći teorem o prodolženju}{12}
\medskip
\tocline{\textbf{Glava III: Abelove grupy}}{13}
\tocsec{14}{Grupovo prsten}{13}
\tocsec{15}{Grupove prsteny abelovyh grup}{14}
\tocsec{16}{Relacije karakterov}{15}
\tocsec{17}{Galoisova teorija abelovyh grup}{15}
\medskip
\tocline{\textbf{Glava IV: Dvustranno proste prsteny}}{18}
\tocsec{18}{Pomočna lema}{18}
\tocsec{19}{Prědstavjenja dvustranno prostyh hiperkompleksnyh sistemov v nekomutativnyh tělah, ktore razširjajut jih koeficientne polja}{19}
\tocsec{20}{Nekomutativne těla}{22}
\tocsec{21}{Galoisova teorija nekomutativnyh těl}{26}
\clearpage
\tocline{\textbf{Glava V: Faktorne sistemy}}{29}
\tocsec{22}{Grupa těl s danym centrom}{29}
\tocsec{23}{Faktorne sistemy}{30}
\tocsec{24}{Množenje faktornyh sistemov}{33}
\tocsec{25}{Normalno prědstavjenje $\mathfrak K_r$ s Galoisovymi maksimalnymi komutativnymi podpoljami}{39}
\tocsec{26}{Množenje skrěženyh prědstavjenij}{43}
\medskip
\tocline{\textbf{Glava VI: Teorija skrěženyh produktov}}{44}
\tocsec{27}{Prědstavjenja $\mathfrak R_r$ kako skrěžene produkty}{44}
\tocsec{28}{Produktny teorem za faktorne sistemy}{47}
\tocsec{29}{Teorem o glavnom rodu v minimalnom}{50}
\tocsec{30}{Specializacija na ciklične razpadne polja}{50}
\tocsec{31}{Priměnjenja cikličnogo specialnogo slučaja}{52}
\endgroup

\clearpage
\section*{Uvod}\label{einleitung}

Ta lekcija obrabja obću teoriju prědstavjenij prstenov, ktora vznikla s
jednoj strany iz teorije prědstavjenij konečnyh grup, a s drugoj strany
iz algebraičnoj i aritmetičnoj teorije hiperkompleksnyh čislovyh
sistemov. To obćejše razuměnje teorije prědstavjenij imaje slědne
prědnosti nad dosěgašnjim pristupom (\foreign{Frobenius},
\foreign{Schur}):

Zato že razgledamo ne samu grupu, ale grupovo prsten, a obćejše
hiperkompleksny sistem ili libovoljno prsten, možno priměniti
teoriju prstenov i idealov i tym osvoboditi teoriju od mnogyh nejasnyh
izčislenij. To priměnjenje obćej teorije prstenov i idealov ne samo
uproščaje teoriju, ale takože vodi ju dalje --- napr. v pytanju o
odnošenjah grupy s jej podgrupami --- i otvara jej nove oblasti
priměnjenja. Pomocju obćej teorije prědstavjenij možno dati novo, vrlo
elegantno osnovanje Galoisovoj teoriji i --- možda ješče važnějše ---
ustanoviti Galoisovu teoriju nekomutativnyh těl.

Osnovne pojmy teorije grup, prstenov i idealov i obće teoremy iz teorije
prědstavjenij možno najdti v dělu gospodičny \foreign{Noether}
\foreign{Hyperkomplexe Größen und Darstellungstheorie},
\foreign{Mathematische Zeitschrift} 30, str.~641. To dělo, citovano
kako \foreign{M.Z.}, trěba razgledati dopolnjenjem k tomu izloženju.
Isty material jest takože v zapisah lekcije gospodičny
\foreign{Noether} \foreign{Über Hyperkomplexe Größen und
Gruppencharaktere}, zimny semestr 1927/28.

Pojmy prědstavjenje i modul prědstavjenja sut tu kratko objasnjene,
da by vvesti recipročny modul prědstavjenja i recipročna
prědstavjenja, ktore sut nove samo formalno.

\section*{Glava I. Prědstavjenja}\label{kapitel-i.-darstellungen}

$\mathfrak{o}$ nehaj bude prsten --- prsten, ktore se prědstavja --- a
$\mathsf T$ drugo prsten --- to, v ktorom se prědstavja.

% BEGIN BOOK_S01
\subsection*{§ 1. Definicija prěmogo prědstavjenja}\label{definition-der-direkten-darstellung}

\emph{Prěmo prědstavjenje} \(n\)-togo stupenja od \(\mathfrak{o}\) v \(\mathsf T\) jest podprsten \(\mathfrak D\) prstena \(\mathsf T_n\) matric \(n\)-togo stupenja s elementami iz \(\mathsf T\), na ktoro jest \(\mathfrak{o}\) obrazovano prstenovym homomorfizmom. V znakah

\[
\mathfrak{o}\to\mathfrak D\subseteq\mathsf T_n.
\]

\subsection*{Definicija recipročnogo prědstavjenja}\label{definition-der-reziproken-darstellung}

\emph{Recipročno prědstavjenje} \(n\)-togo stupenja od \(\mathfrak{o}\) v \(\mathsf T^*\) jest podprsten \(\mathfrak D^*\) prstena \(\mathsf T_n^*\) matric \(n\)-togo stupenja s elementami iz \(\mathsf T^*\), na ktoro jest \(\mathfrak{o}\) obrazovano recipročnym prstenovym homomorfizmom. Recipročny prstenovy homomorfizm znači slědujuće: ako elementam iz \(\mathfrak{o}\), \(c,d\), odpovědajut elementi iz \(\mathfrak D^*\), \(c^*,d^*\), to elementu iz \(\mathfrak{o}\), \(cd\), trěba odpovědati element iz \(\mathfrak D^*\), \(d^*\cdot c^*\), a elementu iz \(\mathfrak{o}\), \(c+d\), element iz \(\mathfrak D^*\), \(c^*+d^*\). V znakah

\[
\mathfrak{o}\nrightarrow\mathfrak D^*\subseteq\mathsf T_n^*.
\]

Kako priměr pojma recipročnogo homomorfizma: ako matricam ręda \(n\) nad komutativnym poljem priręditi jih transponovane matrice, to priręđenje jest recipročny izomorfizm prstenov.

Obće možno k libovoljnomu prstenu \(\mathsf T\) konstruovati recipročno izomorfno prsten \(\mathsf T^*\) slědujućim načinom:

Vsakomu elementu iz \(\mathsf T\), \(c\), prirędimo simbol \(c^*\) i definirujemo složenje tyh simbolov formulom \(c^*+d^*=(c+d)^*\), a množenje formulom

\[
c^*\cdot d^*=(dc)^*.
\]

Množstvo simbolov \(c^*\) jest prsten, recipročno izomorfno s \(\mathsf T\), i označujemo je \(\mathsf T^*\).
% END BOOK_S01

% BEGIN BOOK_S02
\subsection*{§ 2. Klasy prědstavjenij}\label{darstellungsklassen}

Prěmo prědstavjenje \(\mathfrak D\) od \(\mathfrak{o}\) v \(\mathsf T\) možno prětvoriti v izomorfno prědstavjenje \(\mathfrak D'\) tako, že se ono po elementah transformuje regularnoj matricoju \(P\), t.~j. elementu iz \(\mathfrak{o}\), \(c\), vměsto \(C\subseteq\mathfrak D\) prirędi se \(P^{-1}CP\subseteq P^{-1}\mathfrak D P\).

Dva prědstavjenja, ktore možno tako transformovati jedno v drugo, nazyvajut se \emph{ekvivalentnymi}.

Vsa prědstavjenja, ekvivalentne jednomu prědstavjenju, tvoręt \emph{klasu prědstavjenij} (k. p.).

Tako samo definujeme ekvivalentnost recipročnyh prědstavjenij; dva recipročna prědstavjenja nazyvajut se ekvivalentnymi, ako jedno možno dostati iz drugogo transformacijoju regularnoj matricoju. Vsa prědstavjenja, ekvivalentne jednomu ustanovjenomu recipročnomu prědstavjenju, tvoręt \emph{klasu recipročnyh prědstavjenij}.
% END BOOK_S02

% BEGIN BOOK_S03
\subsection*{\texorpdfstring{\(\S\) 3. Moduly prědstavjenij}{\textbackslash S 3. Moduly prědstavjenij}}\label{s-3.-darstellungsmoduln}

Definicija \emph{prěmogo modula prědstavjenja} (p. m. p.).

Prěmy modul prědstavjenja od \(\mathfrak{o}\) po odnošenju k \(\mathsf T\) jest vsaky dvojny \(\mathfrak{o}\)-\(\mathsf T\)-modul \(\mathfrak M\) (t.~j. aditivno zapisana Abelova grupa, ktora jest lěvym modulom po odnošenju k \(\mathfrak{o}\), pravym modulom po odnošenju k \(\mathsf T\) i izpolnja smješany zakon asociativnosti: \(c\in\mathfrak{o}\), \(m\in\mathfrak M\), \(\tau\in\mathsf T\); \((cm)\tau=c(m\tau)\)), ktory možno zapisati kako prěmu sumu konečnogo čisla jednočlennyh \(\mathsf T\)-modulov --- pričem \(\tau=0\) slěduje iz \(x_i\tau=0\) ---

\[
\mathfrak M=x_1\cdot\mathsf T+x_2\cdot\mathsf T+\cdots+x_n\cdot\mathsf T.
\]

Definicija \emph{recipročnogo modula prědstavjenja} (r. m. p.).

Recipročny modul prědstavjenja od \(\mathfrak{o}\) po odnošenju k \(\mathsf T^*\) jest vsaky lěvy \(\mathfrak{o}\)-modul i lěvy \(\mathsf T^*\)-modul \(\mathfrak M^*\), ktory izpolnja zakon komutativnosti (ako \(c\in\mathfrak{o}\), \(m\in\mathfrak M^*\), \(\tau\in\mathsf T^*\), to \(c(\tau^*m^*)=\tau^*(cm)\)) i ktory možno zapisati kako prěmu sumu konečnogo čisla prostyh \(\mathsf T^*\)-modulov --- pričem \(\tau^*=0\) slěduje iz \(\tau^*\cdot x_i^*=0\) ---

\[
\mathfrak M^*=\mathsf T^*\cdot x_1+\cdots+\mathsf T^*\cdot x_n.
\]
% END BOOK_S03

% BEGIN BOOK_S04
\subsection*{§ 4. Svęz medžu modulami prědstavjenij i prědstavjenjami}\label{der-zusammenhang-zwischen-darstellungsmoduln-und-darstellungen}

\textbf{1.} \emph{Vsaky modul prědstavjenja jednoznačno vodi k klase prědstavjenij (p. m. p. k prěmoj klase prědstavjenij, r. m. p. k recipročnoj klase prědstavjenij).}

I to slědujućim načinom:

Množenje prěmogo modula prědstavjenja \(\mathfrak M\) elementom iz \(\mathfrak{o}\), \(c\), znači homomorfizm od \(\mathfrak M\) na sebe samogo, ktory zaradi zakona asociativnosti \((cm)\tau=c(m\tau)\) dopušča elementy iz \(\mathsf T\) kako operatory. Ono jest zato \emph{linearna transformacija} od \(\mathfrak M\) v sebe, t.~j. priręđenje \(m\mapsto cm=m'\) imaje svojstvo:

\[
\text{Ako } y_i\mapsto y_i'(=cy_i) \text{ jest, to takože }
\sum y_i\tau_i\mapsto \sum y_i'\tau_i\left(=c\sum y_i\tau_i\right).
\]

Definujmo sumu \(\mathfrak T_1+\mathfrak T_2\) dvoh linearnyh transformacij

\[
\mathfrak T_1:\ m\mapsto m'(=c_1m)
\quad\text{i}\quad
\mathfrak T_2:\ m\mapsto m''(=c_2m)
\]

kako transformaciju \(m\mapsto m'+m''(=(c_1+c_2)m)\), a jih produkt \(\mathfrak T_1\cdot\mathfrak T_2\) kako složenu transformaciju \(m\mapsto (m'')'=c_1c_2m\). Togda množstvo različnyh linearnyh transformacij od \(\mathfrak M\), ktore stvarjajut elementi iz \(\mathfrak{o}\), jest prstenovo homomorfny obraz \(\overline{\mathfrak D}\) od \(\mathfrak{o}\).

Nyně izberimo libovoljnu bazu \(x_1,\ldots,x_n\) modula \(\mathfrak M\). Linearnu transformaciju, stvorjenu elementom \(c\), od \(\mathfrak M\) očevidno jest dostatočno zadati tym, kako se transformujut bazove elementi \(x_i\), to jest poznanjem matrice \((\gamma_{ij})=C\) sistema jednačenij

\[
x'_j=cx_j=\sum x_i\gamma_{ij}
\]

zato že za druge elementy \(x=\sum x_j\tau_j\) prosto jest \(x'=\sum x'_j\tau_j\). Te matrice \(C\) sut jednoznačno priręđene različnym linearnym transformacijam (zato že jedne oprěděljajut druge) i, kako se vidi neposrědnje, tvoręt prsten izomorfno s \(\overline{\mathfrak D}\), to jest prstenovo homomorfny obraz od \(\mathfrak{o}\), prědstavjenje \(n\)-togo stupnja od \(\mathfrak{o}\).

Izomorfizm medžu \(\overline{\mathfrak D}\) i prědstavjenjem sodrživan jest v slědujućih formulah: ako linearnoj transformaciji, stvorjenoj elementom \(c\in\mathfrak{o}\), odpovědaje matrica \(C\), to

\[
c(x_1,\ldots,x_n)=(x_1,\ldots,x_n)C
\]

i ako na jednaky način \(d\in\mathfrak{o}\) i \(D\in\mathsf T_n\) odpovědajut jedno drugomu

\[
d(x_1,\ldots,x_n)=(x_1,\ldots,x_n)D
\]

to očevidno jest

\[
\begin{aligned}
(c+d)(x_1,\ldots,x_n)
&=(cx_1+dx_1,\ldots,cx_n+dx_n)\\
&=(x_1,\ldots,x_n)(C+D)
\end{aligned}
\]

i

\[
\begin{aligned}
cd(x_1,\ldots,x_n)
&=c(dx_1,\ldots,dx_n)=c((x_1,\ldots,x_n)D)\\
&=(cx_1,\ldots,cx_n)D=(x_1,\ldots,x_n)CD.
\end{aligned}
\]

Tako linearny transformacije, stvorjene elementami iz \(\mathfrak{o}\), od \(\mathfrak M\) v sebe pri jih realizaciji posrědstvom ustanovjene bazy \(x_\nu\) modula \(\mathfrak M\) davajut prědstavjenje \(n\)-togo stupenja od \(\mathfrak{o}\) v \(\mathsf T\).

Ako \((y_1,\ldots,y_n)\) jest druga baza modula \(\mathfrak M\), ona nastava iz \(x_1,\ldots,x_n\) množenjem regularnoj matricoju

\[(y_1,\ldots,y_n)=(x_1,\ldots,x_n)P\]

Ako linearnaja transformacija od \(\mathfrak M\) po odnošenju k baze \(x_\nu\) jest realizovana matricoju \(C\), togda po odnošenju k baze \(y_\nu\) ona jest realizovana matricoju \(P^{-1}CP\), zato že ako

\[c(x_1,\ldots,x_n)=(x_1,\ldots,x_n)C\]
to imaje město

\[
\begin{aligned}
c(y_1,\ldots,y_n)
&=c(x_1,\ldots,x_n)P=((x_1,\ldots,x_n)C)P\\
&=(x_1,\ldots,x_n)CP=((y_1,\ldots,y_n)P^{-1})CP\\
&=(y_1,\ldots,y_n)P^{-1}CP.
\end{aligned}
\]

Iz toga neposrědnje slěduje:

\emph{\(\mathfrak M\) davaje vsa prědstavjenja jednoj klasy prědstavjenij, ako za realizaciju linearnyh transformacij od \(\mathfrak M\), stvorjenyh elementami iz \(\mathfrak{o}\), upotrěbimo vse bazy modula \(\mathfrak M\).}

\begin{enumerate}
\def\labelenumi{\arabic{enumi}.}
\setcounter{enumi}{1}
\tightlist
\item
  \emph{Vsaka klasa prědstavjenij jest stvorjena modulom prědstavjenja na način ukazany pod 1.} (Ponovno se ograničujemo na prěme prědstavjenja.)
\end{enumerate}

\emph{Dokaz.} Izberemo ustanovjeno prědstavjenje iz klasy; nehaj ono bude zadano homomorfnym priręđenjem \(c\mapsto C\). Ako \(n\) jest stupenj prědstavjenja, to s \(n\) neizvěstnymi \(x_1,\ldots,x_n\) obrazujemo pravy \(\mathsf T\)-modul

\[\mathfrak{M} = x_1 \mathsf{T} + \dots + x_n \mathsf{T}\]

(s obyčnymi pravilami računanja \(\sum x_i\tau_i+\sum x_i\overline{\tau}_i=\sum x_i(\tau_i+\overline{\tau}_i)\), \(\sum x_i\tau_i=\sum x_i\overline{\tau}_i\) togda i samo togda, ako \(\tau_i=\overline{\tau}_i\) za vse \(i\), \((\sum x_i\tau_i)\varrho=\sum x_i\tau_i\varrho\)), i činimo ga lěvym \(\mathfrak{o}\)-modulom ustanovjenjem \(c\sum x_i\tau_i=\sum x_i'\tau_i\), ako \((x_1',\ldots,x_n')=(x_1,\ldots,x_n)C\), pričem pod \(C\) razumějemo matricu prědstavjenja, priręđenu elementu iz \(\mathfrak{o}\), \(c\).

Tym ustanovjenjem \(\mathfrak M\) jest učinjeno dvojnym \(\mathfrak{o}\)-\(\mathsf T\)-modulom, zato že prvo iz relacij homomorfizma slěduje

\[
c+d\mapsto C+D \quad \text{i} \quad cd\mapsto CD
\]

\[
(c+d)x_i=cx_i+dx_i
\]

\[
cd\cdot x_i=c\cdot dx_i,\quad\text{nadalje jest}\quad
c\left(\sum x_i\tau_i+\sum x_i\overline{\tau}_i\right)
=c\sum x_i\tau_i+c\sum x_i\overline{\tau}_i,
\]

to sut svojstva lěvogo modula; drugo, očevidno jest

\[c(\sum x_i \, \tau_i \, \varrho) = (c \, \sum x_i \, \tau_i) \cdot \varrho \,,\]

to jest smješany zakon asociativnosti.

\(\mathfrak{M}\) jest zato modul prědstavjenja, a upotrěbjeno prědstavjenje jest nim stvorjeno na način ukazany pod 1., ako za realizaciju njegovyh linearnyh transformacij upotrěbimo bazu \(x_1, \ldots, x_n\). Druge bazy davajut ostale prědstavjenja te klasy.

Za recipročna prědstavjenja i module prědstavjenij vse ide tak samo; dobyva se, že modul

\[\mathfrak{M}^* = \mathsf{T}^* x_1^* + \dots + \mathsf{T}^* x_n^*\]

stvarja recipročno prědstavjenje od \(\mathfrak{o}\) v \(\mathsf T^*\) tako, že elementu iz \(\mathfrak{o}\), \(c\), priręđuje matricu \(C\) v jednačenju

\[
c\begin{pmatrix}x_1\\ \vdots\\ x_n\end{pmatrix}
=\begin{pmatrix}cx_1\\ \vdots\\ cx_n\end{pmatrix}
=C\begin{pmatrix}x_1\\ \vdots\\ x_n\end{pmatrix}
\]

Različne bazy ponovno davajut različna prědstavjenja jednoj klasy.

\section*{Glava II. Galoisova teorija komutativnyh polj}\label{kapitel-ii.-galoissche-theorie-kommutativer-koerper}

Dosěga razvijena teorija prědstavjenij jest dostatočna za osnovanje Galoisovoj teoriji, ktora imaje prědnost nad obyčnoj teorijeju, že ne potrěbuje specialne sistemy tvorečih elementov ni teoremy o polinomah neizvěstnyh.\srcfn{1}{Za slědujuće sravni M. Z. §§ 15 --- 20. M. Z. § 21 sodrživaje prve teoremy slědujućego odděla o Galoisovoj teoriji. Za definiciju hiperkompleksnyh sistemov vidi M. Z. § 6.}
% END BOOK_S04

% BEGIN BOOK_S05
\subsection*{§ 5. Razširjenje koeficientov hiperkompleksnyh sistemov}\label{koeffizientenerweiterung-hyperkomplexer-systeme}

Formulujemo slědujući obći teorem, čij dokaz ostavjamo neizvedenym, zato že on jest skoro samorazumny. Nehaj \(\mathfrak{o}=c_1\mathsf P+\cdots+c_n\mathsf P\) bude hiperkompleksny sistem nad komutativnym poljem \(\mathsf P\). Nehaj \(\Omega\) bude razširjujuće prsten od \(\mathsf P\), čiji centar sodrživaje \(\mathsf P\). Definujemo \emph{razširjeny sistem} \(\mathfrak{o}_{\Omega}\) kako množstvo vsih formalnyh sum \(\sum c_i\omega_i\), \(\omega_i\in\Omega\), ktoro najprvo ustanovjenjami

\[
\begin{array}{l}
\sum c_i\omega_i=\sum c_i\omega_i' \quad\text{togda i samo togda, ako } \omega_i=\omega_i' \text{ za vse } i;\\
\sum c_i\omega_i+\sum c_i\omega_i'=\sum c_i(\omega_i+\omega_i');\\
\omega\sum c_i\omega_i=\sum c_i\omega\omega_i;\quad
\sum c_i\omega_i\cdot\omega=\sum c_i\omega_i\omega;\\
\sum c_i\omega_i=\sum \omega_i c_i
\end{array}
\]

postavaje \(\Omega\)-modulom, komutativno svęzanym s \(\Omega\), a daljšim ustanovjenjem

\[
\sum c_i\omega_i\cdot \sum c_i\omega_i
=\sum_i\sum_j c_ic_j\omega_i\omega_j
\]

--- pričem \(c_ic_j\) jest už definovany kako element iz \(\mathfrak{o}\) --- postavaje razširjujućim prstenom od \(\mathfrak{o}\). \(\mathfrak{o}_{\Omega}\) imaje rang \(n\) nad \(\Omega\) i možno ga v formě

\[
\mathfrak{o}_{\Omega}=c_1\Omega+\cdots+c_n\Omega
\]

zapisati.

\emph{\(\mathfrak{o}_{\Omega}\) jest definovany nezavisno od bazy \(c_1,\ldots,c_n\). Ako \(\Omega\) jest komutativno polje, to \(\mathfrak{o}_{\Omega}\) jest hiperkompleksny sistem nad \(\Omega\).}
% END BOOK_S05

% BEGIN BOOK_S06
\subsection*{§ 6. Ireducibilne prědstavjenja komutativnyh sistemov}\label{die-irreduziblen-darstellungen-kommutativer-systeme}

Nehaj \(\mathfrak Z=z_1\mathsf P+\cdots+z_n\mathsf P\) bude komutativny sistem nad \(\mathsf P\). Hočemo ustanoviti vsa ireducibilna prědstavjenja od \(\mathfrak Z\) v algebraično zatvorjenom razširjenom polju \(\Omega\) od \(\mathsf P\). Zato že vsako prědstavjenje od \(\mathfrak Z\) v \(\Omega\) vodi k prědstavjenju od \(\mathfrak Z_\Omega=z_1\Omega+\cdots+z_n\Omega\) --- zato že prědstavjenja bazovyh elementov, už sodrživanyh v \(\mathsf P\), sut dostatočne za ustanovjenje cělokupnogo prědstavjenja --- i obratno vsako prědstavjenje od \(\mathfrak Z_\Omega\) sodrživaje prědstavjenje od \(\mathfrak Z\), možno pytati takože za prědstavjenja od \(\mathfrak Z_\Omega\). Ona sut po M. Z. § 20 sodrživana v regularnom prědstavjenju od \(\mathfrak Z_\Omega\), t.~j. v prědstavjenju od \(\mathfrak Z_\Omega\), ktoro nastava, ako sam \(\mathfrak Z_\Omega\) upotrěbimo kako modul prědstavjenja. Dostatočno jest čak ograničiti se na \(\mathfrak Z_\Omega/\mathfrak C\) kako modul prědstavjenja, pričem pod \(\mathfrak C\) razumějemo radikal od \(\mathfrak Z_\Omega\). (vidi M. Z. § 19).

Ale \(\mathfrak Z_\Omega/\mathfrak C\) jest polno reducibilny sistem, to jest prěma suma polj konečnogo stupenja nad \(\Omega\). Zato že \(\Omega\) jest algebraično zatvorjeno, te sumandy sut izomorfne s \(\Omega\).

\[
\mathfrak Z_\Omega/\mathfrak C=e_1\Omega+\cdots+e_t\Omega.
\]

Elementi \(e_\nu\) sut komponenty jedinice. zato imaje město:

\emph{\(\mathfrak Z_\Omega\) imaje točno \(t\) različnyh klas ireducibilnyh prědstavjenij v \(\Omega\), vse prvogo stupenja. Pri tom \(t\) jest rang od \(\mathfrak Z_\Omega/\mathfrak C\), zato najvyše ravny \(n\).}

Vsaka iz tyh klas prědstavjenij sodrživaje samo jedno prědstavjenje, zato že transformacija prědstavjenja prvogo stupenja v komutativnom polju ne izměnjaje prědstavjenje: \(\pi^{-1}\cdot\alpha\cdot\pi=\alpha\).

Zato obstajajut točno \(t\) različne homomorfizmy \(\Theta_1,\ldots,\Theta_t\) od \(\mathfrak Z_\Omega\) na podprsteny od \(\Omega\).
% END BOOK_S06

% BEGIN BOOK_S07
\subsection*{§ 7. Izomorfizmy polja}\label{die-isomorphismen-eines-koerpers}

Nyně nehaj \(\mathfrak Z\) bude polje. Vsaky homomorfizm \(\Theta_i\) od \(\mathfrak Z_\Omega\) na podprsten od \(\Omega\) mora izomorfno obrazovati \(\mathfrak Z\) na podpolje \(Z_i\) od \(\Omega\). Zato že ili jest \(Z_i=0\) --- a to tu ne jest slučaj, zato že jedinice od \(\mathfrak Z\) odpovědaje jedinica od \(\Omega\) --- ili \(Z_i\) jest izomorfny obraz od \(\mathfrak Z\). Zato:

\emph{Ako \(\mathfrak Z\) jest konečno razširjenje \(n\)-togo stupenja od \(\mathsf P\), to v \(\Omega\) obstajajut toliko različne polja, izomorfne s \(\mathfrak Z\), \(Z_i\), koliky jest rang \(t\) od \(\mathfrak Z_\Omega/\mathfrak C\). Čislo polj \(Z_i\) jest zato najvyše ravno \(n\).}

Nyně nazyvajemo \(\mathfrak Z\) poljem \emph{prvogo roda} nad \(\mathsf P\), ako čislo polj \(Z_i\) jest ravno \(n\), a poljem \emph{drugogo roda}, ako to čislo jest menje od \(n\).

\emph{\(\mathfrak Z\) jest togda i samo togda prvogo roda nad \(\mathsf P\), ako \(\mathfrak Z_\Omega\) jest polno reducibilno.}

Tu vidimo prědnost toga osnovanja Galoisovoj teoriji nad obyčnym. V obyčnoj Galoisovoj teoriji razgledajut se izomorfizmy jednogo iz polj \(Z_i\) na druge. Tu se izběgaje ta nesimetrija. Za dokaz, že polje ne imaje vyše konjugovanyh polj, než koliko ukazuje jego stupenj, obyčno se upotrěbje primitivny element. Tu toj fakt prosto slěduje iz toga, že vsa ireducibilna prědstavjenja od \(\mathfrak Z_\Omega\) sut sodrživana v regularnom prědstavjenju.
% END BOOK_S07

% BEGIN BOOK_S08
\subsection*{§ 8. Razpadno polje}\label{der-zerfaellungskoerper}

Polja \(Z_i\) sut ireducibilna prědstavjenja od \(\mathfrak Z\) v \(\Omega\). Ako zato polje \(\Gamma\) medžu \(\mathsf P\) i \(\Omega\), \(\mathsf P\subseteq\Gamma\subseteq\Omega\), sodrživaje vsa \(Z_i\), to už \(\mathfrak Z_\Gamma/\mathfrak C\) mora se razpadati na \(t\) polj prvogo stupenja \(\Gamma\). Obratno, da by \(\mathfrak Z_\Gamma/\mathfrak C\) razpadalo se na \(t\) polj, ktore sut togda ireducibilnymi prědstavjenjami prvogo stupenja od \(\mathfrak Z_\Gamma\), polja \(Z_i\) morajut byti sodrživane v \(\Gamma\). Ako razpadnym poljem od \(\mathfrak Z\) nazyvamo polje \(\Gamma\) s svojstvom, že \(\mathfrak Z_\Gamma/\mathfrak C\) postavaje prěmoju sumoju \(t\) polj, to imajemo:

\emph{Galoisovo polje polj, izomorfnyh s \(\mathfrak Z\), \(Z_i\), t.~j. jih kompozitum, jest minimalno razpadno polje od \(\mathfrak Z\).}
% END BOOK_S08

% BEGIN BOOK_S09
\subsection*{§ 9. Izomorfizmy od \(Z_1\) na \(Z_i\)}\label{die-isomorphismen-von-z1-auf-die-zi}

\(\Theta_i\), ako razgledamo samo \(\mathfrak Z\), jest izomorfizm od \(\mathfrak Z\) na \(Z_i\). zato \(S_i=\Theta_i\Theta_1^{-1}\) jest izomorfizm od \(Z_1\) na \(Z_i\). Zato že \(\Theta_i\) sut različne, takože \(S_i\) sut različne. \(S_i\) sut izomorfizmy jednogo polja \(Z_1\) na jego konjugovana polja, ktore se obyčno razgledajut v Galoisovoj teoriji. Drugyh ne imaje, zato že ne imaje drugyh \(\Theta_i\).
% END BOOK_S09

% BEGIN BOOK_S10
\subsection*{§ 10. Galoisova grupa}\label{die-galoissche-gruppe}

Ako \(\mathfrak Z\) jest Galoisovo polje, to jest ako polja \(Z_i\) sběgajut sę kako množstva elementov, to \(S_i\) tvoręt grupu vsih izomorfizmov, ktore ostavjajut \(\mathsf P\) nepodvižnym, od \(Z\) na sebe. Zato že vsaky produkt \(S_iS_j\) jest avtomorfizm od \(Z\), a zato že ne imaje drugyh kromě \(S_i\), \(S_iS_j\) jest ravny nekomu \(S_k\). zato \(S_i\) tvoręt grupu vsih avtomorfizmov, ktore ostavjajut vsaky element iz \(\mathsf P\) nepodvižnym, od \(Z\); to jest \emph{Galoisova grupa od \(Z\) po odnošenju k \(\mathsf P\)}.
% END BOOK_S10

% BEGIN BOOK_S11
\subsection*{§ 11. Glavny teorem Galoisovoj teorije}\label{der-hauptsatz-der-galoisschen-theorie}

Nyně nehaj \(\mathfrak Z\) bude Galoisovo razširjenje prvogo roda od \(\mathsf P\). Dokazujemo glavny teorem Galoisovoj teorije:

\emph{Polja \(\Gamma\) medžu \(\mathsf P\) i \(\mathfrak Z\) možno jednoznačno vzajemno priręditi podgrupam \(\mathfrak H\) Galoisovoj grupy \(\mathfrak G\) od \(Z\) po odnošenju k \(\mathsf P\) tako, že polju \(\Gamma\) priręđena grupa \(\mathfrak H\) sostoji se iz vsih avtomorfizmov od \(Z\), ktore ostavjajut vsaky element iz \(\Gamma\) nepodvižnym, a \(\Gamma\) jest množstvo vsih elementov iz \(Z\), ktore ostavajut nepodvižne pri vsih avtomorfizmah, sodrživanyh v \(\mathfrak H\).}

Kratko to formulujemo tako:

\begin{enumerate}
\item \emph{Ako \(\mathfrak H\) jest invariantna grupa od \(\Gamma\), to \(\Gamma\) jest invariantno polje od \(\mathfrak H\).}
\item \emph{Ako \(\Gamma\) jest invariantno polje od \(\mathfrak H\), to \(\mathfrak H\) jest invariantna grupa od \(\Gamma\).}
\end{enumerate}

\paragraph{1. možno svesti na pomoćnu lemu:}\label{laesst-sich-zurueckfuehren-auf-den-hilfssatz}

\emph{Ako \(\mathsf P \subseteq \Sigma \subseteq \mathsf T \subseteq \mathsf Z\) i \(\mathsf T\) imaje stupenj \(t\) nad \(\Sigma\), to vsaky izomorfizm od \(\Sigma\) na neko konjugovano polje (potrěbno sodrživano v \(\mathsf Z\)) možno prodolžiti do \(t\) različnyh izomorfizmov od \(\mathsf T\) na konjugovana polja.}

Ako prědpoložimo, že \(\mathfrak H\) jest invariantna grupa od \(\Sigma\), a \(\mathsf T\) jest invariantno polje od \(\mathfrak H\), možno zaključiti slědujuće: Elementi iz \(\mathfrak H\) sut prodolženja na cělo \(\mathsf Z\) identičnogo izomorfizma od \(\Sigma\). Ako te prodolženja priměnimo samo na \(\mathsf T\), iz pomoćnoj lemy slěduje, že ona razpadajut se na \(t\) klas tako, že izomorfizmy iste klasy sut identične na \(\mathsf T\), a izomorfizmy različnyh klas sut različne na \(\mathsf T\). Ale zato že \(\mathsf T\) jest invariantno polje od \(\mathfrak H\), mora byti \(t=1\), \(\mathsf T=\Sigma\), tym jest dokazano.

Za dokaz pomoćnoj lemy vrnemo se k poljam \(\mathfrak S\), \(\mathfrak T\) medžu \(\mathsf P\) i \(\mathsf Z\), ktorym sut \(\Sigma\) i \(\mathsf T\) priręđene posrědstvom \(\Theta_1\). Togda trěba dokazati:

\emph{Vsaky izomorfizm \(H_i\) od \(\mathfrak S\) na neko \(\Sigma_i\) možno na točno \(t\) različnyh načinov prodolžiti do izomorfizmov od \(\mathfrak T\) na podpolja \(\mathsf T_i\) od \(\mathsf Z\).}

Zaradi prědpoloženja, že \(\mathfrak Z\) jest prvogo roda nad \(\mathsf P\), \(\mathfrak Z_\Omega/\mathfrak C\) jest identično s \(\mathfrak Z_\Omega\), zato \(\mathfrak S_\Omega\) razpada se na toliko polj, koliko ukazuje njegov stupenj \(s\)
\[
\mathfrak S_\Omega = E_1\Omega+\cdots+E_s\Omega .
\]
Moduly prědstavjenij \(E_\nu\Omega\) stvarjajut \(s\) izomorfizmov od \(\mathfrak S\) na \(s\) podpolj \(\Sigma_1,\ldots,\Sigma_s\). Za najdanje izomorfizmov od \(\mathfrak T\) potrěbujemo \(\mathfrak T_\Omega\). Zato že \(E_\nu\Omega\) sut idealy od \(\mathfrak S_\Omega\), imaje město \(E_\nu\Omega=\mathfrak S_\Omega E_\nu\). Nadalje očevidno jest
\[
\mathfrak T_\Omega=\mathfrak T\cdot\Omega=\mathfrak T\cdot\mathfrak S_\Omega .
\]
Zato jest
\[
\mathfrak T_\Omega
=E_1\mathfrak S_\Omega\cdot\mathfrak T+\cdots+E_s\mathfrak S_\Omega\cdot\mathfrak T
=E_1\mathfrak T_\Omega+\cdots+E_s\mathfrak T_\Omega .
\]
To jest razklad od \(\mathfrak T_\Omega\) na idealy. \(E_\nu\cdot\mathfrak T_\Omega\) razpada se na \(t\) prostyh idealov. Zato že sut pravdive relacije stupnjev
\[
[\mathfrak T_\Omega\cdot E_\nu:\mathfrak S_\Omega E_\nu]
\leqq [\mathfrak T_\Omega:\mathfrak S_\Omega]=[\mathfrak T:\mathfrak S]=t,
\]
zato zaradi \(\mathfrak S_\Omega E_\nu=\Omega E_\nu\simeq\Omega\):
\[
[\mathfrak T_\Omega\cdot E_\nu:\Omega]
=[\mathfrak T_\Omega E_\nu:\mathfrak S_\Omega E_\nu]\leqq t,
\]
i zato že suma vsih stupnjev \([\mathfrak T_\Omega E_\nu:\Omega]\) mora byti ravna \(s\cdot t\):
\[
[\mathfrak T_\Omega\cdot E_\nu:\Omega]=t,\qquad
E_\nu\mathfrak T_\Omega=e_\nu^{(1)}\cdot\Omega+\cdots+e_\nu^{(t)}\Omega .
\]
Vsaky \(e_\nu^{(i)}\Omega\) davaje ireducibilno prědstavjenje od \(\mathfrak T_\Omega\). \emph{Pokažimo ješče}, že vsa ta prědstavjenja, ako jih priměnimo samo na \(\mathfrak S_\Omega\), sběgajut sę s prědstavjenjem, stvorjenym posrědstvom \(E_\nu\cdot\Omega\); togda jesmo gotovi. To ale jest jasno, zato že iz
\[
s\cdot E_\nu=E_\nu\sigma_\nu\quad
(s\in\mathfrak S_\Omega\text{ imaje }\sigma_\nu\text{ kako prědstavjenje})
\]
slěduje
\[
\sum s\cdot e_\nu^{(i)}=\sum e_\nu^{(i)}\sigma_\nu,
\]
zato zaradi jednoznačnosti zapisa sumy
\[
s\cdot e_\nu^{(i)}=e_\nu^{(i)}\sigma_\nu,
\]
nove prědstavjenja takože prirędžajut elementu \(s\) element \(\sigma_\nu\).

\paragraph{2. Nehaj \(\mathsf T\) bude invariantno polje od \(\mathfrak H\).}
Hočemo poznati \(\mathfrak H\) kako invariantnu grupu od \(\mathsf T\); za to \emph{jest dostatočno} pokazati element iz \(\mathsf T\), ktory ne ostavaje nepodvižnym pri nekom avtomorfizmu, ne naležečem k \(\mathfrak H\).

Elementi \(S_i\) Galoisovoj grupy \(\mathfrak G\) od \(\mathsf Z\) sut avtomorfizmy, ktore ostavjajut \(\mathsf P\) nepodvižnym, od \(\mathsf Z\). Činimo jih avtomorfizmami od \(\mathfrak Z_{\mathsf Z}\) ustanovjenjem
\[
S_i z=z,\quad \text{ako } z \text{ jest v } \mathfrak Z.
\]
Priměnjenje od \(S_i\) na
\[
\mathfrak Z_{\mathsf Z}=e_1\mathsf Z+\cdots+e_n\mathsf Z
\]
mora prěměnjivati prěme sumandy \(e_\nu\mathsf Z\), zato že one sut jednoznačno oprěděljene; v osoblivosti mora byti
\[
S_i e_\nu=e_{\nu'},\qquad\text{i}\qquad S_i e_\nu\ne S_k e_\nu
\quad\text{za } i\ne k
\]
Ako \(S_1,\ldots,S_h\) sut elementi iz \(\mathfrak H\), izberemo neko \(e_\nu\), recimo \(e'\), i razgledamo vse \(e_\nu\), ktore iz njega stvarjajut \(S_i\), \(i\leqq h\):
\[
e_1=S_1e',\ldots,e_h=S_he'.
\]
Zato že, ako \(i\leqq h, j\leqq h\),
\[
S_i e_j=S_iS_j e'=S_ke'=e_k \quad \text{pri}\quad k\leqq h
\]
vsaky \(S_i\) iz \(\mathfrak H\) prěměnjuje \(e_1,\ldots,e_h\) medžu soboju,
\[
E_1=e_1+\cdots+e_h
\]
jest zato pri vsih \(S_i\in\mathfrak H\) invariantny element iz \(\mathfrak Z\). Ako prědstavimo \(E_1\) posrědstvom bazy od \(\mathfrak Z\), jego koeficienty morajut byti pri \(\mathfrak H\) invariantne elementi iz \(\mathsf Z\); zato naležet k \(\mathsf T\).

\(E_1\) ne jest invariantny pri nikakom elementu izvan \(\mathfrak H\), \(S_i\). Zato že za tako \(S_i\)
\[
S_i E_1=e_i+\cdots
\]
sodrživaje medžu komponentami, ktore ne nastupajut medžu \(e_1,\ldots,e_h\), komponentu \(e_i\), ale zapis sumy od \(E_1\) jest jednoznačny.
% END BOOK_S11

% BEGIN BOOK_S12
\subsection*{§ 12. Formalno značenje komponentov \(e_i\)}\label{die-formale-bedeutung-der-komponenten-ei}

Za osnovu izberemo ustanovjenu bazu \(z_1,\ldots,z_n\) od \(\mathfrak Z\).
\[
\mathfrak Z=z_1\mathsf P+\cdots+z_n\mathsf P
\]
(Nyně o \(\mathfrak Z\) prědpokladamo samo, že ono jest prvogo roda). Izomorfizm \(\Theta_i\) obrazuje sistem \(z_1,\ldots,z_n\) na bazu \(\zeta_1^{(i)},\ldots,\zeta_n^{(i)}\) od \(\mathsf Z_i\), i to relacijami
\[
z_\nu e_i=e_i\zeta_\nu^{(i)} .
\]

Obratno, komponenta jedinice \(e_j\) mora byti izraziva posrědstvom \(z_\nu\) s koeficientami iz \(\mathsf Z\)
\[
e_j=\sum_\nu \varrho_\nu^{(j)}z_\nu .
\]
zato imajemo
\[
e_j=e_j^2=\sum \varrho_\nu^{(j)}z_\nu e_j
=\sum \varrho_\nu^{(j)}\cdot\zeta_\nu^{(j)}\cdot e_j
\]
i
\[
0=e_je_k=\sum \varrho_\nu^{(j)}\cdot z_\nu\cdot e_k
=\sum \varrho_\nu^{(j)}\cdot\zeta_\nu^{(k)}\cdot e_k,
\]
zato:
\[
\sum_\nu \varrho_\nu^{(j)}\cdot\zeta_\nu^{(j)}=1;
\qquad
\sum_\nu \varrho_\nu^{(j)}\cdot\zeta_\nu^{(k)}=0,\quad
\text{ako } j\ne k.
\]
To znači: Matrice
\[
\begin{pmatrix}
\zeta_1^{(1)} & \cdots & \zeta_n^{(1)}\\
\cdots & \cdots & \cdots\\
\cdots & \cdots & \cdots\\
\cdots & \cdots & \cdots\\
\zeta_1^{(n)} & \cdots & \zeta_n^{(n)}
\end{pmatrix}
\quad\text{i}\quad
\begin{pmatrix}
\varrho_1^{(1)} & \cdots & \varrho_1^{(n)}\\
\cdots & \cdots & \cdots\\
\cdots & \cdots & \cdots\\
\cdots & \cdots & \cdots\\
\varrho_n^{(1)} & \cdots & \varrho_n^{(n)}
\end{pmatrix}
\]
sut vzajemno obratne. To možno takože izraziti slědujuće: \(\varrho_1^{(i)},\ldots,\varrho_n^{(i)}\) jest baza komplementarna k \(\zeta_1^{(i)},\ldots,\zeta_n^{(i)}\). Ona imaje rolu v teoriji diferenty čislovogo polja.
% END BOOK_S12

% BEGIN BOOK_S13
\subsection*{§ 13. Obći teorem o prodolženju}\label{der-allgemeine-fortsetzungssatz}

Teorem o prodolženju iz \(H_i\), dokazan na str.~10, možno obćiti.

Nehaj \(\mathfrak Z=z_1\mathsf P+\cdots+z_n\mathsf P\) bude komutativny, hiperkompleksny, polno reducibilny sistem. Nehaj \(\mathfrak S\) i \(\mathfrak T\) budut -- v toj situaciji potrěbno takože polno reducibilne -- podsistemy od \(\mathfrak Z\):
\[
\mathsf P\subseteq\mathfrak S\subseteq\mathfrak T\subseteq\mathfrak Z .
\]
Po M. Z. § 13 vsaky iz troh sistemov imaje jedinicu. Ale te tri jedinice nikako ne morajut byti jednake. (Sigurno ne sut jednake, ako \(\mathfrak S\) i \(\mathfrak T\) sut idealy od \(\mathfrak Z\), različne od \(\mathfrak Z\).) Nehaj \(E\) bude jedinica od \(\mathfrak S\). Sistem \(\mathfrak T\cdot E\) jest modul konečnogo ranga nad \(\mathfrak S\):
\[
E\mathfrak T=a_1\mathfrak S+\cdots+a_t\mathfrak S
\]
(\(\mathfrak T\) obće ne jest!).

\emph{Dokaz.} \(\mathfrak S\) jest prěma suma polj: \(\mathfrak S=\mathfrak K_1+\cdots+\mathfrak K_r\). Tomu odpovědaje razklad od \(\mathfrak T E\) na prěmu sumu. Vsaky sumand od \(\mathfrak T E\) mora sodrživati jedno iz polj \(\mathfrak K_i\), inače by \(E\) anulovalo go, ale \(E\) jest jedinica od \(\mathfrak T E\). Zato \(\mathfrak T E\) jest prěma suma razširjujućih prstenov od \(\mathfrak K_i\), t.~j. prěma suma formy
\[
\mathfrak T E=a_{11}\mathfrak K_1+\cdots+a_{1p_1}\mathfrak K_1+\cdots+a_{rp_r}\mathfrak K_r,
\]
ili \(E\mathfrak T=a_{11}\mathfrak S+\cdots+a_{rp_r}\mathfrak S\), tym jest dokazano.

\begin{center}
Obći teorem o prodolženju.
\end{center}

\emph{Vsako ireducibilno prědstavjenje od \(\mathfrak S\) v \(\Omega\), zadano homomorfizmom \(H_i\), možno na točno \(t\) različnyh načinov prodolžiti do prědstavjenij od \(\mathfrak T\).}

Dokaz jest analogny dokazu na str.~10--11. Imajemo
\[
\mathfrak S_\Omega=E_1\Omega+\cdots+E_s\Omega,\qquad
E_\nu\Omega=\mathfrak S_\Omega E_\nu,
\]
\[
\mathfrak T_\Omega\cdot E=\mathfrak T\cdot E\Omega
=\mathfrak T\mathfrak S\Omega=\mathfrak T\mathfrak S_\Omega,
\]
\[
\mathfrak T_\Omega\cdot E
=\mathfrak T(E_1\mathfrak S_\Omega+\cdots+E_s\mathfrak S_\Omega)
=E_1\mathfrak T_\Omega+\cdots+E_s\mathfrak T_\Omega,
\]
\[
[\mathfrak T_\Omega E_\nu:\Omega]
=[\mathfrak T_\Omega E_\nu:\Omega E_\nu]
=[\mathfrak T_\Omega E_\nu:E_\nu\mathfrak S_\Omega]
\leqq[\mathfrak T_\Omega\cdot E:\mathfrak S_\Omega]
=[\mathfrak T\cdot E:\mathfrak S]=t,
\]
\[
\sum[\mathfrak T_\Omega\cdot E_\nu:\Omega]=st
\]
zato
\[
[\mathfrak T_\Omega E_\nu:\Omega]=t.
\]
\(\mathfrak T_\Omega E_\nu\) razpada se zato na \(t\) idealov, iz ktoryh vsaky kako sumand od \(\mathfrak T_\Omega E\), zato takože od \(\mathfrak T_\Omega\), davaje prědstavjenje od \(\mathfrak T_\Omega\), ktoro, kako na str.~10--11, okazuje se prodolženjem prědstavjenja zadanogo posrědstvom \(H_\nu\). Ostala ireducibilna prědstavjenja od \(\mathfrak T_\Omega\) -- stvorjena sumandami, ne sodrživanymi v \(\mathfrak T_\Omega\cdot E_\nu\), od \(\mathfrak T_\Omega\) -- ne mogut byti prodolženja prědstavjenij, stvorjenyh posrědstvom \(H_\nu\), od \(\mathfrak S_\Omega\), zato že te idealy sut anulovane posrědstvom \(\mathfrak T_\Omega\cdot E_\nu\), a tym vyše posrědstvom \(\mathfrak S_\Omega\).

\section*{Glava III. Abelove grupy}\label{kapitel-iii.-abelsche-gruppen}
% END BOOK_S13

% BEGIN BOOK_S14
\subsection*{§ 14. Grupovo prsten}\label{der-gruppenring}

Nehaj \(\mathfrak G\) bude konečna grupa s elementami \(a_1,\ldots,a_h\), a \(\mathsf P\) komutativno polje. S grupovym množenjem kako množenjem obrazujemo hiperkompleksny sistem
\[
\mathfrak o[\mathfrak G]=a_1\mathsf P+\cdots+a_h\mathsf P
\]
\emph{grupovo prsten od \(\mathfrak G\) v \(\mathsf P\)}. (M. Z. § 6, Speiser, Theorie d. Gruppen end. Ord. § 48).

Pytajemo za prědstavjenja od \(\mathfrak o[\mathfrak G]\) v \(\mathsf P\) ili v algebraično zatvorjenom polju \(\Omega\) nad \(\mathsf P\). Kako už bylo spomenuto, ireducibilna prědstavjenja sut stvorjena prostymi lěvymi idealami od \(\mathfrak o[\mathfrak G]/\mathfrak C\). Zato nas -- kako v Galoisovoj teoriji -- interesuje pytanje, kogda \(\mathfrak o[\mathfrak G]\) imaje radikal. O tom imaje město obći teorem:

\begin{center}
\emph{\(\mathfrak o[\mathfrak G]\) jest togda i samo togda polno reducibilno, ako ręd \(h\) od \(\mathfrak G\) ne jest dělivy karakteristikoju \(p\) od \(\mathsf P\).}
\end{center}

Najprvo dokazujemo:
\[
\text{Iz } h\equiv0(p) \text{ slěduje } \mathfrak C\ne 0.
\]
\[
\mathfrak a=\mathsf P(a_1+\cdots+a_h)
\]
jest dvustranny ideal od \(\mathfrak o\). Zato že jest
\[
\mathsf P\cdot(a_1+\cdots+a_h)a
=\mathsf P\cdot(a_1a+\cdots+a_ha)
=\mathsf P(a_1+\cdots+a_h)=\mathfrak a
\]
i
\[
a\cdot\mathsf P(a_1+\cdots+a_h)
=\mathsf P(aa_1+\cdots+aa_h)
=\mathsf P(a_1+\cdots+a_h)=\mathfrak a
\]

zaradi grupovogo svojstva elementov \(a_i\). \(\mathfrak a\) očevidno ne jest nula, ale jest nilpotentny:
\[
\left(\sum_i a_i\right)^2=\sum_i\sum_j a_i a_j
=h\sum_j a_j=0.
\]

Obratno tvrdženje:
\[
\text{Iz } h\not\equiv0(p) \text{ slěduje } \mathfrak C=0,
\]
nyně dokazujemo samo za Abelove grupy. (V M. Z. § 26 ta čest teorema jest dokazana obće.)
% END BOOK_S14

% BEGIN BOOK_S15
\subsection*{§ 15. Grupove prsteny Abelovyh grup}\label{die-gruppenringe-abelscher-gruppen}

Pod \(\mathfrak A=\{a_1,\ldots,a_h\}\) razumějemo Abelovu grupu. Togda \(\mathfrak o\) jest komutativno, zato možno priměniti obće teoremy o komutativnyh hiperkompleksnyh sistemah.

\emph{Ako \(h\) ne jest dělivy karakteristikoju \(p\) od \(\mathsf P\), to \(\mathfrak o[\mathfrak A]\) imaje točno \(h\) različnyh ireducibilnyh prědstavjenij v \(\Omega\). (Ona sut prvogo stupenja, zato že \(\mathfrak o\) jest komutativno.)}

Iz toga slěduje, že \(\mathfrak o[\mathfrak A]/\mathfrak C\) imaje tojže rang kako \(\mathfrak o\), a tym samym \(\mathfrak C=0\); to jest vtora čest poprědnogo, ostavjenogo bez dokaza teorema za Abelove grupy.

\emph{Dokaz.} \(\mathfrak A\) jest prěmy produkt cikličnyh grup \(\mathfrak Z_i\) rędov \(h_i\)
\[
\mathfrak A=\mathfrak Z_1\times\cdots\times\mathfrak Z_t,\qquad
h=h_1h_2\cdots h_t .
\]
Ako \(z_i\) jest tvoreči element od \(\mathfrak Z_i\), to jemu odgovarjajuči element ireducibilnogo prědstavjenja (prvogo stupenja) mora byti \(h_i\)-ty korenj iz jedinice \(\varepsilon_\nu^{(h_i)}\). S drugoj strany, ako vsakomu \(z_i\) prirędimo \(h_i\)-ty korenj iz jedinice \(\varepsilon_\nu^{(h_i)}\), dobyva se ireducibilno prědstavjenje od \(\mathfrak o\) v \(\Omega\). Zato že nyně \(h\not\equiv0(p)\), a tym vyše \(h_i\not\equiv0(p)\), obstajajut točno \(h_1h_2\cdots h_t=h\) različne prědstavjenja od \(\mathfrak o\), tym jest dokazano.

(Takože jest legko viděti: ako \(h=0(p)\), to \(h_i=0(p)\) za prinajmenje jedno \(h_i\), obstajajut najvyše \(h_i/p\) različne \(\varepsilon_\nu^{(h_i)}\), zato menje než \(h\) prědstavjenij, zato \(\mathfrak o_\Omega\) mora sodrživati nenulovy radikal. Ale iz toga bezposrědnje ne slěduje ničto o radikalu od \(\mathfrak o\).)

Te \(h\) različne homomorfizmy od \(\mathfrak o_\Omega\) na podprsteny od \(\Omega\) označimo s \(\Theta_1,\ldots,\Theta_h\); vsakomu \(\Theta_i\) odpovědaje sumand \(e_i\Omega\) od \(\mathfrak o_\Omega\) kako modul prědstavjenja. Někdy priměnjujemo \(\Theta_i\) samo na \(\mathfrak A\) i togda govorimo o prědstavjenju od \(\mathfrak A\). Te homomorfizmy od \(\mathfrak A\) na grupy korenjev iz jedinice sut \emph{\(h\) karakterov od \(\mathfrak A\)}. Specialny karakter jest \emph{glavny karakter}, ktory vsakomu elementu iz \(\mathfrak A\) prirędžaje \(1\). Nehaj prinaležeči \(\Theta_i\) bude \(\Theta_1\).

Prědstavjenje glavnym karakterom možno za vsaku grupu, ne samo za Abelove grupy. Zato že to prědstavjenje už leži v \(\mathsf P\), v slučaju, kogda \(\mathfrak o\) jest polno reducibilno, \(\mathfrak o\) mora odcěpiti prosty sumand \(E\cdot\mathfrak o\), ktory davaje to prědstavjenje. \(E\) možno bezposrědnje izračunati. Mora byti \(E=\sum x_i a_i\). Prvo jest
\[
aE=E\cdot 1\quad\text{ili}\quad \sum x_i a a_i=\sum x_i a_i
\]
iz čego za \(a=a_j a_i^{-1}\)
\[
x_1a_ja_i^{-1}a_1+\cdots+x_i a_j+\cdots+x_ha_ja_i^{-1}a_h
=x_1a_1+\cdots+x_ha_h
\]
zato
\[
x_i=x_j
\]
slěduje.

Drugo jest:

\[
E=E^2=x^2\left(\sum a_i\right)^2=x^2h\sum a_i=xhE.
\]

V slučaju \(h=0(p)\) slěduje \(E=0\), t.~j.: \(\mathfrak o\) ne sodrživaje modul prědstavjenja za glavny karakter, zato ne moglo jest byti polno reducibilno. Tym jest ponovno dokazana jedna polovica teorema na str.~13.

Ako ale \(h\not\equiv0(p)\), to \(\mathfrak o\) jest polno reducibilno, neposrědnje slěduje

\[
x=\frac{1}{h},\qquad E=\frac{1}{h}(a_1+\cdots+a_h).
\]
% END BOOK_S15

% BEGIN BOOK_S16
\subsection*{§ 16. Relacije karakterov}\label{die-charakterenrelationen}

Fakt, že \(E=(1/h)\cdot\sum a_i\) davaje glavny karakter, neposrědnje vodi k relacijam karakterov. Mora byti

\[a \cdot e_i = e_i \cdot \Theta_i a\]

zato

\[e_i \cdot \Theta_i e_j = \begin{cases} e_i & i = j \\ 0 & i \neq j \end{cases}\]

ili

\[\Theta_i e_j = \begin{cases} 1 & i = j \\ 0 & i \neq j \end{cases}\]

V osoblivosti jest

\[\Theta_i e_1 = \begin{cases} 1 & i = 1 \\ 0 & i \neq 1 \end{cases}\]

ili, zato že \(e_1 = (1/h) \sum a_i\),

\[
\sum_\nu \Theta_i a_\nu=\begin{cases} h & i=1,\\ 0 & i\ne 1, \end{cases}
\]

To sut poznate relacije karakterov:

\emph{Suma jednogo karaktera po vsih elementah grupy jest nula, kromě glavnogo karaktera, za ktory ona jest h.}
% END BOOK_S16

% BEGIN BOOK_S17
\subsection*{§ 17. Galoisova teorija Abelovyh grup}\label{die-galoissche-theorie-abelscher-gruppen}

Za grupove prsteny Abelovyh grup možno ustanoviti teoremy podobne onym, ktore sut dokazivane v Galoisovoj teoriji komutativnyh polj. V Galoisovoj teoriji \(\Theta_i\Theta_1^{-1}\) tvorili sut grupu \(\mathfrak{G}\). Glavny teorem izražaje jednoznačno vzajemno priręđenje, osnovano na ustanovjenyh invariantnyh svojstvah, medžu podgrupami od \(\mathfrak{G}\) i podpoljami od \(Z\).

Homomorfizmy \(\Theta_i\) od \(\mathfrak{o}[\mathfrak{A}]\) možno sjediniti v grupu \(\mathsf A\), čije podgrupy stojęt v podobnom odnošenju k podgrupam od \(\mathfrak{A}\).

\emph{Ako produkt \(\Theta_i \Theta_k = \Theta\) definujemo formulom}

\[\Theta a = \Theta_i a \cdot \Theta_k a,\]

\emph{to množstvo vsih \(\Theta_i\) jest s \(\mathfrak A\) izomorfna grupa \(\mathsf A\).}

\emph{Dokaz.} \(\Theta = \Theta_i \cdot \Theta_k\) jest zaradi

\[\begin{aligned}
\Theta a \cdot \Theta b &= \Theta_i a \cdot \Theta_k a \cdot \Theta_i b \cdot \Theta_k b \\
&= \Theta_i a \cdot \Theta_i b \cdot \Theta_k a \cdot \Theta_k b \\
&= \Theta_i a b \cdot \Theta_k a b = \Theta a b
\end{aligned}\]

homomorfizm od \(\mathfrak A\) na \(\mathsf A\), a zato že kromě \(\Theta_i\) ne imaje drugyh homomorfizmov od \(\mathfrak A\), \(\Theta\) jest ravny nekomu \(\Theta_j\); zato \(\Theta_i\) naistino tvoręt grupu iz \(h\) elementov.

Nyně pod \(\varepsilon_{h_\nu}\) razumějemo primitivny \(h_\nu\)-ty korenj iz jedinice. Togda \(\Theta_i\) obrazuje tvoreči element \(z_\nu\) cikličnogo faktora \(\mathfrak Z_\nu\) od \(\mathfrak A\) na potenciju od \(\varepsilon_{h_\nu}\)

\[
\Theta_i z_\nu=\varepsilon_{h_\nu}^{\varrho_\nu^{(i)}}.
\]

Homomorfizmu \(\Theta_i\) prirędžajemo element \(\prod_\nu z_\nu^{\varrho_\nu^{(i)}}\) od \(\mathfrak A\). To jest jednoznačno vzajemno priręđenje od \(\mathsf A\) na \(\mathfrak A\), zato že \(\Theta_i\) jest oprěděljen posrědstvom \(\varrho_\nu^{(i)}\), različnym \(\Theta_i\) odpovědajut različne \(\prod_\nu z_\nu^{\varrho_\nu^{(i)}}\). Priręđenje jest izomorfno. Zato že s jednoj strany

\[
\Theta_i\Theta_j z_\nu=\Theta_i z_\nu\cdot \Theta_j z_\nu
=\varepsilon_{h_\nu}^{\varrho_\nu^{(i)}+\varrho_\nu^{(j)}},
\]

a s drugoj strany

\[
\prod_\nu z_\nu^{\varrho_\nu^{(i)}}\cdot
\prod_\nu z_\nu^{\varrho_\nu^{(j)}}=
\prod_\nu z_\nu^{\varrho_\nu^{(i)}+\varrho_\nu^{(j)}},
\quad\text{tym jest dokazano.}
\]

Nyně definujemo

\emph{Invariantnu oblast podgrupy \(\mathsf B\) od \(\mathsf A\):}

Podgrupa \(\mathfrak B\) elementov od \(\mathfrak A\), ktore vsaky \(\Theta\in\mathsf B\) obrazuje na 1.

\emph{Invariantnu grupu podgrupy \(\mathfrak B\) od \(\mathfrak A\):}

Podgrupa \(\mathsf B\) elementov od \(\mathsf A\), ktore obrazujut vsaky \(a\in\mathfrak B\) na 1.

i formulujemo glavny teorem:

\emph{Podgrupy \(\mathfrak B\) od \(\mathfrak A\) i podgrupy \(\mathsf B\) od \(\mathsf A\) možno jednoznačno vzajemno priręditi tako, že ako \(\mathfrak B\) i \(\mathsf B\) odpovědajut jedna drugoj, \(\mathsf B\) jest invariantna grupa od \(\mathfrak B\), a \(\mathfrak B\) jest invariantna oblast od \(\mathsf B\).}

Trěba jest zato pokazati:

\begin{enumerate}
\item \emph{Ako \(\mathsf B\) jest invariantna grupa od \(\mathfrak B\), to \(\mathfrak B\) jest invariantna oblast od \(\mathsf B\).}
\item \emph{Ako \(\mathfrak B\) jest invariantna oblast od \(\mathsf B\), to \(\mathsf B\) jest invariantna grupa od \(\mathfrak B\).}
\end{enumerate}

\paragraph{1.}
slěduje kako na str.~10 iz obćego teorema o prodolženju. Prědpokladamo, že \(\mathsf B\) jest invariantna grupa od \(\mathfrak B\), \(\mathfrak B'\) invariantna oblast od \(\mathsf B\), a \(\mathfrak B'\) sodrživaje \(\mathfrak B\). Elementi od \(\mathsf B\) sut prodolženja na cělu \(\mathfrak A\) jednogo homomorfizma od \(\mathfrak B\), ktory obrazuje \(\mathfrak B\) na 1. Ako te prodolženja priměnimo samo na \(\mathfrak B'\), iz obćego teorema o prodolženju slěduje, že ona razpadajut se na \(t\) klas tako, že homomorfizmy iste klasy sut identične, a homomorfizmy različnyh klas sut različne na \(\mathfrak B'\); pri tom \(t\) jest stupenj od \(\mathfrak o[\mathfrak B']\) nad \(\mathfrak o[\mathfrak B]\), to jest indeks od \(\mathfrak B\) v \(\mathfrak B'\), zato že \(\mathfrak o[\mathfrak B]\) i \(\mathfrak o[\mathfrak B']\) imajut tojže element jedinice, jmenno element jedinice od cěloj \(\mathfrak A\). Ale zato že \(\mathfrak B'\) jest invariantna oblast od \(\mathsf B\), mora byti \(t=1\), zato \(\mathfrak B'=\mathfrak B\), tym jest dokazano.

\textbf{2.} svodimo na 1.

Elementy \(a\) od \(\mathfrak A\) razgledamo homomorfizmami od \(\mathsf A\). Ustanovjujemo, že \(a\Theta\) jest ravno korenju iz jedinice \(\Theta a\), kratko \(a\Theta=\Theta a\). Tym \(a\) jest naistino oprěděljen kako homomorfizm od \(\mathsf A\), zato že
\[
a\Theta\cdot a\Theta'=\Theta a\cdot\Theta' a=\Theta\Theta' a=a\Theta\Theta'.
\]

Produkt dvoh homomorfizmov \(a,b\) jest obyčny grupovy produkt:
\[
\begin{array}{cccccc}
ab\Theta&=&a\Theta\cdot b\Theta&=&\Theta a\cdot\Theta b=\Theta ab&=ab\Theta\\
&&\text{produkt homomorfizmov}&&\text{grupovy produkt}&
\end{array}
\]

Elementi \(a,b\), različne kako grupove elementi, takože sut različne kako homomorfizmy: iz \(a\Theta=b\Theta\) za vse \(\Theta\) slěduje \(\Theta a=\Theta b\) za vse \(\Theta\), ili \(\Theta ab^{-1}=1\) za vse \(\Theta\). zato \(\Theta\) sut prodolženja na cělu \(\mathfrak A\) identičnogo homomorfizma posrědstvom \(ab^{-1}\) stvorjene cikličnoj podgrupy \(\mathfrak Z\) od \(\mathfrak A\). Po teoremu o prodolženju indeks od \(\mathfrak Z\) v \(\mathfrak A\) jest ravny \(h\), t.~j. \(\mathfrak Z=e,\ a=b\).

Nyně prěvodimo tvrdženja, učinjena v novom razuměnju,
\[
\begin{array}{ll}
\text{a)}& \mathfrak B\ \text{jest invariantna grupa od }\mathsf B.\\
\text{b)}& \mathsf B\ \text{jest invariantna oblast od }\mathfrak B.
\end{array}
\]
v staro razuměnje:

K a). \(\mathfrak B\) sostoji se iz vsih \(a\in\mathfrak A\), za ktore \(a\Theta=1\) za vse \(\Theta\in\mathsf B\), t.~j. \(\mathfrak B\) sostoji se iz vsih \(a\in\mathfrak A\), za ktore \(\Theta\cdot a=1\) za vse \(\Theta\in\mathsf B\). zato prěvod od a) jest
\[
\text{a')} \quad \mathfrak B\ \text{jest invariantna oblast od }\mathsf B.
\]

Odgovarjajuči prěvod od b) jest
\[
\text{b')} \quad \mathsf B\ \text{jest invariantna grupa od }\mathfrak B.
\]

zato prěvod tvrdženja už dokazanogo pod 1.:

Ako \(\mathfrak B\) jest invariantna grupa od \(\mathsf B\), to \(\mathsf B\) jest invariantna oblast od \(\mathfrak B\),
jest ravno tvrdženju, ktoro trěba dokazati:

Ako \(\mathfrak B\) jest invariantna oblast od \(\mathsf B\), to \(\mathsf B\) jest invariantna grupa od \(\mathfrak B\).

Ako relacije karakterov iz str.~15 zapišemo za \(\mathsf A\) vměsto za \(\mathfrak A\), dobyvajut se jednačenja
\[
\sum_{i=1}^{h} \Theta_i a_\nu =
\begin{cases}
h & \nu=1,\\
0 & \nu\ne 1.
\end{cases}
\]

Suma vsih karakterov nad grupovym elementom jest nula, kromě elementa jedinice, za ktory ona jest \(h\).

\section*{Glava IV. Dvustranno proste prsteny}\label{kapitel-iv.-zweiseitig-einfache-ringe}
% END BOOK_S17

% BEGIN BOOK_S18
\subsection*{§ 18. Pomoćna lema}\label{ein-hilfssatz}

Pod \(\mathsf K\) razumějemo tělo, a pod \(\mathfrak G\) grupu avtomorfizmov (izomorfnyh obrazovanij na sebe) od \(\mathsf K\). Množstvo elementov od \(\mathsf K\), ktore ostavajut neizměnjene pri vsakom avtomorfizmu naležečem k \(\mathfrak G\), jest podtělo \(\mathsf P\) od \(\mathsf K\), \emph{invariantno tělo od} \(\mathfrak G\).

Ako \(\tau\ne 0\) jest libovoljny element od \(\mathsf K\), avtomorfizm od \(\mathsf K\) dobyva se tako, že elementu od \(\mathsf K\), \(\alpha\), prirędimo element \(\tau^{-1}\alpha\tau\); po odgovarjajučem nazvanju iz teorije grup tako avtomorfizm nazyvamo \emph{vnutrnym avtomorfizmom}. Invariantno tělo \(\mathsf P\) grupy vsih vnutrnyh avtomorfizmov těla \(\mathsf K\) jest njegov centar.

Nyně nehaj budut zadane tělo \(\mathsf K\), grupa avtomorfizmov \(\mathfrak G\) od \(\mathsf K\) i njeno invariantno tělo \(\mathsf P\).

\[
\mathfrak M=x_1\mathsf P+\cdots+x_n\mathsf P
\]

nehaj bude \(\mathsf P\)-modul ranga \(n\). Dějstvo elementa iz \(\mathfrak G\), \(\gamma\), na elementy od \(\mathfrak M\) definujemo formulom

\[
\gamma\cdot x_i=x_i,\qquad\text{zato}\qquad
\gamma\sum \varrho_i x_i=\sum \varrho_i x_i.
\]

Poslědica toga jest, že \(\mathfrak G\) jest postala grupoju avtomorfizmov razširjenogo modula

\[
\mathfrak M_{\mathsf K}=x_1\mathsf K+\cdots+x_n\mathsf K
\]

a invariantna oblast od \(\mathfrak G\) v \(\mathfrak M_{\mathsf K}\) jest togda \(\mathfrak M\). Pri tyh prědpoloženjah imaje město teorem:

\begingroup\itshape
Ako
\[
\mathfrak C=z_1\mathsf K+\cdots+z_r\mathsf K
\]

jest podmodul od \(\mathfrak M_{\mathsf K}\), ktory vsaky element iz \(\mathfrak G\), \(\gamma\), obrazuje v njega samogo (ne potrěbno po elementah), to \(\mathfrak C\) jest razširjeny modul, obrazovany s \(\mathsf K\), nekogo podmodula \(\mathfrak c\) od \(\mathfrak M\).

\[
\mathfrak C=\mathfrak c_{\mathsf K}.
\]
\endgroup

\emph{Dokaz.} 1. Poprědna pripomněnka: Mora byti \(\mathfrak c=\mathfrak C\cap\mathfrak M\). Zato že obće imaje město
\[
\mathfrak c=\mathfrak c_{\mathsf K}\cap\mathfrak M.
\]

\(\mathfrak c\subseteq\mathfrak c_{\mathsf K}\cap\mathfrak M\) jest jasno. Za dokaz od \(\mathfrak c_{\mathsf K}\cap\mathfrak M\subseteq\mathfrak c\) izberemo bazu \(y_1,\ldots,y_r\) od \(\mathfrak c\),

\[
\mathfrak c=y_1\mathsf P+\cdots+y_r\mathsf P,
\]

togda jest

\[
\mathfrak c_{\mathsf K}=y_1\mathsf K+\cdots+y_r\mathsf K.
\]

Vsaky element iz \(\mathfrak c_{\mathsf K}\cap\mathfrak M\), \(c\), dopušča bazovy zapis kako element od \(\mathfrak c_{\mathsf K}\)
\[
c=\sum_{i=1}^{r}y_i\chi_i
\]
i bazovy zapis kako element od \(\mathfrak M\)

\[
c=\sum_{i=1}^{r}y_i\varrho_i+\sum_{j=r+1}^{n}y_j\varrho_j
\]

ako pod \(y_{r+1},\ldots,y_n\) razumějemo sistem elementov od \(\mathfrak M\), ktory dopolnjaje \(y_1,\ldots,y_r\) do bazy od \(\mathfrak M\) (M. Z. § 6).

Zato že oba zapisa možno razgledati bazovymi zapisami v \(\mathfrak M_{\mathsf K}\), oni sut identične, \(\chi_i=\varrho_i\ (i=1,\ldots,r)\), \(\varrho_j=0\ (j=r+1,\ldots,n)\), zato

\[
c=\sum_{i=1}^{r}y_i\varrho_i
\qquad\text{tym jest dokazano.}
\]

\textbf{2.} Da by dokazati razgledany teorem, konstruujemo bazu \(z_1,\ldots,z_r\) od \(\mathfrak C\), čiji elementi \(z_i\) naležet k \(\mathfrak M\). Togda, ako položimo

\[
\mathfrak c=z_1\mathsf P+\cdots+z_r\mathsf P
\]

naistino jest

\[
\mathfrak C=\mathfrak c_{\mathsf K}.
\]

Izberemo libovoljnu bazu \(y_1,\ldots,y_r\) od \(\mathfrak C\) i dopolnimo ju do bazy od \(\mathfrak M_{\mathsf K}\) odgovarjajučimi elementami \(x_{r+1},\ldots,x_n\) jednoj bazy \(x_1,\ldots,x_n\) od \(\mathfrak M\), ktora jest takože baza od \(\mathfrak M_{\mathsf K}\). To možno po M. Z. § 5.

\[
\mathfrak M_{\mathsf K}=\mathfrak C+x_{r+1}\mathsf K+\cdots+x_n\mathsf K.
\]

Iz toga za prvyh \(r\) elementov \(x_1,\ldots,x_r\) bazy od \(\mathfrak M\) slědujut jednačenja

\[
x_i=z_i-\sum_{j=r+1}^{n}x_j\chi_j^{(i)},\qquad
z_i\in\mathfrak C,\quad \chi_j^{(i)}\in\mathsf K.
\]

Zato že \(z_i\) zajedno s \(x_{r+1},\ldots,x_n\) očevidno tvoręt bazu od \(\mathfrak M_{\mathsf K}\), oni sut linearno nezavisne; zato \(z_1,\ldots,z_r\) jest baza od \(\mathfrak C\). \(z_1,\ldots,z_r\) jest baza iskanoj formy, čiji \(z_i\) naležet k \(\mathfrak M\):

Nehaj \(\gamma\) bude element iz \(\mathfrak G\). Togda s jednoj strany

\[
\begin{array}{rcl}
(1)\qquad \gamma z_i
&=&\gamma x_i+\displaystyle\sum_{j=r+1}^{n}\gamma x_j\cdot \gamma\chi_j^{(i)}
=x_i+\displaystyle\sum_{j=r+1}^{n}x_j\gamma\chi_j^{(i)}.
\end{array}
\]

a s drugoj strany
\[
\begin{array}{rcl}
(2)\qquad \gamma z_i&=&\displaystyle\sum_{\mu=1}^{r}z_\mu\sigma_\mu^{(i)},
\qquad \sigma_\mu^{(i)}\in\mathsf K,
\end{array}
\]

zato že \(\mathfrak C\) dopušča vse \(\gamma\in\mathfrak G\). (2) možno zapisati ješče v formě

\[
\begin{array}{rcl}
(3)\qquad \gamma z_i
&=&\displaystyle\sum_{\mu=1}^{r}x_\mu\sigma_\mu^{(i)}
+\displaystyle\sum_{j=r+1}^{n}x_j\sum_{\mu=1}^{r}\chi_j^{(\mu)}\sigma_\mu^{(i)}
\end{array}
\]

Sravnjenje koeficientov v (1) i (3) davaje

\[
\sigma_j^{(i)}=
\begin{cases}
1 & i=j,\\
0 & i\ne j.
\end{cases}
\]

Zato jest

\[
\gamma\cdot z_i=z_i.
\]

zato \(z_i\) naležet k invariantnoj oblasti od \(\mathfrak G\), t.~j. k \(\mathfrak M\), tym jest dokazano.
% END BOOK_S18

% BEGIN BOOK_S19
\subsection*{§ 19. Prědstavjenja dvustranno prostyh hiperkompleksnyh sistemov v nekomutativnyh tělah, ktore razširjajut jih polja koeficientov}\label{darstellungen-zweiseitig-einfacher-hyperkomplexer-systeme-in-nichtkommutativen-erweiterungskoerpern-ihrer-koeffizientenkoerper}

\[
\mathfrak S=x_1\mathsf P+\cdots+x_n\mathsf P
\]

nehaj bude dvustranno prosty, zato polno reducibilny hiperkompleksny sistem, a \(\mathsf K\) tělo s centrom \(\mathsf P\).

\textbf{Teorem 1.} \emph{\(\mathfrak S_{\mathsf K}\) jest dvustranno prosto (zato takože polno reducibilno).}

\emph{Dokaz} 1. Vsaky dvustranny \(\mathfrak S_{\mathsf K}\)-ideal \(\mathfrak a\) jest invariantny po odnošenju k grupě \(\mathfrak G\) vnutrnyh avtomorfizmov od \(\mathsf K\). Zato že ako, pod \(\chi\) razumějuči nenulovy element od \(\mathsf K\), položimo \(\chi^{-1}\mathfrak a\chi=\bar{\mathfrak a}\), to \(\bar{\mathfrak a}\subseteq\mathfrak a\) zaradi svojstva ideala od \(\mathfrak a\); ale zato že \(\bar{\mathfrak a}\), razgledany kako \(\mathsf K\)-modul, imaje nad \(\mathsf K\) tojže rang kako \(\mathfrak a\), mora byti \(\bar{\mathfrak a}=\mathfrak a\).

\textbf{2.} Po upravo ustanovjenoj pomoćnoj lemi dvustranny \(\mathfrak S_{\mathsf K}\)-ideal \(\mathfrak a\) mora zato byti razširjeny ideal nekogo podmodula \(\mathfrak A\) od \(\mathfrak S\), \(\mathfrak a=\mathfrak A_{\mathsf K}\). \(\mathfrak A\) jest dvustranny ideal od \(\mathfrak S\), zato že iz \(\mathfrak A=\mathfrak S\cap\mathfrak a\) dobyva se
\[
\mathfrak S\mathfrak A\subseteq\mathfrak S\mathfrak S\cap\mathfrak S\mathfrak a\subseteq\mathfrak S\cap\mathfrak a=\mathfrak A
\]
\[
\mathfrak A\mathfrak S\subseteq\mathfrak S\mathfrak S\cap\mathfrak a\mathfrak S\subseteq\mathfrak S\cap\mathfrak a=\mathfrak A.
\]

Ale zato že \(\mathfrak S\) sodrživaje samo dvustranne idealy 0 i \(\mathfrak S\), v \(\mathfrak S_{\mathsf K}\) takože obstajajut samo dva dvustranna ideala 0 i \(\mathfrak S_{\mathsf K}\), tym jest dokazano.

Nyně razgledajemo ireducibilna prědstavjenja od \(\mathfrak S\) v \(\mathsf K\), i hočemo ustanoviti recipročna prědstavjenja; to imaje formalnu prědnost. V konečnom jest vsejedno, izbirajemo li prěma ili recipročna prědstavjenja, zato že ona sut jednoznačno vzajemno priręđena.

Nehaj \(\mathfrak M\) bude recipročny modul prědstavjenja od \(\mathfrak S\) v \(\mathsf K\), to jest lěvy \(\mathfrak S\)-modul i lěvy \(\mathsf K\)-modul konečnogo ranga,
\[
\mathfrak M=\mathsf K\cdot y_1+\cdots+\mathsf K\cdot y_r
\]
s zakonom komutativnosti
\[
s(xm)=x(sm).
\]

Dokazujemo, že \(\mathfrak M\) možno takože razgledati kako lěvy \(\mathfrak S_{\mathsf K}\)-modul. Najprvo trěba objasniti množenje elementa od \(\mathfrak S_{\mathsf K}\) elementom od \(\mathfrak M\). Ako \(s\in\mathfrak S,\ x\in\mathsf K,\ m\in\mathfrak M\), ustanovjujemo
\[
sx\cdot m=s\cdot(x\cdot m)
\]
i analogično za obći \(\mathfrak S_{\mathsf K}\)-element \(\sum s_ix_i\)
\[
\sum s_ix_i\cdot m=\sum s_i\cdot(x_i\cdot m).
\]

Možno uvěriti se, že ta definicija jest jednoznačna i ne zavisi od načina, na ktory element od \(\mathfrak S_{\mathsf K}\) zapisujemo kako sumu produktov \(s_ix_i\).

Trěba zato iz
\[
\begin{array}{rcl}
(1)\qquad \displaystyle\sum_{i=1}^{p}s_ix_i&=&0
\end{array}
\]
zaključiti
\[
\sum_{i=1}^{p}s_i(x_i m)=0
\]
Prědpokladamo, že \(s_1,\ldots,s_q\) sut linearno nezavisne i že \(s_{q+1},\ldots,s_p\) možno izraziti posrědstvom \(s_1,\ldots,s_q\)
\[
s_j=\sum_{i=1}^{q}s_i\varrho_i^{(j)},\qquad j=q+1,\ldots,p.
\]

Koeficienty morajut ležati v \(\mathsf P\), zato že \(s_j\) možno izraziti samo na jeden način, a zapis jest už možny v \(\mathfrak S\). Iz (1) dobyva se

\[
x_i+\sum_{j=q+1}^{p}\varrho_i^{(j)}x_j=0,\qquad i=1,\ldots,q.
\]

zato jest

\[
\begin{aligned}
\sum_{i=1}^{p}s_i(x_i m)
&=s_1(x_1m)+\cdots+s_q(x_qm)
  +\left(\sum_i s_i\varrho_i^{(q+1)}\right)(x_{q+1}m)\\
&\quad+\cdots+\left(\sum_i s_i\varrho_i^{(p)}\right)(x_pm)\\
&=s_1(x_1m)+\cdots+s_q(x_qm)\\
&\quad+s_1\varrho_1^{(q+1)}(x_{q+1}m)+\cdots+s_q\varrho_q^{(q+1)}(x_{q+1}m)
  \qquad\text{Zaradi svojstva modula}\\
&\quad+\cdots\\
&\quad\cdots\\
&\quad+s_1\varrho_1^{(p)}(x_pm)+\cdots+s_q\varrho_q^{(p)}(x_pm)\\
&=s_1(x_1m)+\cdots+s_q(x_qm)\\
&\quad+s_1(\varrho_1^{(q+1)}x_{q+1}m)+\cdots+s_q(\varrho_q^{(q+1)}x_{q+1}m)\\
&\quad+\cdots\\
&\quad\cdots\\
&\quad+s_1(\varrho_1^{(p)}x_pm)+\cdots+s_q(\varrho_q^{(p)}x_pm)
  \qquad\begin{gathered}\text{Zaradi zakona}\\ \text{asociativnosti }s'(sm)\\ =s's\cdot m\text{ i}\\ x'(xm)=x'x\cdot m\end{gathered}\\
&=s_1\left(x_1m+\varrho_1^{(q+1)}x_{q+1}\cdot m+\cdots+\varrho_1^{(p)}x_p\cdot m\right)+\cdots\\
&=s_1\left(\left(x_1+\varrho_1^{(q+1)}x_{q+1}+\cdots+\varrho_1^{(p)}x_p\right)m\right)+\cdots\\
&=0.
\end{aligned}
\]


Iz svojstv \(\mathfrak S_{\mathsf K}\)-modula od \(\mathfrak M\) distributivne zakony \((S+S')m=Sm+S'm,\ S(m+m')=Sm+Sm'\) sut jasne. Ješče trěba obratiti vnimanje na zakon asociativnosti \(S(S'm)=SS'\cdot m\); zato že distributivne zakony sut pravdive, možno ograničiti se na slučaj

\[
sx\cdot(s'x' m)=sx(s'x'\cdot m)
\]

S jednoj strany, po definiciji od \(sx\cdot m\), jest

\[
sx\cdot(s'x'\cdot m)=s\cdot(x\cdot(s'\cdot(x'\cdot m)))
\]

a s drugoj strany

\[
sx s'x'\cdot m=ss'xx'\cdot m=ss'\cdot(xx'\cdot m)
=s\cdot(s'\cdot(x\cdot(x'\cdot m)))
\]

zaradi dvoh relacij komutativnosti \(=s\cdot(x\cdot(s'\cdot(x'\cdot m)))\)

\[
x\cdot s'=s'\cdot x
\]
\[
s'\cdot(x\cdot m)=x\cdot(s'\cdot m)
\]

zato \(\mathfrak M\) naistino jest lěvy \(\mathfrak S_{\mathsf K}\)-modul, očevidno konečny. Obratno, vsaky konečny lěvy \(\mathfrak S_{\mathsf K}\)-modul jest recipročny modul prědstavjenja od \(\mathfrak S\) v \(\mathsf K\), zato že očevidno jest konečny \(\mathsf K\)-modul i zaradi

\[
s(x\cdot m)=sx\cdot m=xs\cdot m=x(s\cdot m)
\]

takože izpolnja uslovje komutativnosti. (Sravni takože prostějši dokaz v § 24).

Nyně po M.Z. § 18 vsaky konečny \(\mathfrak S_{\mathsf K}\)-modul jest suma prostyh modulov, ktore sut operatorno izomorfne s prostymi lěvymi idealami od \(\mathfrak S_{\mathsf K}\) po odnošenju k operatoram iz \(\mathfrak S_{\mathsf K}\). Obratno, vsaky lěvy ideal od \(\mathfrak S_{\mathsf K}\) jest modul prědstavjenja od \(\mathfrak S\) v \(\mathsf K\), a zato že vsi lěvi idealy od \(\mathfrak S_{\mathsf K}\) sut operatorno izomorfne, imaje město

\textbf{Teorem 2.} \emph{Jedina klasa ireducibilnyh recipročnyh prědstavjenij od \(\mathfrak S\) v \(\mathsf K\) jest ona, ktora jest stvorjena prostym lěvym idealom od \(\mathfrak S_{\mathsf K}\).}

Ako to ireducibilno prědstavjenje imaje stupenj \(r\), to stupenj libovoljnogo prědstavjenja jest kratny od \(r\); očevidno možno nadovezivanjem ireducibilnyh prědstavjenij za zadany kratny \(rs\) od \(r\) konstruovati prědstavjenje \(rs\)-togo stupenja. Tomu odpovědaje obrazovanje modula prědstavjenja, ktory jest prěma suma od \(s\) prostyh modulov.

Nyně už znajemo jedno prědstavjenje od \(\mathfrak S\) v \(\mathsf K\), jmenno ono, ktoro jest stvorjeno samym \(\mathfrak S\) kako modulom prědstavjenja (glavno prědstavjenje, M. Z. § 20); ono už leži v \(\mathsf P\). Njegov stupenj jest \(n\), zato \(n\) jest dělivy posrědstvom \(r\):

\[
n=rt.
\]

Kvocient \(t\) jest legko oprěděliti: zato že \(r\) jest rang jednogo lěvogo ideala nad \(\mathsf K\), a \(n\) rang cělokupnogo \(\mathfrak S_{\mathsf K}\), \(t\) jest čislo lěvyh idealov od \(\mathfrak S_{\mathsf K}\), zato \(t^2\) jest rang od \(\mathfrak S_{\mathsf K}\) nad tělom avtomorfizmov njegovyh prostyh idealov.

Specializujemo naše razgledanje na slučaj, kogda \(\mathfrak S\) jest komutativno polje.

\textbf{Teorem 3.} \emph{Ako \(\mathfrak S\) jest komutativno polje, to \(\mathfrak S\) jest centar od \(\mathfrak S_{\mathsf K}\).}

\emph{Dokaz.} Zato že vsaky element od \(\mathfrak S\) jest komutativen s vsakym elementom od \(\mathfrak S\) i s vsakym elementom od \(\mathsf K\), zato takože s vsakym elementom od \(\mathfrak S_{\mathsf K}\), \(\mathfrak S\) naleži k centru od \(\mathfrak S_{\mathsf K}\). Element od \(\mathfrak S_{\mathsf K}\) jest komutativen s vsakym elementom od \(\mathfrak S_{\mathsf K}\) samo togda, ako jest komutativen prinajmenje s vsakym elementom od \(\mathsf K\), to jest ako njegovi koeficienty naležet k centru od \(\mathsf K\), ktory jest \(\mathsf P\), to jest ako on leži v \(\mathfrak S\), tym jest dokazano.

Zato že pri komutativnyh prstenah ne imaje razlike medžu recipročnym i prěmym izomorfizmom, vsako recipročno prědstavjenje možno razgledati takože prěmym. Ireducibilno prědstavjenje od \(\mathfrak S\) v \(\mathsf K\) mora zato takože byti dano prěmym modulom prědstavjenja. Toj prěmy modul prědstavjenja jest prosty pravy ideal od \(\mathfrak S_{\mathsf K}\). Zato že \(\mathfrak r\) jest konečny pravy \(\mathsf K\)-modul ranga \(r\) i zaradi \(\mathfrak r\mathfrak S=\mathfrak S\mathfrak r\) takože lěvy \(\mathfrak S\)-modul (samorazumno s smješanym zakonom asociativnosti), zato prěmy modul prědstavjenja ranga \(r\); a zato že obstaja samo jedna klasa ireducibilnyh prědstavjenij, on mora stvarjati dano prědstavjenje.
% END BOOK_S19

% BEGIN BOOK_S20
\subsection*{§ 20. Nekomutativne těla}\label{nichtkommutative-koerper}

Nyně pod

\[
\mathfrak K=y_1\mathsf P+\cdots+y_m\mathsf P
\]

razumějemo nekomutativno tělo s centrom \(\mathsf P\). Za \(\mathfrak K\) razvijemo ręd teoremov, analognyh teoremam iz teorije komutativnyh konečnyh algebraičnyh razširjenij polj.

Ako \(Z=\xi_1\mathsf P+\cdots+\xi_n\mathsf P\) označaje hiperkompleksny sistem nad \(\mathsf P\), dva razširjena sistema \(\mathfrak K_Z\) i \(Z_{\mathfrak K}\) sut identične, zato že za \(\mathfrak K_Z\) možno zapisati

\[
\mathfrak K_Z=y_1\xi_1\mathsf P+\cdots+y_m\xi_n\mathsf P
\]

a za \(Z_{\mathfrak K}\)

\[
Z_{\mathfrak K}=\xi_1y_1\mathsf P+\cdots+\xi_ny_m\mathsf P
\]

ale oba izraza sut identične zaradi \(y_i\xi_j=\xi_jy_i\).

V slědujućem \(Z\) vsečasno bude komutativno polje, zato možno priměniti teoremy iz § 19. Po teoremu 1 iz § 19 \(\mathfrak K_Z\) jest vsečasno dvustranno prosto.

\textbf{Teorem 1.} \emph{Nehaj \(\Omega\) bude algebraično zatvorjeno razširjeno polje od \(\mathsf P\). Togda centar od \(\mathfrak K_\Omega\) jest ravny \(\Omega\).}

\emph{Dokaz.} Oslanjamo se na odgovarjajuči teorem 3 iz § 19. Jasno jest, že \(\Omega\) leži v centru od \(\mathfrak K_\Omega\). Ako \(\alpha\) jest libovoljny element centra od \(\mathfrak K_\Omega\), njegovi koeficienty naležet k nekomu konečnomu algebraičnomu razširjenomu polju \(Z\) od \(\mathsf P\) (jmenno k polju, ktoro oni stvarjajut). Samorazumno \(\alpha\) takože mora naležati k centru od \(\mathfrak K_Z\). Ale toj centar jest \(Z\), zato \(\alpha\) naleži k \(Z\), a zato takože k \(\Omega\), tym jest dokazano.

Na sovsěm podobny način dokazujemo

\textbf{Teorem 2.} \emph{\(\mathfrak K_\Omega\) jest dvustranno prosto.}

\emph{Dokaz.} Nehaj \(\mathfrak a=z_1\Omega+\cdots+z_r\Omega\) bude dvustranny ideal. Koeficienty od \(z_1,\ldots,z_r\) naležet k nekomu konečnomu razširjenomu polju \(Z\) od \(\mathsf P\); v \(\mathfrak K_Z\) jest \(\mathfrak a'=z_1Z+\cdots+z_rZ\) dvustranny ideal, a zato že \(\mathfrak K_Z\) jest dvustranno prosto, \(\mathfrak a'\) jest ili nula ili cělo \(\mathfrak K_Z\); zato \(\mathfrak a\) jest ili 0 ili cělo \(\mathfrak K_\Omega\), t.~j. \(\mathfrak K_\Omega\) jest dvustranno prosto, tym jest dokazano.

Po tom teoremu \(\mathfrak K_\Omega\) jest matricno prsten nad nekym razširjenym poljem svojego centra \(\Omega\). To polje kako podprsten cělogo \(\mathfrak K_\Omega\) mora iměti konečny stupenj nad \(\Omega\), zato mora byti ravno \(\Omega\):

\textbf{Teorem 3.} \emph{\(\mathfrak K_\Omega\) jest matricno prsten nad \(\Omega\),}

\[
\mathfrak K_\Omega=\sum\Omega c_{ik}
\]

\emph{zato stupenj od \(\mathfrak K\) nad njegovym centrom \(\mathsf P\) jest kvadratno čislo \(m=t^2\).}

Čislo \(t\) pravyh (ili lěvyh) idealov, na ktore razpada se \(\mathfrak K_\Omega\), nazyvajmo kratko absolutnym čislom komponentov od \(\mathfrak K\).

Svojstvo od \(\Omega\), že \(\mathfrak K_\Omega\) razpada se na \(t\) prostyh idealov, očevidno imajut už konečna razširjenja \(Z\) od \(\mathsf P\), napr. taka \(Z\), ktore se dobyvajut priloženjem koeficientov od \(c_{ik}\) k \(\mathsf P\). Nyně razgledajemo ta razpadna polja.

\textbf{Definicija.} Konečno algebraično razširjeno polje \(Z\) od \(\mathsf P\) nazyvaje se \emph{razpadnym poljem} od \(\mathfrak K\), ako čislo prostyh pravyh idealov, na ktore razpada se \(\mathfrak K_Z\), už jest ravno absolutnomu čislu komponentov \(t\)\srcfn{1}{Sravni pojem razpadnogo polja na str.~9}.

\textbf{Teorem 4.} \emph{Ako \(Z\) jest razpadno polje od \(\mathfrak K\), to \(Z\) jest tělo avtomorfizmov pravyh idealov od \(\mathfrak K_Z\); obratno, ako tělo avtomorfizmov pravyh idealov od \(\mathfrak K_Z\) jest ravno \(Z\), to \(Z\) jest razpadno polje.}

\emph{Dokaz.} 1. Nehaj \(Z\) bude razpadno polje, a \(\mathfrak T\) tělo avtomorfizmov pravyh idealov od \(\mathfrak K_Z\), to jest

\[
\mathfrak K_Z=\sum\mathfrak T c_{ik}
\]

Iz toga slěduje

\[
\mathfrak K_\Omega=\sum\mathfrak T_\Omega c_{ik}
\]

rang od \(\mathfrak K_\Omega\) nad \(\Omega\) jest zato produkt od \(t^2\) s rangom od \(\mathfrak T_\Omega\); s drugoj strany on jest ravny \(t^2\), zato rang od \(\mathfrak T_\Omega\), a takože od \(\mathfrak T\), jest ravny 1, \(\mathfrak T=Z\), tym jest dokazano.

\textbf{2.} Ako \(Z\) jest tělo avtomorfizmov prostyh idealov od \(\mathfrak K_Z\), to jest

\[
\mathfrak K_Z=\sum Zc_{ik}
\]

zato

\[
\mathfrak K_\Omega=\sum\Omega c_{ik}
\]

čislo elementov \(c_{ik}\) mora zato byti ravno \(t^2\), tym jest dokazano.

\textbf{Teorem 5.} \emph{Ireducibilno prědstavjenje \(\mathfrak Z\) razpadnogo polja \(Z\) od \(\mathfrak K\) v \(\mathfrak K\) -- po teoremu 2 na str.~22 obstaja samo jedna klasa ireducibilnyh prědstavjenij -- jest maksimalno komutativno podpolje matricnogo prstena \(\mathfrak K_r\) (\(r\) jest stupenj ireducibilnogo prědstavjenja od \(Z\)).}

\emph{Ako obratno ireducibilno prědstavjenje \(\mathfrak Z\) komutativnogo razširjenogo polja \(Z\) od \(\mathsf P\) jest maksimalno komutativno podpolje od \(\mathfrak K_r\), to \(Z\) jest razpadno}

\emph{polje.}

\emph{Dokaz.} 1. Ako \(n\) jest stupenj razpadnogo polja \(Z\) nad \(\mathsf P\), to stupenj \(r\) ireducibilnogo prědstavjenja \(\mathfrak Z\) od \(Z\) jest ravny \(n/t\) (str.~22). Nehaj nyně \(\mathfrak Z\subseteq\mathfrak Z^*\subseteq\mathfrak K_r\), a \(\mathfrak Z^*\) bude maksimalno komutativno polje stupenja \(n^*\geq n\) nad \(\mathsf P\). Trěba dokazati \(\mathfrak Z^*=\mathfrak Z\). Vměsto toga dostatočno jest dokazati \(n^*=n\). Ako \(r^*\) označaje stupenj ireducibilnogo prědstavjenja s \(\mathfrak Z^*\) izomorfnogo razširjenogo polja \(Z^*\) od \(Z\), ponovno jest \(n^*=r^*t\); s drugoj strany \(r\) jest dělivy posrědstvom \(r^*\), zato že v \(\mathfrak Z^*\) jest dano prědstavjenje \(r\)-togo stupenja od \(Z^*\), zato \(r^*\leq r\). Iz

\[
n\leq n^*=r^*t\leq rt=n \qquad\text{slěduje}\qquad n^*=n.
\]

2. Nehaj ireducibilno prědstavjenje \(\mathfrak Z\) \(r\)-togo stupenja komutativnogo razširjenogo polja \(Z\) od \(\mathsf P\) v \(\mathfrak K\) bude maksimalno komutativno podpolje od \(\mathfrak K_r\). Nehaj tělo avtomorfizmov prostyh idealov od \(\mathfrak K_Z\) bude \(\mathfrak T\), to jest

\[
\mathfrak K_Z=\sum_{1}^{t'}\mathfrak T c_{ik}.
\]

\(\mathfrak T\) sodrživaje \(Z\). Trěba dokazati \(\mathfrak T=Z\). Nyně prosty pravy ideal \(\mathfrak r\) od \(\mathfrak K_r\) ne jest samo modul prědstavjenja od \(Z\) v \(\mathfrak K\), ale takože modul prědstavjenja od \(\mathfrak T\), zato že jest

\[
\mathfrak r=\mathfrak T c_{i1}+\cdots+\mathfrak T c_{ik},
\]

zato \(\mathfrak T\mathfrak r=\mathfrak r\); \(\mathfrak r\) jest lěvy \(\mathfrak T\)-modul i očevidno pravy \(\mathfrak K\)-modul.

Obstaja zato s \(\mathfrak T\) izomorfno podtělo \(\overline{\mathfrak T}\) od \(\mathfrak K_r\), ktoro sodrživaje \(\mathfrak Z\). Ako by \(\mathfrak T\) bylo različno od \(Z\), to \(\overline{\mathfrak T}\) različno od \(\mathfrak Z\), priloženje ne naležečego k \(\mathfrak Z\) elementa od \(\overline{\mathfrak T}\), \(\alpha\), k \(\mathfrak Z\) dalo by pravo komutativno razširjeno polje od \(\mathfrak Z\) v \(\mathfrak K_r\), a to protivrěči maksimalnosti od \(\mathfrak Z\), tym jest dokazano.

Iz teorema 5 bezposrědnje slěduje

\textbf{Teorem 6.} \emph{Maksimalna komutativna podpolja od \(\mathfrak K\) imajut stupenj \(t\) nad \(\mathsf P\) i sut razpadna polja\srcfn{1}{Prědpoloženje, že vsako razpadno polje takože sodrživaje neko maksimalno komutativno podpolje od \(\mathfrak K\), ne jest pravilno; mogut obstajati razpadna polja bez toga svojstva, a odgovarjajuči \(r\) sigurno jest vyše od 1. Vidi: \foreign{R. Brauer} i \foreign{E. Noether}, \foreign{Über minimale Zerfällungskörper irreduzibler Darstellungen; Sitzungsberichte d. preuss. Akademie der Wisssch. 1927 XXXII, str.~221.}}.}

Nyně razvijemo teorem o izomorfizmu podprstenov jednogo \(\mathfrak K_r\), na ktorom osnovana jest podrobnėjša znajemost struktury Galoisovoj grupy od \(\mathfrak K\) i s tym svęzanyh faktov.

\textbf{Teorem 7.} \emph{Nehaj \(\mathfrak S_1\) i \(\mathfrak S_2\) budut dvustranno proste, \(\mathsf P\) sodrživajuće podprsteny matricnogo prstena \(\mathfrak K_r\) nad \(\mathfrak K\), i nehaj medžu nimi obstaja izomorfizm, ostavjajuči \(\mathsf P\) po elementah nepodvižnym, \(\mathfrak S_1\cong\mathfrak S_2\). Togda toj izomorfizm jest transformacija regularnym \(\mathfrak K_r\)-elementom \(\tau\).}

\emph{Ako \(\sigma_1\in\mathfrak S_1\) i \(\sigma_2\in\mathfrak S_2\) sut vzajemno priręđene, to jest}

\[
\tau^{-1}\sigma_1\tau=\sigma_2.
\]

\emph{Dokaz.} Konstruujemo s \(\mathfrak S_1\) i \(\mathfrak S_2\) recipročno izomorfno prsten \(\Sigma\), čije podpolje izomorfno s \(\mathsf P\) identifikujemo s \(\mathsf P\). \(\mathfrak S_1,\mathfrak S_2\) sut togda dva recipročna prědstavjenja jednakogo stupenja \(r\) od \(\Sigma\) v \(\mathfrak K\). Zato že po teoremu 2 na str.~22 obstaja samo jedna klasa ireducibilnyh recipročnyh prědstavjenij, zato takože samo jedna klasa prědstavjenij zadanogo stupenja od \(\Sigma\) v \(\mathfrak K\) (kratny od ireducibilnoj klasy), \(\mathfrak S_1\) i \(\mathfrak S_2\) morajut naležati k istoj klase; zato jedno nastava iz drugogo transformacijoju, tym jest dokazano.

Iz teorema 7 neposrědnje slěduje

\textbf{Teorem 8.} \emph{Grupa avtomorfizmov, ktore ostavjajut \(\mathsf P\) po elementah nepodvižnym, od \(\mathfrak K\) -- t.~j. Galoisova grupa \(\mathfrak G\) od \(\mathfrak K\) po odnošenju k njegovomu centru -- jest grupa vnutrnyh avtomorfizmov od \(\mathfrak K\).}

V teoremu 7 trěba samo položiti \(r=1\) i \(\mathfrak S_1=\mathfrak S_2=\mathfrak K\).

Slědujuća dva teorema pokazyvajut upotrěbljivost dosěga razvitoj obćej teorije za razgledanje specialnyh pytanij.

\textbf{Teorem 9.} \emph{Jedino nekomutativno tělo konečnogo stupenja, čiji centar jest polje \(R\) realnyh čisl, jest tělo kvaternionov.}

\emph{Dokaz.} Nehaj \(\mathfrak K\) bude nekomutativno tělo s centrom \(R\) (konečnogo ranga nad \(R\)). Po teoremu 6 stupenj od \(\mathfrak K\) nad \(R\) jest kvadrat stupenja \(t\) nekogo maksimalnogo komutativnogo podpolja od \(\mathfrak K\); a zato že ono može byti samo polje izomorfno s poljem kompleksnyh čisl, jest \(t=2\), a \(\mathfrak K\) imaje stupenj 4 nad \(R\). Nehaj \(\mathfrak Z_1\) bude maksimalno komutativno podpolje od \(\mathfrak K\); togda možno položiti \(\mathfrak Z_1=R(i)\), kde \(i^2=-1\). Ale togda takože \(\mathfrak Z_2=R(-i)\) jest -- kako množstvo sběgajuće sę s \(\mathfrak Z_1\) -- polje izomorfno s \(\mathfrak Z_1\), a izo-

morfizm jest zadany posrědstvom \(i\mapsto -i\). Po teoremu 7 obstaja zato element \(j^*\in\mathfrak K\) s

\[
j^{*-1}ij^*=-i.
\]

Rang od \(\mathfrak Z_1(j^*)\) nad \(\mathfrak Z_1\) jest nyně prinajmenje 2, zato že \(j^*\) ne može naležati k \(\mathfrak Z_1\). zato \(\mathfrak Z_1(j^*)\) imaje rang 4 nad \(R\), jest ravno \(\mathfrak K\), i možno položiti:

\[
\mathfrak K=R+Ri+Rj^*+Rij^*.
\]

Zaradi

\[
\begin{aligned}
j^{*-2}1j^{*2}&=1,\\
j^{*-2}ij^{*2}&=i,\\
j^{*-2}j^*j^{*2}&=j^*,\\
j^{*-2}ij^*j^{*2}&=ij^*
\end{aligned}
\]

\(j^{*2}\) jest komutativen s vsimi elementami od \(\mathfrak K\), naleži k centru \(R\), a \(j^{*2}\) jest realno čislo \(r\). \(r\) ne može byti pozitivno, zato že inače polinom \(x^2-r\) v komutativnom polju \(R(j^*)\) imel by četyri korenja \(+\sqrt r,-\sqrt r,+j^*\) i \(-j^*\). zato \(j^{*2}=-m^2,\ m\in R\). Položimo \(j=j^*m^{-1}\); togda \(j^2=-1\), a \(1,i,j,k=ij\) sut jedinice kvaternionov, tym jest dokazano.

Dokaz očevidno imaje město za vsako algebraično realno zatvorjeno polje vměsto \(R\).

\textbf{Teorem 10 (teorem od Vedderburna).} \emph{Vsako tělo iz konečnogo čisla elementov jest komutativno.}

\emph{Dokaz.} Nehaj \(\mathfrak K\) bude tělo s konečnym čislom elementov, a \(\mathsf P\) jego centar. Nehaj \(t^2\) bude stupenj od \(\mathfrak K\) nad \(\mathsf P\). Vsa maksimalna komutativna podpolja od \(\mathfrak K\) imajut stupenj \(t\) nad \(\mathsf P\), zato imajut ravnako mnogo elementov; po poznatom teoremu iz teorije Galoisovyh polj (teorem od Moorea) ona sut izomorfne, a te izomorfizmy, zato že vsako Galoisovo polje jest Galoisovo razširjenje svojego prostogo polja, možno izbrati tako, že vsaky element od \(\mathsf P\) odpovědaje samomu sebe. Po teoremu 7 vsa maksimalna komutativna podpolja od \(\mathfrak K\) nastavajut zato iz jednogo \(\mathfrak Z\) transformacijoju. Ale \(\mathfrak K\) jest množstvo sjedinenja vsih svojih maksimalnyh komutativnyh podpolj. (Priloženje \emph{jednogo} elementa k \(\mathsf P\) stvarjaje komutativno polje). Ako nyně razgledamo samo multiplikativnu grupu \(\mathfrak K^*\) od \(\mathfrak K\) (\(\mathfrak K\) bez nuly), dostali jesmo, že \(\mathfrak K^*\) jest množstvo sjedinenja podgrup konjugovanyh podgrupě \(\mathfrak Z^*\). Ale v slučaju \(t>1\) to davaje protivrěčnost. Zato že ako \(\mathfrak K^*\) sodrživaje \(M\) elementov, \(\mathfrak Z^*\) \(N\) elementov, a \(L\) jest indeks od \(\mathfrak Z^*\) v \(\mathfrak K^*\), to imaje město \(M=NL\), a čislo različnyh konjugovanyh podgrup jest najvyše \(L\). Čak ako to čislo bylo by ravno \(L\), da by vse konjugovane podgrupy zajedno izčrpale cělu grupu, nikake dvě iz njih ne směli by iměti obći element. Ale vse imajut obći element jedinice.
% END BOOK_S20

% BEGIN BOOK_S21
\subsection*{§ 21. Galoisova teorija nekomutativnyh těl}\label{die-galoissche-theorie-der-nichtkommutativen-koerper}

Pod \(G\) razumějemo Galoisovu grupu od \(\mathfrak K\) po odnošenju k njegovomu centru \(\mathsf P\). Po teoremu 8 \(G\) jest izomorfna s \(\overline G=\mathfrak K^*/\mathsf P^*\); razgledamo toj izomorfizm ustanovjenym. Podgrupu \(H\) od \(G\) nazyvamo \emph{zatvorjenoj}, ako posrědstvom izomorfizma \(G\simeq\mathfrak K^*/\mathsf P^*\) jej priręđena podgrupa \(\mathfrak H^*\) od \(\mathfrak K^*\) po dodanju \emph{nuly} postavaje tělom.


\textbf{Teorem 11. Glavny teorem Galoisovoj teorije.} \emph{Zatvorjene podgrupy \(H\) od \(G\) i těla \(\mathfrak S\) medžu \(\mathsf P\) i \(\mathfrak K\) možno jednoznačno vzajemno priręditi tako, že ako \(H\) i \(\mathfrak S\) odpovědajut jedno drugomu, \(H\) jest invariantna grupa od \(\mathfrak S\), a \(\mathfrak S\) invariantno tělo od \(H\).}

\emph{Dokaz.} Razděljamo naš teorem na tri čestne teoremy
\begin{enumerate}
\item Invariantna grupa \(H\) těla \(\mathfrak S\) jest zatvorjena.
\item Ako \(\mathfrak S\) imaje invariantnu grupu \(H\), to \(\mathfrak S\) jest invariantno tělo od \(H\).
\item Ako \(H\) imaje invariantno tělo \(\mathfrak S\), to \(H\) jest invariantna grupa od \(\mathfrak S\).
\end{enumerate}

\begin{enumerate}
\item Invariantnoj grupě \(H\) těla \(\mathfrak S\) posrědstvom izomorfizma \(G\simeq \mathfrak K^*/\mathsf P^*\) priręđeno jest množstvo \(\mathfrak H^*\) vsih nenulovyh elementov od \(\mathfrak K\), ktore sut po elementah komutativne s \(\mathfrak S\). Zajedno s 0 ona očevidno tvoręt tělo, t.~j. \(H\) jest zatvorjena.

\item Treti čestny teorem svodimo na vtory. Tělu \(\mathfrak S\) priręđena jest podgrupa \(S\) od \(G\), jmenno podgrupa, prinaležeča k \(\mathfrak S^*/\mathsf P^*\), od \(G\). Tak samo grupě \(H\) prinaleži podtělo \(\mathfrak H\) od \(\mathfrak K\), i jest \(H\simeq \mathfrak H^*/\mathsf P^*\). Ako \(\mathfrak S\) jest invariantno tělo od \(H\), to \(\mathfrak S\) sostoji se iz vsih s \(\mathfrak H\) po elementah komutativnyh elementov od \(\mathfrak K\); možno to jest takože kazati: \(S\) jest invariantna grupa od \(\mathfrak H\). Analogično: Ako \(H\) jest invariantna grupa od \(\mathfrak S\), to \(\mathfrak H\) jest invariantno tělo od \(S\). Zato treti čestny teorem možno formulovati takože ovako:

\item[\(3'\).] Ako \(\mathfrak H\) imaje invariantnu grupu \(S\), to \(\mathfrak H\) jest invariantno tělo od \(S\), a to jest ravno vtoromu čestnomu teoremu, ktory ješče trěba dokazati.

\item[3.] Toj dokaz teče, kromě několiko malyh razlika, kako v komutativnom slučaju:
\end{enumerate}

Konstruujemo s \(\mathfrak K\) recipročno izomorfno tělo \(K\), čiji centar identifikujemo s \(\mathsf P\). Vměsto avtomorfizmov od \(\mathfrak K\) možno razgledati recipročne izomorfizmy od \(\mathfrak S\) s \(K\).

Ponovno hočemo dokazati pomoćnu lemu, analognu teoremu o prodolženju, upotrěbjenomu v komutativnom slučaju:

Nehaj \(\mathfrak S\) i \(\mathfrak T\) budut těla medžu \(\mathsf P\) i \(\mathfrak K\),
\[
\mathsf P \subseteq \mathfrak S \subseteq \mathfrak T \subseteq \mathfrak K.
\]
Tvrdimo, že vsaky recipročny izomorfizm od \(\mathfrak S\) na podtělo od \(K\) (recipročno prědstavjenje prvogo stupenja od \(\mathfrak S\) v \(K\)) možno na prinajmenje toliko načinov prodolžiti do recipročnyh izomorfizmov od \(\mathfrak T\), koliko ukazuje stupenj \(s\) od \(\mathfrak T\) nad \(\mathfrak S\).

Těla \(\mathsf P\), \(\mathfrak S\), \(\mathfrak T\), \(\mathfrak K\) razširjajut se s \(K\)
\[
K\subseteq \mathfrak S_K\subseteq \mathfrak T_K\subseteq \mathfrak K_K.
\]
Po teoremu 1 na str.~20 vsako \(\mathfrak S_K\) jest dvustranno prosto i polno reducibilno, a prosti lěvi idealy -- medžu soboju operatorno izomorfne -- sut moduly prědstavjenja za jedinu klasu ireducibilnyh prědstavjenij od \(\mathfrak S\) v \(K\). Zato že ta klasa prědstavjenij imaje stupenj 1 -- zato že po konstrukciji od \(K\) v \(K\) obstajajut těla recipročno izomorfne s \(\mathfrak S\) -- vsi prosti lěvi idealy imajut rang 1, t.~j. sistem jest suma toliko lěvyh idealov, koliko ukazuje njegov stupenj nad \(\mathsf P\)
\[
\mathfrak S_K=K E_1+\cdots+K E_p,\qquad K E_\nu=\mathfrak S_K E_\nu.
\]

Togda jest
\[
\begin{aligned}
\mathfrak T_K&=\mathfrak T\cdot K=\mathfrak T\cdot \mathfrak S_K,\\
\mathfrak T_K&=\mathfrak T\mathfrak S_K E_1+\cdots+\mathfrak T\mathfrak S_K E_p
              =\mathfrak T_K E_1+\cdots+\mathfrak T_K E_p.
\end{aligned}
\]
Vsaky ideal \(\mathfrak T_K\cdot E_\nu\) razpada se na \(s\) prostyh idealov, zato že
\[
[\mathfrak T_K E_\nu:K]=[\mathfrak T_K E_\nu:K E_\nu]=[\mathfrak T_K E_\nu:\mathfrak S_K E_\nu]\leqq[\mathfrak T_K:\mathfrak S_K]=[\mathfrak T:\mathfrak S]=s
\]
i
\[
\sum [\mathfrak T_K E_\nu:K]=sp
\]
zato
\[
[\mathfrak T_K E_\nu:K]=s,
\qquad
\mathfrak T_K E_\nu=K e_1^{(\nu)}+\cdots+K e_s^{(\nu)}.
\]

Prědstavjenje, dano posrědstvom \(K e_i^{(\nu)}\), od \(\mathfrak T\) jest prodolženje prědstavjenja, zadanogo posrědstvom \(K E_\nu\), od \(\mathfrak S\), zato že ako \(\alpha\in \mathfrak S\), to jest
\[
\alpha E_\nu=\sigma E_\nu,\qquad
\sum \alpha e_i^{(\nu)}=\sigma\sum e_i^{(\nu)},\qquad
\alpha e_i^{(\nu)}=\sigma e_i^{(\nu)}.
\]

Nyně trěba dokazati samo, že prědstavjenja dane posrědstvom \(K e_i^{(\nu)}\) sut vse različne. V komutativnom slučaju to bylo jasno, zato že ona naležali k različnym klasam prědstavjenij. Tu ona, nasuprot, vse naležet k istoj klase prědstavjenij, zato trěba razgledati same prědstavjenja. Za to razgledajemo dane posrědstvom \(K e_i^{(\nu)}\) \(K\)-homomorfizmy cělogo \(\mathfrak T_K\); oni sut už cělkom zadane, ako znajemo prědstavjenja od \(\mathfrak T\), zato že za konstrukciju cělogo homomorfizma dostatočno jest znati, kako sut prědstavjene bazove elementi. Možno zato ograničiti se na dokaz, že posrědstvom \(K e_i^{(\nu)}\) dane homomorfizmy od \(\mathfrak T_K\) sut vse različne. Ako \(e_i^{(\nu)}\) jeden raz prědstavimo posrědstvom modula prědstavjenja \(K e_i^{(\nu)}\), a drugi raz posrědstvom \(K e_j^{(\nu)}\), togda zaradi
\[
e_i^{(\nu)2}=e_i^{(\nu)}
\]
prvy raz najdemo 1, a zaradi
\[
e_i^{(\nu)}e_j^{(\nu)}=0
\]
drugi raz 0 kako prědstavjajuči element od \(e_i^{(\nu)}\); oba prědstavjenja sut zato naistino različne. (Točno teže zaključky možno priměniti takože v komutativnom slučaju, ale tam sut nepotrěbne).

Od nyně zaključujemo točno kako na str.~10 i str.~16. Nehaj \(H\) bude invariantna grupa od \(\mathfrak S\), a \(\mathfrak T\) invariantno tělo od \(H\). Elementi od \(H\) sut vsa prodolženja identičnogo avtomorfizma od \(\mathfrak S\) na cělo \(\mathfrak K\). Zato že po upravo dokazanoj pomoćnoj lemi medžu nimi obstajajut prinajmenje \(s\) različne na \(\mathfrak T\), mora byti \(s=1\), \(\mathfrak T=\mathfrak S\), tym jest dokazano.

\section*{Glava V. Faktorne sistemy}\label{kapitel-v.-faktorensysteme}
% END BOOK_S21

% BEGIN BOOK_S22
\subsection*{§ 22. Grupa těl s danym centrom}\label{die-gruppe-der-koerper-mit-gegebenem-zentrum}

Za osnovu našego razgledanja izberemo ustanovjeno komutativno polje \(\mathsf P\). Množstvo vsih matricnyh prstenov \(\mathfrak K_r\), čija těla avtomorfizmov \(\mathfrak K\) sut těla konečnogo ranga nad \(\mathsf P\) s samym \(\mathsf P\) kako centrom, jest multiplikativno zatvorjeny sistem po odnošenju k ,,obrazovanju prěmogo produkta``. Prěmy produkt dvoh takyh prstenov \(\mathfrak K_r\) i \(\mathfrak L_s\) jest prosto prsten \(\mathfrak K_r\cdot\mathfrak L_s\), ktoro nastava iz \(\mathfrak K_r\) razširjenjem osnovnogo polja do \(\mathfrak L_s\):
\[
\mathfrak K_r\times\mathfrak L_s=\mathfrak K_r\cdot\mathfrak L_s.
\]
Po pripomněnke na str.~22 to obrazovanje prěmogo produkta jest komutativno:
\[
\mathfrak K_r\times\mathfrak L_s=\mathfrak L_s\times\mathfrak K_r.
\]
Tvrdženje, že sistem vsih \(\mathfrak K_r\) jest multiplikativno zatvorjeny, ne znači ničto drugo než, že \(\mathfrak K_r\times\mathfrak L_s=\mathfrak K_r\mathfrak L_s\) jest vsečasno ponovno matricno prsten, čije tělo avtomorfizmov \(\mathfrak M\) imaje centar \(\mathsf P\) -- konečny rang po odnošenju k \(\mathsf P\) razuměje se sam soboju. Po teoremu 1 na str.~20 \(\mathfrak K_r\times\mathfrak L_s\) jest dvustranno prosto i polno reducibilno, zato matricno prsten nad nekym razširjenym tělom \(\mathfrak M\) od \(\mathsf P\). Fakt, že centar od \(\mathfrak M\) jest ravny \(\mathsf P\), neposrědnje slěduje iz toga, že prvo \(\mathsf P\) sigurno jest sodrživano v centru, a drugo ono mora byti sodrživano v centru \(\mathsf P\) od \(\mathfrak K\) (ili od \(\mathfrak L\)). Množenje \(\mathfrak K_r\times\mathfrak L_s\) očevidno jest asociativno. \(\mathfrak K_r\) tvoręt multiplikativno zatvorjeny sistem, ale ne grupu, zato že ako \(\mathfrak K_r\times\mathfrak L_s=\mathfrak M_m\), to sigurno \(m\geq rs\); zato ne bude možno najdti \(\mathfrak L_s\), ktoro izpolnja jednačenje \(\mathfrak K_r\times\mathfrak L_s=\mathfrak M_m\), ako \(r>m\).

Nyně sdružujemo vse \(\mathfrak K_r\) s jednakym \(\mathfrak K\) v klasu \(\{\mathfrak K\}\). Togda imaje město: Ako \(\mathfrak K_r\) i \(\mathfrak K_{r'}\) naležet k istoj klase \(\{\mathfrak K\}\), a \(\mathfrak L_s\) i \(\mathfrak L_{s'}\) k istoj klase \(\{\mathfrak L\}\), to takože \(\mathfrak K_r\times\mathfrak L_s\) i \(\mathfrak K_{r'}\times\mathfrak L_{s'}\) naležet k istoj klase. Dostatočno jest pokazati, že \(\mathfrak K_r\times\mathfrak L_s\) naleži k istoj klase kako \(\mathfrak K\times\mathfrak L\), to jest že ako \(\mathfrak K\times\mathfrak L=\mathfrak M_m\), to takože \(\mathfrak K_r\times\mathfrak L_s\) jest matricno prsten nad \(\mathfrak M\). To jest legko viděti: Jest
\[
\mathfrak K_r\times\mathfrak L
=\left(\sum_1^r\mathfrak K c_{ik}\right)\mathfrak L
=\sum_1^r\mathfrak K_{\mathfrak L}c_{ik}
=\sum_1^r\mathfrak M_m c_{ik}.
\]
Matricno prsten \(r\)-togo stupenja nad \(\mathfrak M_m\), a to očevidno jest matricno prsten \(mr\)-togo stupenja nad \(\mathfrak M\). Togda prodolžujemo točno tako
\[
\mathfrak K_r\times\mathfrak L_s
=\left(\sum_1^s\mathfrak L d_{ik}\right)_{\mathfrak K_r}
=\sum_1^s\mathfrak L_{\mathfrak K_r}d_{ik}
=\sum_1^s\mathfrak M_{mr}d_{ik}
=\mathfrak M_{mrs}
\]
čim naše tvrdženje už jest dokazano. Zapomnimo, že iz \(\mathfrak L\times\mathfrak K=\mathfrak M_m\) slěduje \(\mathfrak K_r\times\mathfrak L_s=\mathfrak M_{mrs}\).

Iz dokazanogo slěduje: ako produkt dvoh klas \(\{\mathfrak K\}\) i \(\{\mathfrak L\}\) definujemo kako klasu produkta jednogo reprezentanta od \(\{\mathfrak K\}\) s jednim reprezentantom od \(\{\mathfrak L\}\), ta definicija ne zavisi od izbora reprezentantov.

Priręđenje \(\mathfrak K_r\mapsto\{\mathfrak K\}\) jest zato homomorfizm sistema vsih \(\mathfrak K_r\) na množstvo vsih klas \(\{\mathfrak K\}\); množstvo \(\mathscr K\) klas \(\{\mathfrak K\}\) jest zato multiplikativno zatvorjeny sistem. Dokazujemo, že \(\mathscr K\) jest grupa. Kako se neposrědnje vidi, element jedinice jest klasa \(\{\mathsf P\}\), \(\mathfrak K_r\times\mathsf P=\mathfrak K_r\).

Ostaje samo dokazati za vsaku klasu \(\{\mathfrak K\}\) obstajanje klasy \(\{\mathfrak K\}^{-1}\). Pokazuje se, že trěba položiti \(\{\mathfrak K\}^{-1}=\{\overline{\mathfrak K}\}\), kde \(\{\overline{\mathfrak K}\}\) označaje tělo recipročno izomorfno s \(\mathfrak K\). Drugymi slovami: \(\mathfrak K\times\overline{\mathfrak K}\) mora byti matricno prsten nad \(\mathsf P\), ili ako \(\mathfrak K\) imaje rang \(t^2\), zato \(\mathfrak K\times\overline{\mathfrak K}\) rang \(t^4\) nad \(\mathsf P\), to \(\mathfrak K\times\overline{\mathfrak K}\) mora razpadati se na \(t^2\) prostyh lěvyh idealov. To slěduje iz toga, že prosty lěvy ideal od \(\mathfrak K\times\overline{\mathfrak K}\), kako modul prědstavjenja nekogo ireducibilnogo recipročnogo prědstavjenja od \(\overline{\mathfrak K}\) v \(\mathfrak K\) -- ktoro jest prvogo stupenja -- imaje rang 1 nad \(\mathfrak K\), to jest rang \(t^2\) nad \(\mathsf P\).

\(\mathfrak K\) hočemo nazyvati \emph{prostym tělom}, ako medžu \(\mathsf P\) i \(\mathfrak K\) ne imaje těla, čiji centar takože jest \(\mathsf P\).

\textbf{Teorem.} \emph{Vsako tělo \(\mathfrak K\) jest prěmy produkt prostyh podtěl \(\mathfrak K^{(i)}\),}
\[
\mathfrak K=\mathfrak K^{(1)}\times\cdots\times\mathfrak K^{(p)}.
\]

\emph{Dokaz.} Dostatočno jest pokazati, že ako \(\mathfrak K^{(1)}\) jest podtělo od \(\mathfrak K\) s centrom \(\mathsf P\), to obstaja podtělo \(\mathfrak S\) od \(\mathfrak K\) s centrom \(\mathsf P\), ktoro izpolnja jednačenje
\[
\mathfrak K^{(1)}\times\mathfrak S=\mathfrak K
\]
V vsakom slučaju obstaja klasa \(\{\mathfrak S\}\) s svojstvom
\[
\{\mathfrak K^{(1)}\}\times\{\mathfrak S\}=\{\mathfrak K\}
\]
jmenno
\[
\{\mathfrak S\}=\{\overline{\mathfrak K^{(1)}}\}\times\{\mathfrak K\}.
\qquad
\text{Nehaj bude}\qquad
\overline{\mathfrak K^{(1)}}\times\mathfrak K=\mathfrak S_\lambda.
\]
Togda jest
\[
\mathfrak K^{(1)}\times\mathfrak S_\lambda
=\mathfrak K^{(1)}\times\overline{\mathfrak K^{(1)}}\times\mathfrak K
=\mathsf P_{t^2}\times\mathfrak K
=\mathfrak K_{t^2},
\qquad
t^2=\text{rang od }\mathfrak K^{(1)}.
\]
Ale \(\mathfrak K^{(1)}\times\mathfrak S_\lambda\) jest matricno prsten stupenja prinajmenje ravnogo \(\lambda\), zato \(\lambda\leq t^2\). \(\mathfrak S_\lambda=\overline{\mathfrak K^{(1)}}\times\mathfrak K\) mora razpadati se na proste idealy ranga 1 nad \(\mathfrak K\), zato že \(\overline{\mathfrak K^{(1)}}\) dopušča recipročna prědstavjenja prvogo stupenja v \(\mathfrak K\); zato \(\lambda\geq t^2\), \(\lambda=t^2\),
\[
\mathfrak K_{t^2}=\mathfrak K^{(1)}\times\mathfrak S_{t^2}.
\]
zato samo \(\mathfrak K^{(1)}\times\mathfrak S\) mora byti tělo, i
\[
\mathfrak K^{(1)}\times\mathfrak S_{t^2}
=\bigl(\mathfrak K^{(1)}\times\mathfrak S\bigr)_{t^2}
=\mathfrak K_{t^2}.
\]
\(\mathfrak K^{(1)}\times\mathfrak S\) jest tělo avtomorfizmov prostyh lěvyh idealov v jednom razkladu od \(\mathfrak K_{t^2}\), a \(\mathfrak K\) tělo avtomorfizmov prostyh lěvyh idealov v drugom razkladu od \(\mathfrak K_{t^2}\). Ale zato že vsaka dva različna lěva ideala od \(\mathfrak K_{t^2}\) možno vměstiti v jeden razklad, \(\mathfrak K\) i \(\mathfrak K^{(1)}\times\mathfrak S\) sut izomorfne, a toj izomorfizm možno izbrati tako, že \(\mathsf P\) ostavaje po elementah nepodvižno. Zato po teoremu 7 na str.~25 toj izomorfizm jest stvorjen vnutrnym avtomorfizmom od \(\mathfrak K_{t^2}\).
\[
\mathfrak K
=\lambda^{-1}\bigl(\mathfrak K^{(1)}\times\mathfrak S\bigr)\lambda
=\lambda^{-1}\mathfrak K^{(1)}\lambda\times\lambda^{-1}\mathfrak S\lambda,
\qquad
\lambda\in\mathfrak K_{t^2}.
\]
\(\lambda^{-1}\mathfrak K^{(1)}\lambda\) jest s \(\mathfrak K^{(1)}\) izomorfno podtělo od \(\mathfrak K\); zato po teoremu 7 na str.~25 s odgovarjajučim \(\mu\) iz \(\mathfrak K\):
\[
\mu^{-1}\lambda^{-1}\mathfrak K^{(1)}\lambda\mu=\mathfrak K^{(1)}
\]
zato
\[
\mu^{-1}\mathfrak K\mu
=\mathfrak K
=\mathfrak K^{(1)}\times\mu^{-1}\lambda^{-1}\mathfrak S\lambda\mu,
\]
čim jest tvrdženje dokazano.
% END BOOK_S22

% BEGIN BOOK_S23
\subsection*{§ 23. Faktorne sistemy}\label{faktorensysteme}

V slědujućih razgledanjah prědpokladamo, že osnovno polje \(\mathsf P\) jest doskonalo; vměsto toga možno učiniti obćejše prědpoloženje, že razgledajut se samo taka nekomutativna \(\mathfrak K\) s centrom \(\mathsf P\), ktore imajut svojstvo, že vsako \(\mathfrak K_r\) sodrživaje maksimalno komutativno podpolje prvogo roda nad \(\mathsf P\). (Kako medžutym pokazal \foreign{Köthe}, to imaje město za vsako \(\mathfrak K\).)

Pod \(Z\) razumějemo razpadno polje nekomutativnogo těla \(\mathfrak K\) s centrom \(\mathsf P\); nehaj ireducibilno prědstavjenje \(\mathfrak Z\) od \(Z\) v \(\mathfrak K\) bude stupenja \(r\), togda ono jest maksimalno komutativno podpolje od \(\mathfrak K_r\). Nehaj \(\Gamma\) bude Galoisovo polje od \(Z\), zato
\[
\mathfrak Z_\Gamma=e_1\Gamma+\cdots+e_n\Gamma
\qquad\text{(str.~9)}
\]
\(n\) jest stupenj od \(Z\) nad \(\mathsf P\). Za s \(\Gamma\) razširjeno prsten \(\mathfrak K_r\) dobyva se takože razklad
\[
\mathfrak K_{r\Gamma}=\mathfrak K_{r\Gamma}e_1+\cdots+\mathfrak K_{r\Gamma}e_n
\]
\[
\mathfrak K_{r\Gamma}=\mathfrak l_1+\cdots+\mathfrak l_n,\qquad
\mathfrak l_i=\mathfrak K_{r\Gamma}e_i.
\]
\(\mathfrak l_i\) očevidno sut lěvi idealy. Oni sut čak prosti lěvi idealy, zato že \(Z\) jest razpadno polje, \(\mathfrak K_{r\Gamma}\) jest matricno prsten stupenja \(r\cdot t=n\) nad \(\Gamma\). (Teorem 4, str.~23, \(n=rt\) na str.~22.) Tako dostany razklad od \(\mathfrak K_{r\Gamma}\) na proste lěve idealy osoblivo jest prilagodženy algebraičnym svojstvam od \(\mathfrak K\) i \(Z\). Ako, kako na str.~18, dějstvo substitucije \(S\) iz Galoisovoj grupy \(\mathfrak G\) od \(\Gamma/\mathsf P\) na vse elementy od \(\mathfrak K_{r\Gamma}\) definujemo tako, že za vsaky \(\mathfrak K_r\)-element \(z\) položimo
\[
S(z)=z
\]
to idealy \(\mathfrak l_i\) sut konjugovane, t.~j. priměnjenje jednogo \(S\) na neko \(\mathfrak l_j\) davaje drugo
\[
S(\mathfrak l_j)=\mathfrak l_i,
\]
zato že vsaka nerazloživa komponenta jedinice \(e_j\) mora prěhoditi v neku drugu \(e_i\). (Vidi takože str.~11.)

Modul prědstavjenja \(e_i\Gamma\) obrazuje \(\mathfrak Z\) na podpolje \(Z_i\) od \(\Gamma\). \(e_i\) už leži v \(\mathfrak Z_{Z_i}\), zato že prědstavjenje \(Z_i\) od \(\mathfrak Z\) už jest sodrživano v \(Z_i\), zato \(\mathfrak Z_{Z_i}\) mora odcěpiti modul prědstavjenja \(e_iZ_i\). Nehaj \(\{Z_i,Z_k\}\) bude kompozitum od \(Z_i\) i \(Z_k\). \(\mathfrak K_{r\{Z_i,Z_k\}}\) odcěpja dva lěva ideala
\[
\bar{\mathfrak l}_i=\mathfrak K_{r\{Z_i,Z_k\}}e_i,\qquad
\bar{\mathfrak l}_k=\mathfrak K_{r\{Z_i,Z_k\}}e_k
\]
ktore sut proste, zato že imaje město \(\mathfrak l_i=\mathfrak K_{r\Gamma}\bar{\mathfrak l}_i\). zato jest
\[
\bar{\mathfrak l}_i=\sum_{\lambda=1}^m\{Z_i,Z_k\}c_{\lambda i}^{(ik)},\qquad
\bar{\mathfrak l}_k=\sum_{\lambda}\{Z_i,Z_k\}c_{\lambda k}^{(ik)}.
\]
\(c_{\lambda\rho}^{(ik)}\) jest sistem matricnyh jedinic od \(\mathfrak K_{r\{Z_i,Z_k\}}\).

Vsaky izomorfizm \(a_i\mapsto a_k\) od \(\bar{\mathfrak l}_i\) na \(\bar{\mathfrak l}_k\) s \(\mathfrak K_{r\{Z_i,Z_k\}}\) kako oblastju operatorov oprěděljen jest elementom \(p_{ik}\) s \(e_i\mapsto p_{ik}\), zato že togda jest \(a_i=a_ie_i\mapsto a_ip_{ik}\), zato \(\bar{\mathfrak l}_i\cdot p_{ik}=\bar{\mathfrak l}_k\). Vsi možne taki \(p_{ik}\) očevidno sut
\[
p_{ik}=c_{ik}^{(ik)}\gamma_{ik},\qquad \gamma_{ik}\ne0\quad\text{iz}\quad \{Z_i,Z_k\}.
\]
Posrědstvom \(a_i\mapsto a_ip_{ik}\) togda jest zadany takože operatorny izomorfizm od \(\mathfrak l_i\) na \(\mathfrak l_k\). Nyně izbirajemo ustanovjene sistemy takyh \(p_{ik}\), nalagajuči jim uslovje konjugovanosti.

Dějstvo jednogo \(S\) na indeks \(i\) nehaj bude definovano tako
\[
S(i)=k,\qquad\text{ako}\qquad S(Z_i)=Z_k.
\]
Za vsaky par indeksov \((\nu,\mu)\) obstaja prinajmenje jeden konjugovany par \(S(\nu,\mu)\) formy \(S(\nu,\mu)=(1,k)\). Iz vsakoj klasy konjugovanyh parov indeksov izbiraje se ustanovjeny par \((1,k)\), a togda za toj \((1,k)\):
\[
p_{1k}=c_{1k}^{(1k)}\gamma_{1k},\qquad \gamma_{1k}\ne0\quad\text{iz}\quad \{Z_1,Z_k\},
\]
inače se položi libovoljno, a togda, ako \(S(1,k)=(\nu,\mu)\),
\[
p_{\nu\mu}=S(p_{1k}).
\]
\(p_{\nu\mu}\) togda naistino imaje formu
\[
p_{\nu\mu}=c_{\nu\mu}^{(\nu\mu)}\gamma_{\nu\mu},\qquad
\gamma_{\nu\mu}\ne0\quad\text{iz}\quad \{Z_\nu,Z_\mu\}.
\]
Zato že dějstvo od \(S\) na izomorfizm
\[
a_1\mapsto a_1p_{1k},\qquad\text{specialno}\qquad e_1\mapsto p_{1k}
\quad\text{od }\bar{\mathfrak l}_1\text{ na }\bar{\mathfrak l}_k
\]
davaje izomorfizm
\[
a_\nu\mapsto a_\nu p_{\nu\mu},\qquad\text{specialno}\qquad e_\nu\mapsto p_{\nu\mu}
\quad\text{od }\bar{\mathfrak l}_\nu\text{ na }\bar{\mathfrak l}_\mu.
\]

Ako pod \(c_{\lambda\rho}\) razumějemo sistem matricnyh jedinic od \(\mathfrak K_{r\Gamma}\), jasno jest, že \(p_{\nu\mu}\) možno takože zapisati v formě
\[
p_{\nu\mu}=c_{\nu\mu}\bar\gamma_{\nu\mu}
\]
relacije
\[
c_{ij}c_{1k}=0,\quad\text{ako}\quad j\ne1,
\qquad
c_{ij}c_{jk}=c_{ik}
\]
imajut zato kako poslědicu relacije
\[
\begin{aligned}
p_{ij}p_{1k}&=0,\quad &&\text{ako }j\ne1,\\
p_{ij}p_{jk}&=\alpha_{ik}^{(j)}p_{ik},\quad
&&\alpha_{ik}^{(j)}\text{ iz }\Gamma\text{ (točnėje iz }\{Z_i,Z_j,Z_k\}\text{)}
\end{aligned}
\]
Te \(n^3\) veličiny \(\alpha\) tvoręt \emph{faktorny sistem od \(\mathfrak K\) pri izbranju od \(Z\)}, ktory prinaleži k sistemu \(p_{ik}\) ,,\emph{pseudomatricnyh jedinic}``.

\(\alpha\) sut konjugovane: Iz \(S(i,k,j)=(\nu,\mu,\tau)\) slěduje
\[
S\bigl(\alpha_{ik}^{(j)}\bigr)=\alpha_{\nu\mu}^{\tau}.
\]

Vsaky sistem \(p_{ik}^*\) pseudomatricnyh jedinic nastava iz nekogo sistema \(p_{ik}\) tako, že početne \(p_{ik}\) množimo libovoljnymi \(\delta_{ik}\ne0\) iz \(\{Z_i,Z_k\}\):
\[
p_{1k}^*=p_{1k}\delta_{1k}
\]
i togda ponovno, ako \(S(1,k)=(\nu,\mu)\), položimo
\[
p_{\nu\mu}^*=S(p_{1k}^*)
\]
zato jest
\[
p_{\nu\mu}^*=p_{\nu\mu}\delta_{\nu\mu}
\]
s \(\delta_{\nu\mu}\) iz \(\{Z_\nu,Z_\mu\}\); \(\delta_{\nu\mu}\) sut konjugovane, iz \(S(i,k)=(\nu,\mu)\) slěduje \(S(\delta_{ik})=\delta_{\nu\mu}\). Ako \(\alpha^*\) jest faktorny sistem prinaležeči k \(p_{ik}^*\), to jest
\[
\alpha_{ik}^{(j)*}=\frac{\delta_{ij}\delta_{jk}}{\delta_{ik}}\alpha_{ik}^{(j)}.
\]

Faktorne sistemy \(\alpha^*\) i \(\alpha\) nazyvajut se \emph{asociovanymi} (tak samo \(p_{ik}^*\) i \(p_{ik}\)); vsi faktorne sistemy od \(\mathfrak K\) pri izbranju od \(Z\) tvoręt klasu \(\{\alpha\}\) asociovanyh faktornyh sistemov.

Klasa \(\{\alpha\}\) ne zavisi od toga, ktoro prědstavjenje \(\mathfrak Z\) od \(Z\) v \(\mathfrak K_r\) izberemo za osnovu jej definicije. Zato že prěhod od jednogo \(\mathfrak Z\) k drugomu možno učiniti transformacijoju elementom \(\lambda\) iz \(\mathfrak K_r\), ktory jest zato \(\mathfrak G\)-invariantny; zato iz sistema \(p_{ik}\), prinaležečego k \(\mathfrak Z\), nastava sistem
\[
p'_{ik}=\lambda^{-1}p_{ik}\lambda,
\]
ktory prinaleži k \(\mathfrak Z'\); faktorne sistemy od \(p_{ik}\) i \(p'_{ik}\) ale sut identične.

Faktorny sistem ne jest ničto drugo než tablica množenja od \(\mathfrak K_{n\Gamma}\), prilagodžena osoblivym algebraičnym svojstvam od \(\mathfrak K\).
% END BOOK_S23

% BEGIN BOOK_S24
\subsection*{§ 24. Množenje faktornyh sistemov}\label{multiplikation-von-faktorensystemen}

Vměsto o razpadnom polju nekomutativnogo těla \(\mathfrak K\) govorimo takože o razpadnom polju od \(\{\mathfrak K\}\); to imaje smysl, zato že spolu s \(\mathfrak K_Z\) takože vsako prsten \(\mathfrak K_{rZ}\) razpada se na absolutno ireducibilne komponenty, i obratno.

Vsaka dva těla \(\mathfrak K\), \(\mathfrak S\) s centrom \(\mathsf P\) imajut obća razpadna polja, napr. kompozitum razpadnogo polja od \(\mathfrak K\) i razpadnogo polja od \(\mathfrak S\).

\textbf{Teorem 1.} \emph{Klasy \(\{\mathfrak K\}\), ktore imajut polje \(Z\) kako razpadno polje, tvoręt podgrupu \(\mathscr K_Z\) grupy \(\mathscr K\) vsih \(\{\mathfrak K\}\).}

\emph{Dokaz.} Ako \(\mathfrak K\) i \(\mathfrak S\) imajut \(Z\) kako razpadno polje, to \(\mathfrak K_Z\) i \(\mathfrak S_Z\), zato takože \((\mathfrak K\times\mathfrak S)_Z\), razpadajut se na absolutno ireducibilne komponenty, t.~j. \(Z\) jest razpadno polje od \(\{\mathfrak K\times\mathfrak S\}\) (teorem 4 na str.~23).

Recipročno izomorfne těla imajut vsa razpadna polja obće.

Pomoćno razgledanje.
\[
\mathfrak K_r=\mathfrak L_1+\cdots+\mathfrak L_r,\qquad
\mathsf P_m=\mathfrak z_1+\cdots+\mathfrak z_m
\]
nehaj budut razklady od \(\mathfrak K_r\), \(\mathsf P_m\) na proste lěve idealy; togda
\[
\mathfrak K_{rm}=\mathfrak K_r\times\mathsf P_m
=\mathfrak L_1\mathfrak z_1+\cdots+\mathfrak L_i\mathfrak z_j+\cdots+\mathfrak L_r\mathfrak z_m
\]
jest razklad od \(\mathfrak K_{rm}\) na \(rm\) nenulovyh, zato prostyh lěvyh idealov. Zato vsaky \(r\)-členny lěvy ideal \(A\) od \(\mathfrak K_{rm}\) jest svęzany s \(r\)-člennym lěvym idealom
\[
\mathfrak L_1\mathfrak z_1+\cdots+\mathfrak L_r\mathfrak z_1=\mathfrak K_r\mathfrak z_1
\]
posrědstvom vnutrnogo avtomorfizma od \(\mathfrak K_r\)
\[
A=\tau^{-1}\mathfrak K_r\tau\cdot\tau^{-1}\mathfrak z_1\tau
=\mathfrak K'_r\cdot\mathfrak z'_1.
\]

Ako \(\mathfrak K_{r\Gamma}=\mathfrak l_1+\cdots+\mathfrak l_n\) jest razklad od \(\mathfrak K_{r\Gamma}\) na konjugovane lěve idealy, prinaležeči k nekomu \(\mathfrak Z\) v \(\mathfrak K_r\), \(\mathfrak l_i=\mathfrak K_{r\Gamma}e_i\), to zaradi
\[
S(\tau^{-1}\mathfrak l_i\tau\cdot\mathfrak z'_1)
=S(\tau^{-1}\mathfrak l_i\tau)\,\mathfrak z'_1
=\tau^{-1}(S\mathfrak l_i)\tau\cdot\mathfrak z'_1
=\tau^{-1}\mathfrak l_k\tau\cdot\mathfrak z'_1
\]
v
\[
\tag{1}
A_\Gamma=\mathfrak l'_1\mathfrak z'_1+\cdots+\mathfrak l'_n\mathfrak z'_1
\]
imajemo razklad od \(A_\Gamma\) na konjugovane lěve idealy. Jest
\[
\mathfrak l'_i\mathfrak z'_1=\mathfrak K_{rm\Gamma}\cdot e'_iE',
\]
ako položimo \(\tau^{-1}e_i\tau=e'_i\), a \(E'\) označaje jedinicu od \(\mathfrak z'_1\).

Vsaky drugy razklad
\[
\tag{2}
A_\Gamma=\mathfrak q_1+\cdots+\mathfrak q_n,\qquad
S\mathfrak q_i=\mathfrak q_k\quad\text{ako}\quad
S\mathfrak l'_i\mathfrak z'_1=\mathfrak l'_k\mathfrak z'_1
\]
na konjugovane lěve idealy nastava iz (1) transformacijoju s \(\mathfrak K_{rm}\)-elementom \(\lambda\)
\[
\mathfrak q_i=\lambda^{-1}\mathfrak l'_i\mathfrak z'_1\lambda,
\qquad
\text{to jest, ako \(e_i^*\) jest jedinica od \(\mathfrak q_i\),}\quad
e_i^*=\lambda^{-1}e'_iE'\lambda.
\]

Dokaz jest neposrědnja poslědica slědujućej pomoćnoj lemy:

\textbf{Pomoćna lema.} \emph{Nehaj \(A\) bude lěvy ideal iz \(\mathfrak K_n\), i nehaj}
\[
A=\mathfrak l_1+\cdots+\mathfrak l_n
=\mathfrak K_{n\Gamma}e_1+\cdots+\mathfrak K_{n\Gamma}e_n
\]
\emph{bude -- po prědpoloženju obstajajuči -- razklad na konjugovane lěve idealy; pri tom možno -- bez ograničenja obćosti -- prědpokladati, že \(\mathfrak l_i\) sut konjugovane.}

\emph{Nehaj \(\mathfrak z\) bude operatorno izomorfny s \(A\); i nehaj budut izpolnjene teže prědpoloženja:}
\[
\mathfrak z=\mathfrak q_1+\cdots+\mathfrak q_n
=\mathfrak K_{n\Gamma}e_1^*+\cdots+\mathfrak K_{n\Gamma}e_n^*
\]
\emph{(zato \(\mathfrak q_i\) i \(e_i^*\) sut konjugovane). Togda obstaja regularny element \(\lambda\) iz \(\mathfrak K_n\), tako že \(\mathfrak q_i=\lambda^{-1}\mathfrak l_i\lambda\).}

\emph{Dokaz. 1. Fakt, že \(e_i\) možno prědpokladati konjugovanymi, slěduje tako:} Ako \(e_1\) jest libovoljny tvoreči idempotent od \(\mathfrak l_1\), to \(\mathfrak l_1=\mathfrak K_{n\Gamma}e_1\), priměnjenjem avtomorfizma na \(\Gamma\) slěduje \(\mathfrak l_i=\mathfrak K_{n\Gamma}e_i\), kde \(e_i\) jest konjugovany s \(e_1\) i ponovno tvoreči idempotent.

2. \emph{Vsako \(\mathfrak l\) jest izomorfno s vsakym \(\mathfrak q\);} zato že \(\mathfrak l\), kako konjugovane, vse imajut jednaku (absolutnu) dolžinu, čisla od \(\mathfrak l\) i \(\mathfrak q\) sut jednaka, a \(\mathfrak z\) i \(A\) sut po prědpoloženju izomorfne.

3. Priręđenjem
\[
\tag{3}
e_1\mapsto s_1=e_1s_1;\qquad
e_1^*\mapsto t_1=e_1^*t_1;\qquad
s_1t_1=e_1;\qquad
t_1s_1=e_1^*
\]
nehaj budut ustanovjene izomorfizm medžu \(\mathfrak l_1\) i \(\mathfrak q_1\) i njegov obrat. Priměnjenjem vsih avtomorfizmov (3) prěhodi v
\[
e_i\mapsto s_i=e_is_i;\qquad
e_i^*\mapsto t_i=e_i^*t_i;\qquad
s_it_i=e_i;\qquad
t_is_i=e_i^*,
\]
čim sut zato ustanovjene izomorfizmy \(\mathfrak q_i=\mathfrak l_is_i\) i \(\mathfrak l_i=\mathfrak q_it_i\).

4. Nyně položimo:
\[
s=s_1+\cdots+s_n;\qquad t=t_1+\cdots+t_n;
\]
togda dobyva se:
\[
\mathfrak q_i=\mathfrak l_is;\qquad
\mathfrak l_i=\mathfrak q_it;\qquad
st=e_1+\cdots+e_n=\xi;\qquad
ts=e_1^*+\cdots+e_n^*=\xi^*,
\]
pričem \(\xi\) i \(\xi^*\) ležet v \(\mathfrak K_n\); dobyva se
\[
A\xi=A,\qquad \mathfrak z\xi^*=\mathfrak z\quad
(\text{zaradi }\mathfrak l_i\xi=\mathfrak l_i\mathfrak l_i=\mathfrak l_i),
\]
\[
\xi^2=\xi,\qquad \xi^{*2}=\xi^*.
\]
Ako nyně položimo: \(e=\overline E+E=\overline E^*+E^*\) i s tym:
\[
\mathfrak K_n=\overline A+A
=\mathfrak K_n\overline\xi+\mathfrak K_n\xi
=\overline{\mathfrak z}+\mathfrak z
=\mathfrak K_n\overline\xi^*+\mathfrak K_n\xi^*,
\]
to \(\overline A\) i \(\overline{\mathfrak z}\) sut takože izomorfne; ako zato
\[
\overline E\mapsto\overline s=\overline E\,\overline s
\qquad\text{i}\qquad
\overline E^*\mapsto\overline t=\overline E^*\overline t
\]
sut obratne izomorfizmy, i ako položimo:
\[
\lambda=s+\overline s,\qquad \mu=t+\overline t,
\]
dobyva se:
\[
\lambda\mu=\mu\lambda=e,\qquad
\mathfrak l_i\lambda=\mathfrak q_i,\qquad
\lambda^{-1}\mathfrak q_i=\mathfrak q_i
\quad
(\lambda^{-1}\mathfrak q_i\subseteq\mathfrak q_i
=\lambda\cdot\lambda^{-1}\mathfrak q_i\subseteq\lambda^{-1}\mathfrak q_i)
\]
i zato
\[
\lambda^{-1}\mathfrak l_i\lambda=\mathfrak q_i,\qquad \text{tym jest dokazano.}
\]

K razkladu (1) možno ponovno oprěděliti elementy
\[
r_{ik}=c_{ik}^{\prime(ik)}\gamma_{ik},\qquad \gamma_{ik}\in\{Z_i,Z_k\},
\]
ktore posrědstvom \(e'_iE'\mapsto r_{ik}\) davajut operatorne izomorfizmy medžu idealami \(\mathfrak l'_i\mathfrak z'_1\) i \(\mathfrak l'_k\mathfrak z'_1\), i takože izpolnjajut uslovje konjugovanosti \(S(r_{ik})=r_{\nu\mu}\), ako \(S(i,k)=(\nu,\mu)\). Oni posrědstvom
\[
r_{ij}r_{jk}=\alpha_{ik}^{(j)}r_{ik}
\]
oprěděljajut faktorny sistem \(\alpha\) od \(A_\Gamma\). Nyně jest važno, že za oprěděljenje faktornyh sistemov od \(A_\Gamma\) možno takože upotrěbiti libovoljny razklad (2). Zato že ako sistem \(r_{ik}\), prinaležeči k (1), zaměnimo s \(r'_{ik}=\lambda^{-1}r_{ik}\lambda\), to \(e_i^*\mapsto r'_{ik}\) davaje operatorny izomorfizm od \(\mathfrak q_i\) na \(\mathfrak q_k\); \(r'_{ik}\) izpolnjajut uslovje konjugovanosti, zato že \(\lambda\) jest \(\mathfrak q\)-invariantny, a \(r'_{ik}\) imajut tojže faktorny sistem kako \(r_{ik}\).

Ako izberemo sistem pseudomatricnyh jedinic \(p_{ik}\) od \(\mathfrak K\), to
\[
r_{ik}=\tau^{-1}p_{ik}\tau\cdot E'
\]
jest \(r_{ik}\)-sistem od \(A_\Gamma\) s jednakym faktornym sistemom. zato imajemo pomoćnu lemu:

\textbf{Teorem.} \emph{Faktorne sistemy od \(A_\Gamma\), obrazovane s pomoćju libovoljnogo razklada}
\[
A_\Gamma=\mathfrak q_1+\cdots+\mathfrak q_n
\]
\emph{od \(A_\Gamma\) na konjugovane lěve idealy, sut jednaka faktornym sistemam od \(\mathfrak K_r\).}

Slědujući teorem pokazuje, že ako \(\alpha_{ik}^{(j)}\), \(\beta_{ik}^{(j)}\) sut faktorne sistemy od \(\{\mathfrak K\}\), \(\{\mathfrak S\}\) pri izbranju od \(Z\), to \(\varepsilon_{ik}^{(j)}=\alpha_{ik}^{(j)}\cdot\beta_{ik}^{(j)}\) jest faktorny sistem od \(\{\mathfrak K\times\mathfrak S\}\) pri izbranju od \(Z\). Zato imaje smysl definirati \emph{produkt \(\{\alpha\}\times\{\beta\}\) dvoh klas asociovanyh faktornyh sistemov} \(\{\alpha\}\), \(\{\beta\}\) kako klasu \(\{\alpha_{ik}^{(j)}\beta_{ik}^{(j)}\}\).

\textbf{Teorem 2.} \emph{Nehaj \(\{\mathfrak K\}\) pri izbranju od \(Z\) imaje klasu \(\{\alpha\}\) asociovanyh faktornyh sistemov. Klasy \(\{\alpha\}\) vsih \(\{\mathfrak K\}\) s razpadnym poljem \(Z\) tvoręt grupu, ktora jest posrědstvom \(\{\mathfrak K\}\mapsto\{\alpha\}\) izomorfno odnesena na \(\mathscr K_Z\).}

\textbf{Za dokaz} trěba pokazati dvoje:

1. Ako \(\{\mathfrak K\}\) imaje faktorny sistem \(\alpha_{ik}^{(j)}\), a \(\{\mathfrak S\}\) \(\beta_{ik}^{(j)}\), to \(\varepsilon_{ik}^{(j)}=\alpha_{ik}^{(j)}\beta_{ik}^{(j)}\) jest faktorny sistem od \(\{\mathfrak K\times\mathfrak S\}\).

2. Ako \(\{\mathfrak K\}\) imaje faktorny sistem \(\alpha_{ik}^{(j)}=1\), to \(\mathfrak K=\mathsf P\).

1. Nehaj stupenji ireducibilnyh prědstavjenij od \(Z\) v \(\mathfrak K\), \(\mathfrak S\) budut \(r,s\)
\[
\mathfrak K_{r\Gamma}=\mathfrak l_1+\cdots+\mathfrak l_n,\qquad
\mathfrak S_{s\Gamma}=\mathfrak m_1+\cdots+\mathfrak m_n
\]
razklady na konjugovane idealy, pričem \(\mathfrak l_i\) i \(\mathfrak m_i\) trěba da budut vzajemno odgovarjajuči idealy (prve izvedene iz odgovarjajučih komponentov od \(\mathfrak Z_r\), a druge iz \(\mathfrak Z'_r\)). Togda jest
\[
(\mathfrak K_r\times\mathfrak S_s)_\Gamma
=\mathfrak l_1\mathfrak m_1+\cdots+\mathfrak l_i\mathfrak m_j+\cdots+\mathfrak l_n\mathfrak m_n
\]

razklad na \(n^2\) nenulovyh, zato prostyh idealov. Zato jest
\[
A_1=\mathfrak l_1\mathfrak m_1+\mathfrak l_2\mathfrak m_2+\cdots+\mathfrak l_n\mathfrak m_n
\]
ideal absolutnoj dolžiny \(n\). \(A_1\) jest invariantny pri vsih \(S\), zato že \(\mathfrak m_i\) transformujut se tako kako \(\mathfrak l_i\); zato obstaja \(\mathfrak K_r\times\mathfrak S_s\)-ideal \(A\) s
\[
A_1=A_\Gamma
\]
(pomoćna lema, str.~18).

Ako \(u\) jest stupenj ireducibilnogo prědstavjenja od \(\mathfrak Z\) v tělu avtomorfizmov \(\mathfrak M\) prostyh idealov od \(\mathfrak K\times\mathfrak S\), to \(A\) razpada se na \(u\) prostyh idealov (zato že absolutna dolžina jest \(n\)). Naša pomoćna lema davaje zato, že faktorne sistemy od \(A\) sut takože faktorne sistemy od \(\mathfrak M\).

Ale razklad od \(A_\Gamma\) na konjugovane idealy, upotrěbjeny za definiciju, jest:
\[
\tag{3}
A_\Gamma=\mathfrak l_1\mathfrak m_1+\mathfrak l_2\mathfrak m_2+\cdots+\mathfrak l_n\mathfrak m_n.
\]
Ako \(\alpha_{ik}^{(j)}\), \(\beta_{ik}^{(j)}\) prinaležet k pseudomatricnym jedinicam \(p_{ik}\), \(q_{ik}\) od \(\mathfrak K\), \(\mathfrak S\), to očevidno \(r_{ik}=p_{ik}q_{ik}\) jest sistem pseudomatricnyh jedinic od \(A_\Gamma\) po odnošenju k razkladu (3); k tomu prinaležeči faktorny sistem \(\varepsilon_{ik}^{(j)}=\alpha_{ik}^{(j)}\beta_{ik}^{(j)}\) jest togda takože jeden iz faktornyh sistemov od \(\mathfrak M\).

2. Nehaj \(\{\mathfrak K\}\) imaje faktorny sistem \(\alpha_{ik}^{(j)}=1\) s \(p_{ik}\) kako prinaležečimi pseudomatricnymi jedinicami i razkladom
\[
\mathfrak K_{r\Gamma}=\mathfrak l_1+\cdots+\mathfrak l_n
\]
na konjugovane idealy. Trěba dokazati \(\mathfrak K=\mathsf P\), ili \(\mathfrak K_r=\mathsf P_r\), ili obstajanje \(n\)-člennogo lěvogo ideala \(\mathfrak l\) od \(\mathfrak K_r\) -- zato že prosty lěvy ideal od \(\mathfrak K_r\) imaje rang \(t^2r\), zato \(n=t^2rh\) s cělym \(h\), ili mora byti \(t=1\).

Ideal od \(\mathfrak K_{r\Gamma}\)
\[
\Omega=\mathfrak l_1(p_{11}+\cdots+p_{1n})
\]
ne jest nula -- \(e_1\sum_i p_{1i}\ne0\) -- i jest prosty, zato ranga \(n\). \(\Omega\) jest \(\mathfrak G\)-invariantny,
\[
S\Omega
=\mathfrak l_i\sum_\nu S p_{1\nu}
=\mathfrak l_i\sum_\mu p_{i\mu}
=\mathfrak l_i\sum_\mu p_{i1}p_{1\mu}
\]
zato že \(\alpha_{ik}^{(j)}=1\). Zato
\[
S\Omega=\mathfrak l_ip_{i1}\sum_\mu p_{1\mu}
=\mathfrak l_1\sum_\mu p_{1\mu}
=\Omega.
\]
zato \(\Omega=\mathfrak l_\Gamma\), \emph{kde \(\mathfrak l\) označaje \(\mathfrak K_r\)-ideal \(n\)-togo ranga.}

\textbf{Teorem 3.} \emph{Vsaky element \(\{\mathfrak K\}\) grupy \(\mathscr K\) imaje konečny ręd \(\lambda\), ktory jest dělitelj absolutnogo čisla komponentov \(t\) od \(\mathfrak K\).}

\emph{Prvy dokaz (Brauer).} Izberemo dva razklady od \(\mathfrak K_{r\Gamma}\) na konjugovane lěve idealy
\[
\mathfrak K_{r\Gamma}=\mathfrak l_1+\cdots+\mathfrak l_n=\mathfrak z_1+\cdots+\mathfrak z_n,\qquad
\mathfrak l_i=\mathfrak K_{r\Gamma}e_i,\quad
\mathfrak z_i=\mathfrak K_{r\Gamma}e_i^*.
\]
Nehaj posrědstvom \(a_i\mapsto a_ir_{ik}\), specialno \(e_i\mapsto r_{ik}\), bude zadany operatorny izomorfizm od \(\mathfrak l_i\) na \(\mathfrak z_k\). Pri njegovom obratu \(e_k^*\) prěhodi v neko \(s_{ki}\) iz \(\mathfrak l_i\), \(s_{ki}\cdot r_{ik}=e_k^*\). Ako izberemo \(r_{ik}\) konjugovanymi
\[
S(r_{ik})=r_{\nu\mu}\quad\text{ako}\quad S(i,k)=(\nu,\mu),
\]
to \(s_{ik}\) takože sut konjugovane:
\[
S(s_{ki})S(r_{ik})=S(s_{ki}r_{ik})=S(e_k^*)=e_\mu^*
=s_{\mu\nu}r_{\nu\mu}=s_{\mu\nu}S(r_{ik}),
\]
zato \(S(s_{ki})=s_{\mu\nu}\). Zato že posrědstvom \(e_i\mapsto e_ir_{ij}s_{jk}\) jest zadany operatorny izomorfizm od \(\mathfrak l_i\) na \(\mathfrak l_k\), mora byti
\[
r_{ij}s_{jk}=p_{ik}\xi_{ik}^{(j)},\qquad
\xi_{ik}^{(j)}\ne0\in\Gamma
\]

Zato že \(r\)-elementi, \(s\)-elementi i \(p\)-elementi sut konjugovane, \(\xi\) takože sut konjugovane: Iz \(S(i,j,k)=(\nu,\tau,\mu)\) slěduje
\[
S(\xi_{ik}^{(j)})=\xi_{\nu\mu}^{\tau}.
\]

Iz
\[
r_{i\lambda}s_{\lambda j}r_{j\lambda}s_{\lambda k}
=r_{i\lambda}e_\lambda^*s_{\lambda k}
=r_{i\lambda}s_{\lambda k}
\]
slěduje
\[
p_{ij}\xi_{ij}^{(\lambda)}\cdot p_{jk}\xi_{jk}^{(\lambda)}
=p_{ik}\xi_{ik}^{(\lambda)}
\]
ili, zato že \(p_{ij}p_{jk}=p_{ik}\alpha_{ik}^{(j)}\):
\[
\alpha_{ik}^{(j)}
=\frac{\xi_{ik}^{(\lambda)}}{\xi_{ij}^{(\lambda)}\xi_{jk}^{(\lambda)}}.
\]
Ako položimo \(\delta_{\nu\mu}=\prod_{\lambda}\xi_{\nu\mu}^{(\lambda)}\), to \(\delta_{\nu\mu}\) jest invariantny pri vsih \(S\), ktore ostavjajut \(\nu\) i \(\mu\) nepodvižnymi, leži zato v \(\{Z_\nu,Z_\mu\}\); nadto \(\delta_{\nu\mu}\) sut konjugovane, zato že \(\xi_{\nu\mu}^{(j)}\) sut konjugovane. Zato že nyně
\[
\alpha_{ik}^{(j)n}
=\frac{\prod_{\lambda}\xi_{ik}^{(\lambda)}}
       {\prod_{\lambda}\xi_{ij}^{(\lambda)}\prod_{\lambda}\xi_{jk}^{(\lambda)}}
=\frac{\delta_{ik}}{\delta_{ij}\delta_{jk}}
\]
iz formuly za asociovane faktorne sistemy (str.~32) slěduje, že \(\alpha_{ik}^{(j)n}\) jest asociovany s sistemom \(\beta_{ik}^{(j)}=1\); zato, zato že ręd od \(\{\mathfrak K\}\) jest ravny rędu od \(\{\alpha\}\), ręd od \(\{\mathfrak K\}\) jest dělitelj od \(n\). Ako osoblivo izberemo \(\mathfrak Z\) kako maksimalno komutativno podpolje od \(\mathfrak K\), to jest \(n=t\), to \(\lambda\) jest dělitelj od \(t\).

\emph{Drugy dokaz (Schur).} Toj dokaz osnovan jest na prědstavjenju \(p_{ik}\) posrědstvom regularnyh matric ręda \(n\), \(P_{ik}\), v \(\Gamma\), ktore izpolnjajut jednake uslovja konjugovanosti kako \(p_{ik}\). \(P_{ik}\) leži v \(\{Z_i,Z_k\}\); ako \(p_{ik}\) prinaleži k \(P_{ik}\), to \(S(p_{ik})\) prinaleži k \(S(P_{ik})\). To prědstavjenje jest homomorfno; t.~j. ako \(P_{ik}\) prinaleži k \(p_{ik}\), a \(P_{kj}\) k \(p_{kj}\), to \(P_{ik}P_{kj}\) prinaleži k \(p_{ik}p_{kj}\), ili
\[
P_{ik}P_{kj}=\alpha_{ij}^{(k)}P_{ij}
\]
\((P_{ik_1}P_{k_2j}=0\) za \(k_1\neq k_2\) ne može byti pravdivo). Ako imajemo to prědstavjenje od \(p_{ik}\), zaključujemo tako: Prěhod k determinantam davaje:
\[
|P_{ik}|\cdot |P_{kj}|=\alpha_{ij}^{(k)n}\cdot |P_{ij}|,
\]
tym samym, zato že matrice \(P_{ik}\) sut regularne, \(|P_{\nu\mu}|\ne0\),
\[
\alpha_{ij}^{(k)n}=\frac{\delta_{ik}\delta_{kj}}{\delta_{ik}},
\]
kde jest položeno \(|P_{\nu\mu}|=\delta_{\nu\mu}\). I potom dalje kako v Brauerovom dokazu.

\emph{Prědstavjenje od \(p_{ik}\) posrědstvom matric \(P_{ik}\):}

Idealy \(\mathfrak l_i\) razklada
\[
\mathfrak K_{r\Gamma}=\mathfrak l_1+\cdots+\mathfrak l_n
\]
sut recipročne moduly prědstavjenja od \(\mathfrak K_r\) v \(\Gamma\) (teorem 2 na str.~32). Ako za osnovu izberemo ustanovjenu bazu \(k_1,\ldots,k_n\) od \(\mathfrak K_r\) po odnošenju k \(\mathfrak Z\), to jest
\[
\mathfrak K_{r\Gamma}=k_1\mathfrak Z_r+\cdots+k_n\mathfrak Z_r,
\]
zato
\[
\mathfrak l_i=\mathfrak K_{r\Gamma}e_i
=k_1\mathfrak Z_r e_i+\cdots+k_n\mathfrak Z_r e_i
=k_1e_i\Gamma+\cdots+k_ne_i\Gamma.
\]
Iz toga slěduje
\[
\mathfrak l_k=\mathfrak l_i p_{ik}
=k_1p_{ik}\Gamma+\cdots+k_np_{ik}\Gamma.
\]
Dvě bazy \((k_1e_i,\ldots,k_ne_i)\) i \((k_1p_{ik},\ldots,k_np_{ik})\) od \(\mathfrak l_k\) možno prěvesti jednu v drugu množenjem regularnoj matricoju \(P_{ik}\):
\[
P_{ik}\cdot
\begin{pmatrix}
k_1e_i\\ \vdots\\ k_ne_i
\end{pmatrix}
=
\begin{pmatrix}
k_1p_{ik}\\ \vdots\\ k_np_{ik}
\end{pmatrix},
\qquad
\text{kratko}\quad
P_{ik}(k_\varrho e_i)=(k_\varrho p_{ik}).
\]
Iz toga slěduje konjugovanost od \(P_{ik}\):
\[
S(P_{ik})(k_\varrho S(e_k))=(k_\varrho S(p_{ik}))
\quad\text{ili}\quad
S(P_{ik})(k_\varrho e_\mu)=(k_\varrho p_{\nu\mu})
=P_{\nu\mu}(k_\varrho e_\mu),
\]
t.~j.
\[
S(P_{ik})=P_{\nu\mu}\quad\text{ako}\quad S(i,k)=(\nu,\mu).
\]
Iz toga ponovno slěduje, že \(P_{ik}\) leži v \(\{Z_i,Z_k\}\) -- jmenno v invariantnom polju tyh \(S\), ktore ostavjajut \(i\) i \(k\) nepodvižnymi. Homomorfnost takože jest legko viděti:
\[
\begin{split}
P_{ik}P_{kj}(k_\varrho e_j)
&=P_{ik}(k_\varrho p_{kj})
=(k_\varrho p_{ik}p_{kj})\\
&=\alpha_{ij}^{(k)}(k_\varrho p_{ij})
=\alpha_{ij}P_{ij}(k_\varrho e_j),
\quad\text{zato}\\
&\cdot\, P_{ik}P_{kj}=\alpha_{ij}^{(k)}P_{ij},
\quad\text{tym jest dokazano.}
\end{split}
\]

Teorem 3 možno posiliti:

\textbf{Teorem 4.} \emph{Ako \(p\) jest prosty faktor absolutnogo čisla komponentov \(t\) od \(\mathfrak K\), to \(p\) vhodi takože v ręd \(\lambda\) od \(\{\mathfrak K\}\). (Brauer)}

\emph{Dokaz.} Nehaj \(p^\sigma\) bude najvyšša potencija od \(p\), ktora se nalazi v stupenju \(n\) razpadnogo polja \(Z\), \(\mathfrak H\) podgrupa ręda \(p^\sigma\) od \(\mathfrak G\) (ona obstaja po \emph{teoremu od Sylova}, Speiser 1. izd. § 19, str.~42, teorem 43), a \(T\) invariantno polje od \(\mathfrak H\).

Togda \([\Gamma:T]=p^\sigma\) i \([T:\mathsf P]=n/p^\sigma\) ne jest dělivy posrědstvom \(p\), a tym vyše ne posrědstvom \(t\). \(T\) zato ne jest razpadno polje; tělo avtomorfizmov \(\mathfrak S\) od \(\mathfrak K_T=\mathfrak S_m\) jest različno od \(\mathsf P\). Ale zato že \(\mathfrak S\) imaje razpadno polje \(\Gamma\), čiji stupenj nad centrom \(T\) od \(\mathfrak S\) jest ravny \(p^\sigma\), absolutno čislo komponentov od \(\mathfrak S\) jest potencija \(p^\beta>1\), zato po teoremu 3 ręd od \(\{\mathfrak S\}\) jest potencija \(p^\alpha>1\) od \(p\). Zato že iz
\[
\{\mathfrak K\}^{\lambda}=\{\mathsf P\},
\qquad
\{\mathfrak K_T\}^{\lambda}=\{T\}
\]
slěduje, to \(\lambda\equiv0\pmod {p^\alpha}\), tym jest dokazano. (Brauer dal priměry, že ne mora byti \(\lambda=t\)\srcfn{1}{\foreign{Brauer: Untersuchungen über die arithmetischen Eigenschaften von Gruppen linearer Substitutionen II. Math. Zeitschr. 31, str.~733 § 5, (1930).}}).

\textbf{Teorem 5.} \emph{Nehaj \(t\), razloženo na proste faktory, bude \(t=\prod_i p_i^{\varrho_i}\), \(\varrho_i\ge1\), \(p_i\ne p_j\), ako \(i\ne j\). Togda jest \(\mathfrak K=\mathfrak K^{(1)}\times\mathfrak K^{(2)}\times\cdots\), kde \(\mathfrak K^{(i)}\) jest tělo s absolutnym čislom komponentov \(p_i^{\varrho_i}\).}

\emph{Dokaz.} Po teoremah 3 i 4 ręd \(\lambda\) od \(\{\mathfrak K\}\) jest
\[
\lambda=\prod_i p_i^{\alpha_i},\qquad 1\le\alpha_i\le\varrho_i.
\]
Posrědstvom \(\{\mathfrak K\}\) stvorjena ciklična podgrupa \(\mathfrak Z\) od \(\mathscr K\) jest prěmy produkt podgrup \(\mathfrak Z_i\) rędov \(p_i^{\alpha_i}\), to jest
\[
\{\mathfrak K\}=\prod_i\{\mathfrak K^{(i)}\},
\]

iz čego dobyva se, že vsako razpadno polje \(Z\) od \(\mathfrak K\) jest takože razpadno polje od \(\mathfrak K^{(i)}\), zato absolutno čislo komponentov \(p_i^{\sigma_i}\) od \(\mathfrak K^{(i)}\) děli \(t\), \(\sigma_i\le\varrho_i\). Nehaj \(\prod_i \mathfrak K^{(i)}=\mathfrak K_m\), to jest
\[
\prod_i p_i^{\varrho_i}=\prod_i p_i^{\sigma_i}\cdot m,
\qquad
\sigma_i\ge\varrho_i,
\]
zato
\[
\sigma_i=\varrho_i,\qquad m=1,\qquad
\mathfrak K=\prod_i \mathfrak K^{(i)},
\quad\text{tym jest dokazano.}
\]
% END BOOK_S24

% BEGIN BOOK_S25
\subsection*{§ 25. Normalno prědstavjenje od \(\mathfrak K_r\) s Galoisovym maksimalnym komutativnym podpoljem\protect\footnotemark}
\label{normaldarstellung-von-r_r-mit-galoisschem-maximalem-kommutativem-teilkoerper-1}
\footnotetext{Prostějše osnovanje od § 27 dalje.}

Nehaj maksimalno komutativno podpolje \(\mathfrak Z\) od \(\mathfrak K_r\) bude Galoisovo. \(n\) avtomorfizmov \(H_i\) od \(\mathfrak Z/\mathsf P\) sut po teoremu 7 na str.~25 vnutrne avtomorfizmy od \(\mathfrak K_r\), t.~j. obstajajut elementi \(u_i\) s
\[
H_i z=u_i^{-1}zu_i\quad\text{za vse }z\text{ iz }\mathfrak Z.
\]
Vsaky \(u_i\) jest oprěděljen do nenulovogo faktora \(y\) iz \(\mathfrak Z\). Zato že ako \(u^{-1}zu=u'^{-1}zu'\) za vse \(z\) iz \(\mathfrak Z\), to slěduje \((u^{-1}u')^{-1}z(u^{-1}u')=z\) za vse \(z\in\mathfrak Z\). Priloženje od \(u^{-1}u'\) k \(\mathfrak Z\) davaje zato komutativno podpolje od \(\mathfrak K_r\), ktoro zaradi maksimalnosti od \(\mathfrak Z\) sběgaje sę s \(\mathfrak Z\). zato \(u^{-1}u'\) leži v \(\mathfrak Z\).

Dokazujemo slědujući teorem:
\[
\mathfrak K_r=u_1\mathfrak Z+\cdots+u_n\mathfrak Z.
\]
Za to očevidno jest potrěbna samo linearna nezavisnost od \(u_i\) po odnošenju k \(\mathfrak Z\). Dokaz se izvodi neposrědnim zadanjem od \(u_i\). Nehaj bude
\[
\mathfrak Z_Z=e_1Z+\cdots+e_nZ,\qquad
\mathfrak K_{rZ}=\mathfrak l_1+\cdots+\mathfrak l_n,\quad
\mathfrak l_i=\mathfrak K_{rZ}e_i
\]
i nehaj \(p_{ik}\) bude sistem pseudomatricnyh jedinic po odnošenju k tomu razkladu. Na str.~8 označili jesmo s \(\Theta_i\) izomorfizm od \(\mathfrak Z\) na \(Z\), dany posrědstvom modula prědstavjenja \(e_iZ\), tako že
\[
S_i=\Theta_i\Theta_1^{-1}
\]
sut \(n\) avtomorfizmov od \(Z\) (ale samo koliko se oni priměnjujut na samo \(Z\), ne vyše za tu vsečasno prědpokladano dějstvo od \(S_i\) na cělo \(\mathfrak K_{rZ}\)!). Zaradi prostoty hočemo vměsto \(\Theta_i\) v \(S_i=\Theta_i\Theta_1^{-1}\) pisati \(\Theta_{S_i}\), a togda takože pri \(p_{ik}\) zaměniti indeksy \(i,k\) prinaležečimi avtomorfizmami \(S\)
\[
p_{ik}=p_{S_iS_k}.
\]

\(n\) avtomorfizmov od \(\mathfrak Z\) sut (\(E\) jest identična substitucija)
\[
H=\Theta_E^{-1}\Theta.
\]
Numerovanje od \(H\) nyně izbirajemo tako:
\[
H_S=\Theta_E^{-1}\Theta_{S^{-1}}.
\]

Pri tyh uslovjah možno položiti

\[
u_T=\sum_S p_{S,ST}=\sum_S S(p_{E,T})
\]



1. \(u_T^{-1}\) obstaja:
\[
u_T^{-1}=\sum_S p_{ST,S}
\frac{1}{\alpha_{S,S}^{ST}\overline{\gamma}_{S,ST}\overline{\gamma}_{ST,S}}
\]

(Značenje od \(\bar{\gamma}\) vidi na str.~32).

2. \(u_T\) jest element od \(\mathfrak K_r\), zato že jest \(\mathfrak G\)-invariantny:
\[
R(u_T)=\sum_S RS(p_{E,T})=\sum_{S'} S'(p_{E,T})=u_T.
\]

3. \(u_T\) stvarjaje \(H_T\), t.~j. ako \(z\in\mathfrak Z\), to \(u_T^{-1}zu_T=H_T(z)\), ili
\[
z u_T=u_T H_T(z).
\]

Imamo

\[
z=\sum_S\Theta_S(z)e_S,\qquad
u_T=\sum_S S(e_Ep_{E,T})=\sum_S e_S S(p_{E,T})
\]

zato

\[
zu_T=\sum_{S,S'}\Theta_S(z)e_Se_{S'}S(p_{E,T})
=\sum_S\Theta_S(z)S(p_{E,T}).
\]

S drugoj strany imajemo

\[
u_T=\sum_S S(p_{E,T}e_T)=\sum_S S(p_{E,T})e_{ST}
\]

i imaje město slědujuće jednačenje za dějstvo na elementy od \(\mathfrak Z\)

\[
\Theta_S H_T=\Theta_S\Theta_E^{-1}\Theta_{T^{-1}}
=S\Theta_{T^{-1}}=ST^{-1}\Theta_E=\Theta_{ST^{-1}}
\]

zato

\[
H_T(z)=\sum_S e_S\Theta_S H_T(z)=\sum_S e_S\Theta_{ST^{-1}}(z)
\]

t.~j.

\[
\begin{aligned}
u_TH_T(z)
&=\sum_{S,S'} S(p_{E,T})e_{ST}\cdot e_{S'}\Theta_{S'T^{-1}}(z)\\
&=\sum_S S(p_{E,T})\Theta_S(z)
\end{aligned}
\]

zato

\[
z\cdot u_T=u_T\cdot H_T(z).
\]

4. \(u_S\) sut linearno nezavisne po odnošenju k \(\mathfrak Z\).

Iz
\[
\sum_T z_Tu_T=0
\]
zaradi \(z=\sum_S\Theta_S(z)e_S\) slěduje

\[
\sum_{T,S,L}\Theta_S(z_T)e_Sp_{L,LT}
=\sum_{S,T}\Theta_S(z_T)p_{S,ST}\quad\text{zato}\quad
\Theta_S(z_T)=0,
\]

zato: \(z_T=0\).

Od nyně vměsto \(H_S(z)\) pišemo \(z^S\), ,,simbolična potencija``. Za \(u_S\) morajut byti izpolnjene relacije slědujućego vida

\[
u_Ru_T=u_{RT}a_{R,T},\qquad a_{R,T}\ne0\quad\text{iz }\mathfrak Z.
\]

Zato že

\[
\begin{aligned}
u_Ru_T
&=\sum_{S,S'} S(p_{E,R})S'(p_{E,T})
=\sum_S p_{S,SR}p_{SR,SRT}\\
&=\sum_S p_{S,SRT}\alpha_{S,SRT}^{(SR)}
=\sum_S p_{S,SRT}\alpha_{S,SRT}^{(SR)}e_{SRT}
\end{aligned}
\]

i

\[a_{R,\,T} = \sum\limits_{S} e_{SRT}\, \Theta_{SRT} \left(a_{R,\,T}\right),\]

zato

\[u_{RT} a_{R,T} = \sum\limits_{S} p_{S,SRT} \Theta_{SRT} (a_{R,T})\]

slěduje

\[\Theta_{SRT}(a_{R,\,T}) = \alpha_{S,SRT}^{(SR)}, ~~ a_{R,\,T} = \sum\limits_{S} e_{SRT} \alpha_{S,SRT}^{(SR)} = \sum\limits_{L} e_{L} \alpha_{LT^{-1}R^{-1},L}^{(LT^{-1})}.\]

Sistem od \(a_{R,\,T}\) hočemo nazyvati malym faktornym sistemom od \(\mathfrak K_r\); veliki faktorny sistem \(\alpha_{BC}^{(A)}\) i mali faktorny sistem \(a_{R,\,T}\) oprěděljajut se vzajemno. K produktu \(\varepsilon_{BC}^{(A)} = \alpha_{BC}^{(A)} \beta_{BC}^{(A)}\) dvoh velikyh faktornyh sistemov prinaleži produkt \(f_{R,\,T} = a_{R,\,T}b_{R,\,T}\) malyh faktornyh sistemov \(a_{R,\,T}\), \(b_{R,\,T}\), priręđenyh k \(\alpha_{BC}^{(A)}\), \(\beta_{BC}^{(A)}\); to neposrědnje slěduje iz gornjih formul. K velikomu jedinčnomu faktornomu sistemu \(\alpha_{BC}^{(A)} = 1\) prinaleži mali jedinčny faktorny sistem \(a_{R,\,T} = 1\), i obratno.

Možno zaměniti \(u_S\) i \(u_S'=u_Sr_S\), kde \(r_S\neq 0\) jest izbrano libovoljno iz \(\mathfrak Z\). K \(u_S'\) prinaleži mali faktorny sistem \(a_{R,T}'\)

\[
\text{(1)}\qquad u_R'u_T'=u_{RT}'a_{R,T}'
\]

\(a_{R,T}^{\prime}\) nazyvajemo asociovanym s \(a_{R,T}\). Jest

\[u_R r_R u_T r_T = u_{RT} r_{RT} a'_{R,T} = u_R u_T r_R^T r_T = u_{RT} a_{R,T} r_R^T r_T\]

iz čego

\[
\text{(2)}\qquad a'_{R,T}=a_{R,T}\frac{r_R^T r_T}{r_{RT}}
\]

slěduje.

K \(u_S' = u_S r_S\) prinaleži slědujući sistem pseudomatricnyh jedinic

\[p_{S,ST}^{\prime}=p_{S,ST}r_{T}\]

mali faktorny sistem \(a'_{R,\,T}\) prinaleži zato k velikomu faktornomu sistemu \(\alpha^{(A)\prime}_{BC}\), ktory jest asociovany s faktornym sistemom \(\alpha^{(A)}_{BC}\), prinaležečim k \(a_{R,\,T}\); asociovane male faktorne sistemy odpovědajut asociovanym velikym faktornym sistemam i obratno. Priręđenje klas \(\{\alpha\}\) asociovanyh velikyh faktornyh sistemov k prinaležečim klasam asociovanyh malyh faktornyh sistemov, \(\{a\}\), jest izomorfizm grup.

Nyně prěhodimo k prědstavjenju od \(u_S\) posrědstvom matric ręda n v \(\mathfrak{Z}\), ktore se sućstveno razlikuje od dosěga razgledanyh prědstavjenij.

Ako \(\mathfrak Z\)-bazu

\[
\text{(I)}\qquad k_{S_1}, \ldots, k_{S_n}
\]

od \(\mathfrak K_r\) množimo s prava posrědstvom \(u_S\), nastava

\[
\text{(II)}\qquad k_{S_1}u_S, \ldots, k_{S_n}u_S.
\]

(Regularnu) matricu \(A_S\), ktora prěvodi (I) v (II),

\[
\text{(III)}\qquad (k_{S_1}u_S, \ldots, k_{S_n}u_S) = (k_{S_1}, \ldots, k_{S_n})A_S
\]

nazyvamo matricoju, ktora prědstavja \(u_S\). Prědstavjenje se odnosi na vse elementy od \(\mathfrak K_{r}\), ktore kako \(u_S\) stvarjajut avtomorfizm od \(\mathfrak Z/P\), t.~j. na vse elementy \(u_Sr_S\), kde \(r_S \neq 0\) jest izbrano libovoljno iz \(\mathfrak Z\). Vsi ti elementi tvoręt grupu \(\mathfrak G^*\), ktora kako normalnu podgrupu sodrživaje multiplikativnu grupu \(\mathfrak Z^*\) od \(\mathfrak Z\); faktorna grupa \(\mathfrak G^*/\mathfrak Z^*\) jest izomorfna s Galoisovoj grupoju \(\mathfrak G\) od \(\mathfrak Z/P\) tako, že vsaky element jednogo pobočnogo klasa od \(\mathfrak G^*\) po odnošenju k \(\mathfrak Z^*\) stvarjaje avtomorfizm od \(\mathfrak Z\), priręđeny tomu klasu. Ta grupa \(\mathfrak G^*\) prědstavja se upravo ukazanym načinom

\(u_S \mapsto A_S\)

Ale to prědstavjenje ne jest homomorfno, ni recipročno homomorfno; iz

\(u_S \mapsto A_S\), \(u_T \mapsto A_T\)

ne slěduje \(u_Su_T \mapsto A_SA_T\), ale zaradi

\[
\begin{aligned}
(k_{S_1}, \ldots, k_{S_n})u_Su_T
&=(k_{S_1}, \ldots, k_{S_n})A_Su_T
=(k_{S_1}, \ldots, k_{S_n})u_TA_S^T\\
&=(k_{S_1}, \ldots, k_{S_n})A_TA_S^T:\ u_Su_T\mapsto A_TA_S^T.
\end{aligned}
\]

To prědstavjenje od \(\mathfrak{G}^*\) nazyvamo prěkriženym prědstavjenjem. Ako \(h_{S_1}, \ldots, h_{S_n}\) jest druga baza od \(\mathfrak{K}_r\) po odnošenju k \(\mathfrak{Z}\), i

\[
\text{(IV)}\qquad (h_{S_1}, \ldots, h_{S_n}) = (k_{S_1}, \ldots, k_{S_n})M
\]

to iz (III) i (IV) dobyva se:

\[(h_{S_1},\ldots,h_{S_n})\,u_S=(h_{S_1},\ldots,h_{S_n})\,M^{-1}A_S\,M^S,\]

t.~j.: Ako \(u_S\) jest pri \(k\)-baze prědstavjeno posrědstvom \(A_S\), to pri \(h\)-baze jest prědstavjeno posrědstvom \(M^{-1}A_SM^S\), ako \(M\) jest matrica, ktora transformuje \(k\)-bazu v \(h\)-bazu.

Formula (III), ktora definujeme prěkriženo prědstavjenje, pokazuje, že prěkriženo prědstavjenje stoji v podobnom odnošenju k prstenu \(\mathfrak K_r\), kako obyčno recipročno prědstavjenje (ale zapisano s prava vměsto s lěva) k svojemu modulu prědstavjenja. Razlika nastava, zato že vměsto relacije komutativnosti

\[(m^*\tau^*)c = (m^*c)\tau^*\]
(vidi str.~5)

pri recipročnom modulu prědstavjenja tu

\[
(kz)u_S=(ku_S)z^S\qquad k\text{ v }\mathfrak K_r,\quad z\text{ v }\mathfrak Z
\]

nastupaje.

Prěkriženo prědstavjenje od \(\mathfrak{G}^*\) možno takože obrazovati posrědstvom bazy nekogo pravogo ideala \(\mathfrak r\) od \(\mathfrak{K}_r\). Vsaky pravy ideal od \(\mathfrak{K}_r\) stvarjaje na toj način klasu ekvivalentnyh prědstavjenij; pri tom prědstavjenja \(u_S \mapsto A_S, u_S \mapsto B_S\) nazyvajut se ekvivalentnymi, ako obstaja matrica \(M\) s \(B_S = M^{-1}A_SM^S\).

Vměsto ograničenja na prědstavjenja od \(\mathfrak G^*\) v \(\mathfrak Z\), stvorjene pravymi idealami od \(\mathfrak K_r\), nyně možno obće definirati prěkriženo prědstavjenje \(m\)-togo stupenja od \(\mathfrak G^*\) v \(\mathfrak Z\) tako:

\begin{enumerate}
\item Vsakomu \(u_S\) odpovědaje matrica \(A_S\) ręda \(m\) nad \(\mathfrak Z\).
\item Iz \(u_S \mapsto A_S\), \(u_T \mapsto A_T\) slěduje \(u_Su_T \mapsto A_TA_S^T\).
\end{enumerate}

Izberemo pravy \(\mathfrak Z\)-modul ranga \(m\)

\[\mathfrak{M}=x_1\mathfrak{Z}+\cdots+x_m\mathfrak{Z}\]

i ustanovjenjem

\[(x_1u_S,\ldots,x_mu_S)=(x_1,\ldots,x_m)A_S\]

činimo go pravym \(\mathfrak G^*\)-modulom s relacijoju komutativnosti

\[(mz) u_S = (mu_S) z^S\]

cělkom analogno konstrukciji modula prědstavjenja k zadanomu prědstavjenju na str.~5. Nyně pokazujemo, kako na str.~20, že \(\mathfrak{M}\) postavaje konečnym pravym \(\mathfrak K_r\)-modulom posrědstvom ustanovjenja

\[m \cdot z u_S = mz \cdot u_S = mu_S \cdot z^S\]

obće

\[m \cdot \sum z_S u_S = \sum m z_S \cdot u_S = \sum m u_S \cdot z_S^S.\]

Pri tom nikako ne jest potrěbno računanje analogno računanju na str.~20, zato že pri našem ustanovjenju upotrěbili jesmo ustanovjenu bazu.

Ale po M.Z. § 18 vsaky konečny prosty \(\mathfrak K_r\)-modul \(\mathfrak{M}\) jest operatorno izomorfny s prostym pravym idealom od \(\mathfrak K_r\); iz toga analogno teoremu 2 na str.~22 zaključujemo, že obstaja samo \emph{jedna} klasa ireducibilnyh prěkriženyh prědstavjenij od \(\mathfrak{G}^*\) v \(\mathfrak Z\) --- jmenno ona, ktora prinaleži k prostym pravym idealam.

Stupenj toga ireducibilnogo prědstavjenja jest \(t\), kde \(t\) jest indeks (absolutno čislo eksponenta) od \(\mathfrak K\). To se dobyva prostym sčitanjem rangov. Zato že stupenj jest ravny rangu \((\mathfrak r:\mathfrak Z)\) prostogo pravogo ideala nad \(\mathfrak Z\). Nyně na str.~22 imaje město:

\[
(\mathfrak Z:P)=n=rt;\qquad (\mathfrak K_r:P)=n^2;\qquad (\mathfrak r:P)=n\cdot t
\]

poslědnje zaradi \(\mathfrak K_r = \mathfrak{r}_1 + \cdots + \mathfrak{r}_r\).

Iz
\[(\mathfrak{r}: P) = (\mathfrak{r}: \mathfrak Z) \cdot (\mathfrak Z: P)\]
ili \(n \cdot t = (\mathfrak{r}: \mathfrak Z) \cdot n\) slěduje zato \((\mathfrak{r}: \mathfrak Z) = t\).
% END BOOK_S25

% BEGIN BOOK_S26
\subsection*{§ 26. Množenje prěkriženyh prědstavjenij}\label{multiplikation-der-verschraenkten-darstellungen}

Doseli jesmo vsegda razgledali prěkriženo prědstavjenje kako prědstavjenje od \(\mathfrak{G}^*\). Ale jego možno razgledati takože kako prědstavjenje Galoisovoj grupy \(\mathfrak{G}\) od \(\mathfrak Z/P\), i to tako, že vsakomu elementu \(S\) od \(\mathfrak{G}\) priręđajemo matricu \(A_S\), ktora prědstavja element \(u_S\) iz \(\mathfrak K_r\), stvarjajuči avtomorfizm \(S\) posrědstvom transformacije -- \(u_S^{-1}zu_S=z^S\) --.

Za označenje takogo prědstavjenja \(S\mapsto A_S\)\srcfn{1}{Tako prědstavjenje (nezavisno od hiperkompleksnyh sistemov)
\[
S\mapsto A_S;\quad T\mapsto A_T;\quad ST\mapsto A_TA_S^T=A_{ST}\cdot a_{S,T}^{-1}
\]
(ili transponovano prědstavjenje) bylo odpravnym punktom Speisera (Math. Zschr. 5) pri oprěděljenju faktornyh sistemov \(a_{S,T}\). Te faktorne sistemy pokazyvajut se zato identičnymi s tu iz drugoj osnovy uvedenymi.} trěba skazati, v ktorom \(\mathfrak K_r\) ono jest vzęto, ktore \(u_S\) i ktore \(k_{S_1},\ldots,k_{S_n}\) od \(\mathfrak K\) nad \(\mathfrak Z\) sut upotrěbjene. Zaměna bazy davaje ekvivalentno prědstavjenje, a zaměna od \(u\) davaje asociovane faktorne sistemy \(a_{S,\,T}\); sućstvene sut to jest \(\mathfrak K_r\), v ktoryh prědstavjenje jest vzęto, t.~j. klasa \(\{a_{S,\,T}\}\) asociovanyh faktornyh sistemov.

Nehaj nyně budut dane dva taka prědstavjenja od \(\mathfrak G\).

\[
\begin{array}{lll}
S\mapsto u_S\mapsto A_S & \text{naležeči mali faktorny sistem} & a_{R,T},\\
S\mapsto v_S\mapsto B_S & \text{naležeči mali faktorny sistem} & b_{R,T}.
\end{array}
\]

Pod Kronekerovym produktom \(A\times B\) dvoh matric \(A=(a_{ik})\), \(B=(b_{\mu\nu})\) razumějemo matricu \(A\times B=(a_{ik}b_{\mu\nu})\).

\noindent\textbf{Tvrdženje.} \emph{Takože \(S\mapsto A_S\times B_S\) jest prěkriženo prědstavjenje od \(\mathfrak G\), i ono prinaleži k produktu \(\mathfrak K_r\times\mathfrak S_s\), ako \(A_S\) prinaleži k \(\mathfrak K_r\), a \(B_S\) k \(\mathfrak S_s\).}

\noindent\emph{Dokaz.} Prědstavjenje \(S\mapsto A_S\times B_S\) jest prěkriženo prědstavjenje. Nehaj bude
\[
S\mapsto A_S\times B_S,\qquad T\mapsto A_T\times B_T
\]
i
\[
ST\mapsto A_{ST}\times B_{ST}
=(A_TA_S^T\times B_TB_S^T)a_{S,T}^{-1}b_{S,T}^{-1},
\]
zato trěba dokazati:
\[
(A_T\times B_T)(A_S\times B_S)^T a_{S,T}^{-1}b_{S,T}^{-1}
=(A_TA_S^T\times B_TB_S^T)a_{S,T}^{-1}b_{S,T}^{-1}.
\]
To ale slěduje iz slědujućego razvažanja:

Ako matrice \(A\) i \(B\) prědstavimo posrědstvom matričnyh jedinic \(c_{ik}\) i \(d_{\mu\nu}\) (ktore ne imajut ničto obće s matričnymi jedinicami od \(\mathfrak K_r\) ili \(\mathfrak S_s\))
\[
A=\sum c_{ik}\alpha_{ik};\qquad
B=\sum d_{\mu\nu}\beta_{\mu\nu},\qquad\text{togda}
\]
\[
A\times B=B\times A
=\sum c_{ik}d_{\mu\nu}\cdot\alpha_{ik}\cdot\beta_{\mu\nu},
\]
t.~j. ono jest ravno matrici v novyh matričnyh jedinicah
\(c_{ik}d_{\mu\nu}=d_{\mu\nu}c_{ik}\). Množenje \(A\)-matric i \(B\)-matric zato jednoznačno odpovědaje množenju elementov v prěmom produktu dvoh matričnyh prstenov:
\[
\sum c_{ik}\mathfrak Z\times\sum d_{\mu\nu}\mathfrak Z.
\]
Zato vsaka \(A\)-matrica jest komutativna s vsakoj \(B\)-matricoj (ale \(A\)-matrice ni \(B\)-matrice ne sut medžu soboju komutativne); i ta komutativnost vsakoj \(A\)-matrice s vsakoj \(B\)-matricoju davaje upravo --- po množenju s \(a_{S,T}^{-1}b_{S,T}^{-1}\) --- relaciju, ktoru trěba dokazati.

\section*{Glava VI. Teorija prěkriženyh produktov}\label{kapitel-vi.-theorie-der-verschraenkten-produkte}
% END BOOK_S26

% BEGIN BOOK_S27
\subsection*{§ 27. Prědstavjenje od \(\mathfrak K_r\) kako prěkriženyh produktov}

Teorija prěkriženyh produktov, razvijena v § 24 iz faktornyh sistemov, nyně bude ponovno nezavisno osnovana. Značenje tyh prěkriženyh produktov jest v tom, že Galoisovo polje i jego grupy sut \emph{zajedno} opisane v \(\mathfrak K_r\).

\noindent\textbf{Teorem 1.} \emph{Ako \(\mathfrak Z\) jest Galoisovo razpadno polje od \(\mathfrak K_r\), to jest maksimalno komutativno podpolje, i ako \(\mathfrak G\) jest jego Galoisova grupa, to \(\mathfrak K_r\) jest ravno prěkriženomu produktu \(\mathfrak Z\cdot\mathfrak G\). Obratno, vsaky prěkriženy produkt \(\mathfrak Z\cdot\mathfrak G\) stvarjaje \(\mathfrak K_r\) takogo vida.}

\noindent\textbf{Dokaz i definicije. Definicija 1.} Nehaj \(\mathfrak G=\{E,S,\ldots,T\}\) bude Galoisova grupa od \(\mathfrak Z\); ravnost od \(\mathfrak K_r\) prěkriženomu produktu znači:
\[
\mathfrak K_r=\mathfrak Zu_E+\cdots+\mathfrak Zu_T
\]
s \(zu_S=u_Sz^S\) i \(u_Su_T=a_{S,T}u_{ST}\) (\(a_{S,T}\) v \(\mathfrak Z\) sut male faktorne sistemy).

\noindent\textbf{Definicija 2.} Stvarjajuča grupa \(\mathfrak G^*\) jest oprěděljena kako množstvo regularnyh elementov \(g\) iz \(\mathfrak K_r\), ktore posrědstvom transformacije prěvodet \(\mathfrak Z\) kako cělost v sebe. \(g^{-1}\mathfrak Zg=\mathfrak Z\), to jest \(g^{-1}zg=z^S\), s \(S\) v \(\mathfrak G\). zato \(\mathfrak G\) jest grupovo-homomorfny obraz od \(\mathfrak G^*\), zato že vsaky avtomorfizm od \(\mathfrak Z\) jest stvorjen nekym \(g\), \(\mathfrak G^*\to\mathfrak G\), i dalje
\[
\mathfrak G\cong\mathfrak G^*/\mathfrak T^*
\quad\text{s}\quad
\mathfrak T^*\supseteq\mathfrak Z^*,
\]
(že \(\mathfrak T^*=\mathfrak Z^*\), bude dokazano pozděje --- ješče ne bylo dokazano, že \(\mathfrak Z\) bylo takože maksimalno komutativno podprsten. To možno dokazati i nezavisno).

\noindent\textbf{1.} \emph{Dokaz prědstavjenja kako prěkriženogo produkta.} Nehaj \(u_E,u_S,\ldots,u_T\) budut elementi iz \(\mathfrak K_r\), ktore stvarjajut avtomorfizmy \(E,\ldots,T\), to jest \(zu_S=u_Sz^S\), osoblivo \(u_E=e\) zaradi 2. Trěba pokazati, že suma
\[
\mathfrak M=\{\mathfrak Zu_E,\ldots,\mathfrak Zu_T\}
\]
jest prěma i ravna \(\mathfrak Zu_E+\cdots+\mathfrak Zu_T\); togda ona jest ravna \(\mathfrak K_r\) zaradi jednakosti rangov. To ale slěduje tako (\foreign{Schwarz}):

\(\mathfrak M\) jest lěvy i pravy \(\mathfrak Z\)-modul, to jest modul prědstavjenja od \(\mathfrak Z\) v \(\mathfrak Z\), i posrědstvom \(\mathfrak M\) stvarjajut se \(n\) različne prědstavjenja; tym samym --- zato že \(\mathfrak Z\) jest komutativno polje prvogo vida --- \(\mathfrak M\) imaje rang prinajmenje \(n\); inače dokaz teče upravo tako.

\noindent\textbf{2.} \(\mathfrak G\cong\mathfrak G^*/\mathfrak Z^*\), t.~j. \(\mathfrak T^*=\mathfrak Z^*\); samo elementi iz \(\mathfrak Z\) inducirajut identitet na \(\mathfrak Z\) posrědstvom transformacije.

Nehaj imenno bude \(tz=zt\) za vsako \(z\) iz \(\mathfrak Z\); položimo
\[
t=z_0+z_1u_S+\cdots+z_{n-1}u_T,
\]
zato
\[
zt=zz_0+\cdots+zz_{n-1}u_T
\]
i
\[
tz=z_0z+\cdots+z_{n-1}u_Tz
=z_0z+z^{S^{-1}}z_1u_S+\cdots+z^{T^{-1}}z_{n-1}u_T,
\]
zato zaradi prěmoj sumy:
\[
(z-z^{S^{-1}})z_1u_S=0,\ldots,
(z-z^{T^{-1}})z_{n-1}u_T=0
\]
i iz toga, zato že elementi \(u_S,\ldots\) sut regularne:
\[
(z-z^{S^{-1}})z_1=0,\ldots,
(z-z^{T^{-1}})z_{n-1}=0,
\qquad\text{za vsako }z.
\]
(ili zato že rang jest ravny \(n\)). Ale po definiciji za \(S,\ldots\), i tako dalje (vse neravne \(E\)), obstaja prinajmenje jedno \(z\), tako že \(z-z^{S^{-1}}\neq0\); i neko \(\bar z\), tako že \(\bar z-\bar z^{T^{-1}}\neq0\); zato dobyva se \(z_i=0\), \(i=1,\ldots,(n-1)\), i zato \(t=z_0\) v \(\mathfrak Z^*\), tym jest dokazano.

\noindent\textbf{3.} Elementi \(u\) izpolnjajut zakony množenja: \(u_Su_T=a_{S,T}u_{ST}\), pri čem \(a_{S,T}\) izpolnjajut prěkriženy asociativny zakon
\[
a_{R,S}\cdot a_{RS,T}=a_{S,T}^{R^{-1}}\cdot a_{R,ST}
\]
(\(a_{S,T}\) sut mali faktorny sistem).

Zato že po 2. dobyva se:
\[
\mathfrak Z^*u_S\cdot\mathfrak Z^*u_T=\mathfrak Z^*u_{ST};
\]
zato
\[
u_Su_T=a_{S,T}u_{ST};
\]
s \(a_{S,T}\) iz \(\mathfrak Z^*\), to jest iz \(\mathfrak Z\) i \(\neq0\). Zaradi \(u_R\cdot u_Su_T=u_Ru_S\cdot u_T\) dalje dobyva se
\[
\begin{aligned}
u_Ra_{S,T}u_{ST}
&=a_{S,T}^{R^{-1}}\cdot u_Ru_{ST}
 =a_{S,T}^{R^{-1}}a_{R,ST}u_{RST}\\
&=a_{R,S}u_{RS}\cdot u_T
 =a_{R,S}a_{RS,T}u_{RST},
\end{aligned}
\]
Tym sut podane zakony množenja.

\noindent\textbf{4. Obrat:} prěkriženy produkt \(=\mathfrak K_r\). Prěma suma:
\[
\mathfrak M=\mathfrak Zu_E+\cdots+\mathfrak Zu_T
\]
bude oprěděljena kako prsten posrědstvom relacij: \(zu_S=u_Sz^S\), \(u_Su_T=a_{S,T}u_{ST}\), (kde \(a_{S,T}\) izpolnjajut prěkrižene asociativne relacije; to jest \(u_Ru_Su_T\) jest jednoznačno oprěděljeno):
\[
\begin{array}{ll}
1.&z_1\cdot z_2u_S=z_1z_2\cdot u_S;\\[2pt]
2.&zu_S\cdot u_T=z\cdot u_Su_T
\end{array}
\]
tym \(\mathfrak M\) postavaje asociativnym prstenom.

\noindent\textbf{Obratny teorem 2.} \emph{Prěkriženy produkt \(\mathfrak Z\cdot\mathfrak G\) jest dvustranno prosty, s \(P\) kako centrom (po pristupah R. Brauera).} Dokaz osnovyvaje se na slědujućih pomoćnyh lemah:

\noindent\textbf{Pomoćna lema 1.} \emph{Vsaky element, ktory jest po elementah komutativny s \(\mathfrak Z\), prinaleži k \(\mathfrak Z\).}

\noindent\emph{Dokaz.} Nehaj \(w=\sum_S c_{\lambda_S}u_S\), (\(c_{\lambda_S}\) v \(\mathfrak Z\)) i nehaj bude prědpoloženo: \(wz=zw\) za vsako \(z\) iz \(\mathfrak Z\). zato
\[
\sum_S zc_{\lambda_S}u_S
 =\sum_S c_{\lambda_S}u_Sz
 =\sum_S c_{\lambda_S}z^{S^{-1}}u_S
\]
zato \(zc_{\lambda_S}=c_{\lambda_S}z^{S^{-1}}\), ili \(c_{\lambda_S}(z-z^{S^{-1}})=0\), za vsako \(z\) i \(S\) (linearna nezavisnost od \(u\)!). Iz toga ale slěduje: \(c_{\lambda_S}=0\) za \(S\neq E\), zato že za vsako tako \(S\) obstaja prinajmenje \emph{jedno} \(z\), tako že \(z-z^{S^{-1}}\neq0\); zato \(w\) jest v \(\mathfrak Z\).

\noindent\textbf{Pomoćna lema 2.} \emph{Vsaky \(\mathfrak Z\)-dvojny modul \(\mathfrak A\), ktory jest dvojno izomorfny s \(\mathfrak Zu_S=u_S\mathfrak Z\) (za \(\mathfrak Z\) kako lěvy i pravy operator), jest identičny s \(\mathfrak Zu_S\).}

\noindent\emph{Dokaz 2.} \(\mathfrak A\) jest izomorfny s \(\mathfrak Zu_S\) kako lěvy modul; ako zato \(u_S\mapsto a\) iz \(\mathfrak A\), to \(\mathfrak Zu_S\mapsto\mathfrak Za\), to \(\mathfrak A=\mathfrak Za\). Iz dvojnogo izomorfizma slěduje (1) \(za=az^S\) (zaradi \(u_Sz\mapsto az\)); i iz toga \(\mathfrak A=\mathfrak Zu_S\). Ako imenno položimo: \(a=u_Su_S^{-1}a=u_Sw\), dobyva se: \(za=zu_Sw=u_Sz^Sw\), ale zaradi (1): \(za=u_Swz^S\), zato (po množenju s \(u_S\)): \(z^Sw=wz^S\) za vsako \(z\) iz \(\mathfrak Z\). Zato takože \(z^S\) proběgaje cělo \(\mathfrak Z\), \(w\) jest po elementah komutativny s \(\mathfrak Z\), zato po 1. prinaleži k \(\mathfrak Z\).

\noindent\textbf{Pomoćna lema 3.} \emph{Vsaky \(\mathfrak Z\)-dvojny modul iz \(\mathfrak Z\cdot\mathfrak G\) imaje formu: \(\mathfrak Zu_{S_1}+\cdots+\mathfrak Zu_{S_k}\), kde \(u_{S_1}\) označaje nykoje \(k\) različne iz \(u_S,\ldots,u_T\).}

\noindent\emph{Dokaz.} Prěkriženy produkt \(\mathfrak Z\cdot\mathfrak G\) jest polno reducibilny kako \(\mathfrak Z\)-dvojny modul, zato že vsaky sumand \(\mathfrak Zu_S\) imaje rang 1. Kako dvojne moduly, vsi sumandy takože prinaležet k različnym klasam, zato ne sut dvojno izomorfne. Zaradi polnoj reducibilnosti vsaky \(\mathfrak Z\)-dvojny modul \(\mathfrak M\) jest prěmy sumand i sam polno reducibilny. \(\mathfrak M=\sum\mathfrak A_{S_i}\), kde \(\mathfrak A_{S_i}\) jest izomorfny s \(\mathfrak Zu_{S_i}\); i zato po 2. jest identičny s \(\mathfrak Zu_{S_i}\). (Vsi \(\mathfrak A_{S_i}\) prinaležet k različnym klasam.)

\noindent\textbf{Pomoćna lema 4.} \emph{Vsako prsten iz \(\mathfrak Z\cdot\mathfrak G\), ktoro sodrživaje \(\mathfrak Z\), osoblivo samo \(\mathfrak Z\cdot\mathfrak G\), jest dvustranno prosto. --- Podprsteny, ktore sodrživajut \(\mathfrak Z\), jednoznačno odpovědajut podgrupam od \(\mathfrak G\).}

\noindent\emph{Dokaz.} Vsako prsten \(\mathfrak o\), ktoro sodrživaje \(\mathfrak Z\), jest \(\mathfrak Z\)-dvojny modul, zato po 3. imaje formu \(\sum\mathfrak Z\cdot u_{S_i}\); zato že \(\mathfrak o\) jest prsten, produkt dvoh \(u_{S_i}\) ponovno prinaleži k \(\mathfrak o\), zato takože prinaležeči \(\mathfrak Z\)-modul. Zato \(u_{S_i}\) tvoręt grupu do faktornyh sistemov, zato že jih jest konečno mnogo. zato \(\mathfrak o=\mathfrak Z\cdot\mathfrak H\), kde \(\mathfrak H\) jest podgrupa od \(\mathfrak G\). (Točnije, \(\mathfrak Z\cdot u_{S_i}\) tvoręt multiplikativny sistem.) Ako nyně \(\mathfrak a\) jest libovoljny dvustranny ideal \(\neq0\) iz \(\mathfrak o\), to on jest \(\mathfrak o\)-dvojny modul, i tym \(\mathfrak Z\)-dvojny modul i podmodul od \(\mathfrak o\). zato vměstě s \(\mathfrak Zu_{S_i}\) sodrživaje cělu grupu \(\mathfrak H\), i zato postavaje \(\mathfrak o=\mathfrak Z\cdot\mathfrak H\).

\noindent\textbf{Pomoćna lema 5.} \emph{\(P\) jest centar od \(\mathfrak Z\cdot\mathfrak G\).}

\noindent\emph{Dokaz.} Ako \(w\) jest po elementah komutativny s \(\mathfrak Z\cdot\mathfrak G\), to osoblivo s \(\mathfrak Z\), zato po 1. prinaleži k \(\mathfrak Z\). Zato iz \(wu_S=u_Sw^S\) slěduje \(w=w^S\), ako jest komutativny. zato \(w\) prinaleži k \(P\).

\noindent\emph{Primětka 1.} Dokaz pokazuje, že teorem imaje město takože, kogda \(\mathfrak Z\) jest beskonečno algebraično Galoisovo polje nad \(P\). Zato že teoremy o polnoj reducibilnosti ostavajut v silě (Krull Zahlringe, Math. Ann. 99). Pri tom \(\mathfrak Z\cdot\mathfrak G\) jest sistem vsih konečnyh sum \(\sum_{S_\lambda}u_{S_\lambda}\cdot c_{S_\lambda}\).

Samo pri 4. trěba za \(\mathfrak o\) prědpoložiti \(\mathfrak o=\mathfrak Z\cdot\mathfrak H\), a to vsegda imaje město za polny produkt \(\mathfrak Z\cdot\mathfrak G\), zato že iz naleženja od \(u_{S_i}u_{S_k}\) nyně ne možno zaključiti grupovo svojstvo sistema, ako \(\mathfrak o\) jest beskonečno nad \(\mathfrak Z\).

\noindent\emph{Primětka 2.} Pri prěmom osnovanju tu trěba prisojediniti teoriju prěkriženogo matričnogo prědstavjenja, razvijenu v §§ 25/26. V slědujućem, na koncu § 28, upotrěbja se, že ireducibilno prěkriženo prědstavjenje imaje stupenj \(t\), kde \(t\) jest absolutno čislo komponentov (indeks) od \(\mathfrak K\). Ponovno trěba ukazati, že ta teorija prědstavjenja davaje identitet faktornyh sistemov, uvedenyh tu kako konstanty množenja, s pojmom obyčnym v literaturě.
% END BOOK_S27

% BEGIN BOOK_S28
\subsection*{§ 28. Produktny teorem za faktorne sistemy}\label{produktsatz-fuer-faktorensysteme}

Nehaj bude
\[
\mathfrak K_r=\mathfrak Z\cdot\mathfrak G
=\mathfrak Zu_E+\cdots+\mathfrak Zu_T,
\qquad u_Su_T=u_{ST}a_{S,T},
\]
\[
\overline{\mathfrak K}_r
=\overline{\mathfrak Z}\,\overline{\mathfrak G}
=\overline{\mathfrak Z}\bar u_E+\cdots+\overline{\mathfrak Z}\bar u_T,
\qquad
\bar u_S\bar u_T=\bar u_{ST}\bar b_{S,T},
\]
(\(\mathfrak Z\) i \(\overline{\mathfrak Z}\) sut izomorfne Galoisove polja, a \(\mathfrak G\), \(\overline{\mathfrak G}\) njih grupy), zato dobyva se
\[
\mathfrak K_r\times\overline{\mathfrak K}_r
=\widetilde{\mathfrak Z}\,\widetilde{\mathfrak G}\times P_n
=\mathfrak A_f\times P_n
=\mathfrak A_j
\]
s \(\tilde u_S\tilde u_T=\tilde u_{ST}\tilde c_{S,T}\) (\(\widetilde{\mathfrak Z}\) jest izomorfno s \(\mathfrak Z\) i \(\overline{\mathfrak Z}\)). To slěduje iz računanja rangov i iz toga, že \(\mathfrak Z\) jest takože razpadno polje).

\noindent\textbf{Tvrdženje.} \emph{Dobyva se \(\tilde c_{S,T}=\tilde a_{S,T}\tilde b_{S,T}\), kde \(\tilde a,\tilde b\) sut s \(a\) i \(\bar b\) izomorfne elementi v \(\widetilde{\mathfrak Z}\).}

Dokaz osnovyvaje se na tom, že prsten avtomorfizmov libovoljnogo lěvogo ideala iz \(\mathfrak A_j\), ktory imaje absolutnu dolžinu \(n\), jest izomorfno s \(\mathfrak A_f=\widetilde{\mathfrak Z}\widetilde{\mathfrak G}\), i že to prsten avtomorfizmov možno prěmo oprěděliti pri pogodnom lěvom idealu iz \(\mathfrak A_j\). Obće imaje město slědujući teorem, imenno s \(s\) vměsto \(n\), ale zaradi udobnosti izrečeny za \(\mathfrak K_r\):

\noindent\textbf{1. Teorem 1.} \emph{Prsten avtomorfizmov lěvogo ideala \(\mathfrak l\) absolutnoj dolžiny \(n\) iz \(\mathfrak K_{rs}=\mathfrak K_r\times P_s\) jest izomorfno s \(\mathfrak K_r\).}

Dokaz slěduje iz slědujućih pomoćnyh lem.

\noindent\textbf{Pomoćna lema 1.} \emph{Vsi lěve idealy absolutnoj dolžiny \(n\) iz \(\mathfrak K_{rs}\) sut operatorno izomorfne. Njih prsteny avtomorfizmov sut zato prstenovo izomorfne. Osoblivo \(\mathfrak l=\mathfrak K_r\times\mathfrak b\) imaje absolutnu dolžinu \(n\), pri čem pod \(\mathfrak b\) razumějemo prosty ideal iz \(P_s\).}

\noindent\emph{Dokaz.} To, že lěve idealy imajut absolutnu dolžinu \(n\), t.~j. razpadajut se na \(n\) absolutno nerazložimyh prostyh idealov, imaje za poslědicu jednaku dolžinu v \(\mathfrak K_{rs}\), t.~j. v \(\mathfrak K_{rs}\) nastupaje ravnako mnogo prostyh idealnyh sumandov. Iz toga slěduje njih operatorny izomorfizm, a tym i prstenovy izomorfizm prstenov avtomorfizmov. Dalje \(\mathfrak K_r=\mathfrak Z\cdot\mathfrak G\) imaje absolutnu dolžinu \(n\), t.~j. rang jest \(n^2\), zato \(\mathfrak K_r\times P_s\) imaje absolutnu dolžinu \(n\cdot s\), i zato \(\mathfrak K_r\times\mathfrak b\) imaje absolutnu dolžinu \(n\).

\noindent\textbf{Pomoćna lema 2.} \emph{Ako \(\mathfrak o\) jest prsten s jedinicoju, \(\mathfrak l=\mathfrak o e_1\) jest lěvy ideal i prěmy sumand, to jest \(e_1\) jest idempotent, to prsten avtomorfizmov od \(\mathfrak l\) jest izomorfno s \(e_1\mathfrak o e_1\).}

\noindent\emph{Dokaz.} Množenje s elementom \(e_1re_1\) davaje operatorny homomorfizm od \(\mathfrak l\) v sebe; različne elementi \(e_1re_1\) davajut različne homomorfizmy, zato že pri tom k \(e_1\) sut priręđene različne elementi, i vsaky homomorfizm jest stvorjen tako: \(e_1\mapsto re_1\) i \(re_1=e_1re_1\) zaradi \(e_1e_1=e_1\). Suma i produkt v \(e_1\mathfrak o e_1\) ale odpovědajut sumi i produktu priręđenyh homomorfizmov.

\noindent\textbf{Pomoćna lema 3.} \emph{Ako položimo \(P_s=\sum_1^s c_{ik}P\) s matričnymi jedinicami \(c_{ik}\) i \(\mathfrak b=c_{11}P+\cdots+c_{s1}P\), to kako prsten avtomorfizmov od \(\mathfrak K_r\times\mathfrak b\) dobyva se prsten \(\mathfrak K_rc_{11}\), ktoro jest izomorfno s \(\mathfrak K_r\).}

\noindent\emph{Dokaz.} Zato že \(c_{11}\) jest po elementah komutativny s \(\mathfrak K_r\) i nijedin element iz \(\mathfrak K_r\) ne jest anulovany, zato že \(c_{11}\) jest stvarjajuči idempotent od \(\mathfrak K_r\times P_s\), i jest
\[
c_{11}\cdot\mathfrak K_r\times P_s\cdot c_{11}
=\mathfrak K_r\times c_{11}P_sc_{11}
=\mathfrak K_rc_{11}.
\]
Iz tyh pomoćnyh lem prěmo slěduje teorem 1.

\noindent\textbf{2. Oprěděljenje i razklad od \(\mathfrak K_r\times\overline{\mathfrak K}_r\) na lěve idealy absolutnoj dolžiny \(n\).} Po definiciji dobyva se
\[
\mathfrak A_f=\mathfrak K_r\times\overline{\mathfrak K}_r
=\mathfrak Z\cdot\mathfrak G\times\overline{\mathfrak Z}\cdot\overline{\mathfrak G},
\]
pri čem element s črtoju jest po elementah komutativny s elementom bez črty, zato
\[
\mathfrak A_f
=(\mathfrak G\times\overline{\mathfrak G})\cdot(\mathfrak Z\times\overline{\mathfrak Z})
=\mathfrak G\times\overline{\mathfrak G}\cdot(\mathfrak Ze_1+\cdots+\mathfrak Ze_n)
\]
s \(\mathfrak Ze_i=\overline{\mathfrak Z}e_i\), (množenje Galoisovogo polja samogo s soboju). Zato že vsi sumandy imajut jednaky rang po odnošenju k \(\mathfrak Z\), a zato takože jednaku absolutnu dolžinu --- absolutna dolžina od \(\mathfrak A_f\) ale jest \(=n^2\) --- vsaky sumand imaje absolutnu dolžinu \(n\). V slědujućem za osnovu bude vzęto
\[
\mathfrak l=(\mathfrak G\times\overline{\mathfrak G})\mathfrak Ze_1,
\]
pri čem relacija \(ze_1=e_1\bar z\) ustanovjaje oprěděljeny izomorfizm medžu \(\mathfrak Z\) i \(\overline{\mathfrak Z}\).

\noindent\textbf{3. Prěkriženo prědstavjenje od \(e_1\mathfrak A_fe_1\), i tym dokaz produktnogo teorema.} --- Po 1. i 2. faktorne sistemy produkta \(\mathfrak K_r\times\overline{\mathfrak K}_r\) sut izomorfne s faktornymi sistemami prěkriženogo prědstavjenja od \(e_1\mathfrak A_fe_1=\mathfrak A\).

\noindent\textbf{Teorem 2.} \emph{\(\mathfrak A=\widetilde{\mathfrak Z}\cdot\widetilde{\mathfrak G}\) s \(\widetilde{\mathfrak Z}=\mathfrak Ze_1=\overline{\mathfrak Z}e_1\), \(\widetilde{\mathfrak G}=\{\ldots,e_1u_S\bar u_Se_1,\ldots\}\), pri čem
\[
e_1u_S\bar u_Se_1\cdot e_1u_T\bar u_Te_1
=e_1u_{ST}\bar u_{ST}e_1\cdot a_{S,T}e_1\bar b_{S,T}e_1,
\]
tako že s \(\tilde a_{S,T}=a_{S,T}e_1\) i \(\tilde b_{S,T}=\bar b_{S,T}e_1\) produktny teorem jest dokazan; (\(\tilde a\) i \(\tilde b\) sut v tom samom polju \(\widetilde{\mathfrak Z}\)).}

\noindent\emph{Dokaz.} \(e_1\) postavaje jedinica prstena; elementi konjugovane s \(\tilde z=ze_1\) sut zato dane posrědstvom \(\tilde z^Se_1\); trěba pokazati, že elementi \(u_S\bar u_Se_1\) inducirajut te substitucije, čime bude dokazano prěkriženo prědstavjenje.

\noindent\textbf{Pomoćna lema 1.} \(e_1u_S\bar u_S=u_S\bar u_Se_1\), zato \(=e_1u_S\bar u_Se_1\).

\noindent\emph{Dokaz.} Položimo \(u_S^{-1}e_1u_S=e^S\)\srcfn{1}{Tym ustanovjenjem dobyva se \(e_1=e^E\).} ili \(e_1u_S=u_Se^S\); togda \(e^S\) proběgajut konjugovane \(e_1,e_2,\ldots,e_n\). Togda imaje město (1)\srcfn{2}{zaradi \(\sum a_i\bar A_i\bar u_R=\sum a_i\bar u_R\bar A_i^R=\bar u_R\sum(a_i^R)^{R^{-1}}\bar A_i^{-R}\).}
\[
e_1\bar u_R=\bar u_Re^{R^{-1}}.
\]
Zato že \(e_1=\sum a_i\bar A_i=\sum a_i^R\bar A_i^R\), kde \(a_i\) proběgaje libovoljnu \(\mathfrak Z\)-bazu, a \(\bar A_i\) prinaležeču komplementarnu \(\overline{\mathfrak Z}\)-bazu, zato to samo imaje město za \(a_i^R\) i \(\bar A_i^R\). zato \(e_1\bar u_R=\bar u_Re^{R^{-1}}\), čime (1) jest dokazano. Iz (1) slěduje
\[
e_1u_S\bar u_S=u_Se^S\bar u_S=u_S\bar u_S\cdot e^{SS^{-1}}=u_S\bar u_Se_1,
\]
čime jest dokazana pomoćna lema 1.

\noindent\textbf{Pomoćna lema 2.} \(ze_1\cdot e_1u_S\bar u_Se_1=e_1u_S\bar u_Se_1\cdot z^Se_1\).

\noindent\emph{Dokaz.} 
\[
ze_1\cdot e_1u_S\bar u_Se_1
=zu_S\bar u_Se_1\quad(\text{Pomoćna lema 1})
=u_Sz^S\bar u_Se_1
=u_S\bar u_Se_1\cdot z^Se_1
=b=e_1b,
\]
(\(z\) jest komutativny s \(\bar u\) i \(e_1\)). --- Nehaj \(u'=u_S\bar u_Se_1\). Pomoćna lema 2 izreka, že \(u'\) indiciraje substitucije od \(\mathfrak Ze_1\), tako že imaje město prěkriženo prědstavjenje podano v teoremu.

\noindent\textbf{Pomoćna lema 3.} \emph{Faktorny sistem postavaje ravny \(a_{S,T}e_1\cdot\bar b_{S,T}e_1\).}

\noindent\emph{Dokaz.}
\[
\begin{aligned}
e_1u_S\bar u_Se_1\cdot e_1u_T\bar u_Te_1
&=u_S\bar u_S\cdot u_T\bar u_Te_1
 =u_{ST}a_{S,T}\cdot\bar u_{ST}\bar b_{S,T}e_1\\
&=u_{ST}\bar u_{ST}\cdot a_{S,T}\bar b_{S,T}e_1
 =e_1u_{ST}\bar u_{ST}e_1\cdot a_{S,T}e_1\cdot\bar b_{S,T}e_1.
\end{aligned}
\]
Tym jest dokazan teorem iz 2., a zato i produktny teorem za faktorne sistemy.

\noindent\textbf{4. Produktny teorem za klasy asociovanyh faktornyh sistemov:}
\[
\{a_{S,T}\}\cdot\{b_{S,T}\}=\{a_{S,T}\cdot b_{S,T}\}.
\]
To prěmo slěduje iz 3., zato že prěhod k novym stvarjajučim elementam \(v\) i \(\bar v\) označaje takože taky prěhod k \(v\bar ve_1\), zato že \(u=vc\) i \(\bar u=\bar v\bar d\) davajut
\[
u\bar ue_1=vc\cdot\bar v\bar de_1=v\bar v\cdot c\bar de_1=v\bar ve_1\cdot c\bar de_1.
\]

\noindent\textbf{5.} Ako \(\{a\}\) jest faktorny sistem od \(\{\mathfrak K\}\), a \(t\) jest indeks od \(\{\mathfrak K\}\), to \(\{a\}^t\) jest ravny jedinčnomu klasu. To slěduje iz toga, že ireducibilno prěkriženo prědstavjenje od \(\mathfrak G\) imaje stupenj \(t\), upravo kako Šurov dokaz po prědstavjenju od \(p_{ik}\) (§ 24, str. 37).

\noindent\textbf{6.} Izomorfizm medžu grupoju \(\mathscr K_{\mathfrak Z}\) od \(\{\mathfrak K\}\) k \(\mathfrak Z\), t.~j. s \(\mathfrak Z\) kako razpadnym poljem, i grupoju klas \(\{a_{S,T}\}\) asociovanyh faktornyh sistemov. Po konstrukciji v § 27 obstaja jednoznačna korespondencija medžu \(\{\mathfrak K\}\) i \(\{a_{S,T}\}\); po 4. priręđenje jest homomorfno, zato imamo izomorfizm. Iz toga osoblivo slěduje, že jedinčne klasy odpovědajut jedne drugim:

\noindent\emph{Prěkriženy produkt jest togda i samo togda ravny \(P\), to jest matričnomu prstenu v centru, kogda faktorny sistem jest asociovany s jedinčnym sistemom.}
% END BOOK_S28

% BEGIN BOOK_S29
\subsection*{§ 29. Teorem o glavnom rodu v minimalnom}\label{hauptgeschlechtssatz-im-minimalen}

Nehaj \(\mathfrak K_r=\mathfrak Z\cdot\mathfrak G\) bude prěkriženy produkt, a \(\mathfrak G^*\) grupa vsih elementov, ktore posrědstvom vnutrnyh avtomorfizmov prěvodet \(\mathfrak Z\) kako cělost v sebe (\(\mathfrak Z^*u_E,\mathfrak Z^*u_S\), --- zapisano v pobočnyh klasah).

\noindent\textbf{Definicija.} „\emph{Glavny rod v minimalnom}“ \(\mathscr H\) sostoji se iz vsih grupovyh avtomorfizmov od \(\mathfrak G^*\), ktore sut prodolženje identiteta na \(\mathfrak Z^*\). Razlog za to nazvanje slěduje iz specializacije na ciklične polja. Vidi § 30.

\noindent\textbf{Teorem.} \emph{Vsaky avtomorfizm glavnogo roda imaje formu \(u_S\mapsto u_S\cdot b^Sb^{-1}=u_Sb^{S-1}\), kde \(b\) jest v \(\mathfrak Z\), i obratno, vsako tako priręđenje stvarjaje avtomorfizm glavnogo roda.}

\noindent\emph{Dokaz.} Nehaj pri grupovom avtomorfizmu bude \(u_S\mapsto v_S\), \(z\mapsto z\), za vsako \(z\) iz \(\mathfrak Z\), i nehaj pri tom produkty odpovědajut jedni drugim; togda avtomorfizm možno dopolniti do avtomorfizma prstena \(\mathfrak K_r\) posrědstvom priręđenja
\[
\sum u_Sc_{\lambda_S}\mapsto\sum v_Sc_{\lambda_S},
\]
pri čem \(c_{\lambda_S}\) sut v \(\mathfrak Z\). Ale vsaky taky avtomorfizm jest, kako znano, vnutrny.

Nehaj \(b\) bude stvarjajuči element, zato \(v_S=bu_Sb^{-1}\), \(z=bzb^{-1}\); zato že \(b\) jest komutativny so vsimi \(z\), on leži v \(\mathfrak Z\), i zato dobyva se
\[
v_S=bu_Sb^{-1}=u_Sb^Sb^{-1}.
\]
Obratno, ako položimo \(v_S=u_Sb^Sb^{-1}\) s \(b\) v \(\mathfrak Z\), to dobyva se \(v_S=bu_Sb^{-1}\); ide zato o avtomorfizm od \(\mathfrak K_r\), ktory prěvodi \(\mathfrak Z\) po elementah v sebe, a tym po definiciji prěvodi takože \(\mathfrak G^*\) kako cělost v sebe; on jest zato avtomorfizm glavnogo roda.
% END BOOK_S29

% BEGIN BOOK_S30
\subsection*{§ 30. Specializacija na ciklične razpadne polja}\label{spezialisierung-auf-zyklische-zerfaellungskoerper}

Ako prěkriženy produkt obrazujemo osoblivo s cikličnym poljem \(\mathfrak Z\), prihodimo k tolkovanju znanyh normnyh teoremov. Najprvo imaje město

\noindent\textbf{Teorem 1.} \emph{Grupa \(\mathscr K_{\mathfrak Z}\) od \(\{\mathfrak K\}\) s ustanovjenym cikličnym razpadnym poljem \(\mathfrak Z\) nad \(P\) jest izomorfna s grupoju \(P^*/N\mathfrak Z^*\), kde \(N\mathfrak Z^*\) proběgaje grupu vsih norm \(Nz\) elementov \(z\neq0\) iz \(\mathfrak Z\).}

\noindent\emph{Dokaz.} Dokaz osnovyvaje se na normovanju prěkriženogo prědstavjenja i na několiko pomoćnyh lemah o tom.

\noindent\textbf{Definicija.} Ciklično prěkriženo prědstavjenje nazyvaje se \emph{normovanym}, ako imaje formu
\[
\mathfrak K_r=\mathfrak Z+\mathfrak Zu+\cdots+\mathfrak Zu^{n-1},
\]
s \(zu=uz^S\), \(u^n=a\), kde \(S\) jest stvarjajuči element iz \(\mathfrak G\) (ako zato položimo \(u_{S^r}=u^r\), \(r=0,\ldots,n-1\), vměsto obćego \(u_{S^r}=b_ru^r\) s \(b_r\) v \(\mathfrak Z\)).

\noindent\textbf{Pomoćna lema 1.} \emph{Ciklično prěkriženo prědstavjenje možno vsegda vzęti normovanym. Pri tom \(a\) leži v \(P\), i vsako \(a\) iz \(P\) davaje tako prědstavjenje.}

\noindent\emph{Dokaz.} Prva čest jest jasna, zato že iz \(u^{-1}zu=z^S\) slěduje \(u^{-r}zu^r=z^{S^r}\), zato \(u^r\) stvarjaje substituciju \(z\mapsto z^{S^r}\). Dalje, zato že \(u^n\) induciraje identitet \(z\mapsto z^{S^n}\), dobyva se \(u^n=a\) s \(a\) v \(\mathfrak Z\). Ale \(u^n\) jest komutativny so vsimi \(u^r\), to samo imaje město za \(a\), i zato \(a^{S^r}=a\); zato \(a\) leži v \(P\).

zato \(a\) jest jedina sućstvena konstanta množenja. Zato uslovja asociativnosti ne imajut sodržanje; vsako \(a\) iz \(P\) davaje asociativny sistem, a tym i neko \(\mathfrak K_r\).

\noindent\textbf{Pomoćna lema 2.} \emph{Vměsto polnogo klasa \(\{a_{S,T}\}\) asociovanyh faktornyh sistemov dostatočno jest razgledati normovane faktorne sistemy, ktore sut sodrživane v klasu --- t.~j. ktore nastavajut iz normovanyh prědstavjenij. Tu podmnožinu klasa nazyvamo normovanym klasom. Grupa polnyh klas \(\{a_{S,T}\}\) jest izomorfna s grupoju normovanyh klas.}

\noindent\emph{Dokaz.} Zato že po pomoćnoj lemi 1 nijedin normovany klas ne jest prazdny, priręđenje jest zato jednoznačno. Operacija ostavaje ista, zato že jest oprěděljena posrědstvom libovoljnogo prědstavitelja.

\noindent\textbf{Pomoćna lema 3.} \emph{Grupa normovanyh faktornyh sistemov jest izomorfna s multiplikativnoj grupoju \(P^*\), a grupa normovanyh klas jest izomorfna s \(P^*/N\mathfrak Z^*\).}

\noindent\emph{Dokaz.} Iz pomoćnoj lemy 1 slěduje, že vsaky normovany faktorny sistem imaje vid \([1,1,\ldots,a,a,\ldots]\), zato že
\[
u^\nu\cdot u^\mu=u^{\nu+\mu}\quad(\nu+\mu<n),
\qquad
u^\nu\cdot u^\mu=au^\varrho\quad(\nu+\mu=n+\varrho,\ 0\leq\varrho\leq n-2).
\]
Ale priręđenje
\[
[1,\ldots,a,a,\ldots]\mapsto a,
\qquad
[1,1,\ldots,b,b,\ldots]\mapsto b,
\]
\[
[1,\ldots,ab,\ldots]\mapsto ab,
\]
jest grupovy izomorfizm na \(P^*\), zato že, ponovno po pomoćnoj lemi 1, \(a\) proběgaje vse elementy iz \(P^*\). --- Dva asociovane normovane faktorne sistemy nastavajut jeden iz drugogo posrědstvom transformacije \(v=uc\) (\(c\) v \(\mathfrak Z\)). To ale davaje
\[
v=uc,\qquad v^2=u^2\cdot c^S\cdot c;
\]
\[
v^r=u^r c^{S^{r-1}}c^{S^{r-2}}\cdots c;
\qquad
v^n=u^nN(c).
\]
Iz \(u^n=a\) zato slěduje \(v^n=aN(c)\). Normovany klas sostoji se zato iz faktornyh sistemov \([1,\ldots,aN(c),aN(c),\ldots]\), kde \(c\) proběgaje vse elementy \(\neq0\) iz \(\mathfrak Z\). Tym jest dokazana druga čest pomoćnoj lemy.

Iz pomoćnyh lem 2 i 3 slěduje gornji teorem 1, s obzirom na izomorfizm tamošnjih \(\{\mathfrak K\}\) s klasami \(\{a_{S,T}\}\).

\noindent\textbf{Teorem 2.} \emph{Grupa glavnogo roda v minimalnom jest izomorfna s grupoju tyh elementov \(c\) iz \(\mathfrak Z\), za ktore \(N(c)=1\). Iz \(N(c)=1\) zato slěduje obstajanje nekogo \(b\neq0\) iz \(\mathfrak Z\), tako že \(c=b^S\cdot b^{-1}=b^{S-1}\) (simbolična \(S-1\)-ta potencija).}

\noindent\emph{Dokaz.} Pri dokazu pomoćnoj lemy 3 dostalo se: \(N(c)=1\) jest identično s tym, že transformacija \(v=uc\) prědstavja avtomorfizm glavnogo roda; različnym \(c\) odpovědajut različne avtomorfizmy i obratno, a produktu odpovědaje produkt.

\noindent\textbf{Primětka.} Teorem 2 opravdava nazvanje glavny rod. Po smyslu faktorgrupu
\[
\mathfrak Z^*/\mathscr H
\quad(\mathscr H=\text{glavny rod})
\cong\mathfrak Z^*/\mathfrak Z^{*S-1}
\]
možno nazvati grupoju rodov v minimalnom; za tu grupu imaje město, že ona jest izomorfna s \(N(\mathfrak Z^*)\), zato v svezi s teoremom 1,
\[
\begin{aligned}
\text{grupa rodov v minimalnom}
&=\mathfrak Z^*/\mathscr H
 \cong\mathfrak Z^*/\mathfrak Z^{*S-1}
 \cong N(\mathfrak Z^*)\\
&=\text{grupa od }\{\mathfrak K\}\text{ k }\mathfrak Z
 =\mathscr K_{\mathfrak Z}
 \cong P^*/N(\mathfrak Z^*).
\end{aligned}
\]
% END BOOK_S30

% BEGIN BOOK_S31
\subsection*{§ 31. Priměnjenja cikličnogo specialnogo slučaja}\label{anwendungen-des-zyklischen-spezialfalles}

Ciklično prěkriženo prědstavjenje davaje nove dokazy za teoremy dokazane v § 20: že vsako konečno tělo jest komutativno i že kromě kvaternionov ne obstaja nikako pravo nekomutativno razširjenje polja realnyh čisel. --- Te dokazy možno by bylo nazvati dokazami teorije klasovyh polj v minimalnom.

\noindent\textbf{1. Vsako tělo iz konečnogo čisla elementov jest komutativno (teorem 10 § 20).}

Dokaz osnovyvaje se na dvoh znanyh faktah o komutativnyh konečnyh poljah (Galoisovyh poljah).

\noindent\textbf{1.} \emph{Galoisove polja sut ciklične nad vsakym svojim podpoljem.}

\noindent\textbf{2.} \emph{Ako \(\mathfrak Z\) jest Galoisovo polje, to \(\mathfrak Z\) jest ciklično nad \(P\), to vsaky element iz \(P\) jest norma nekogo elementa od \(\mathfrak Z\). \(P^*=N(\mathfrak Z^*)\).}

Nehaj nyně \(\mathfrak K\) bude konečno, \(P\) jego centar, a \(\mathfrak Z\) neko ciklično razpadno polje; togda (§ 30 konec), medžu drugym, grupa \(\mathscr K_{\mathfrak Z}\) od \(\{\mathfrak K\}\) k \(\mathfrak Z\) jest izomorfna s \(P^*/N\mathfrak Z^*\); zato po 2. ona jest ravna jedinici \(\{P\}\). To ale znači, že ne obstaja od \(P\) različno tělo s \(P\) kako centrom, \(\mathfrak K=P\).

\noindent\textbf{2. Jedino pravo nekomutativno razširjenje nad poljem \(P\) realnyh čisel sut kvaterniony (teorem 9 § 20).}

\noindent\emph{Dokaz.} \(P\) mora byti centar, zato že \(P\) imaje kako komutativno razširjenje samo polje \(\mathfrak Z\) kompleksnyh čisel, nad ktorym kako centrom, zato že \(\mathfrak Z\) jest algebraično zatvorjeno, ne obstaja nekomutativno razširjenje těla. Jednočasno samo \(\mathfrak Z\) prihodi v obzir kako razpadno polje. I \(\mathfrak Z\) jest ciklično nad \(P\). Dalje, vsako pozitivno čislo jest norma, zato grupa \(P^*/N\mathfrak Z^*\) imaje stupenj 2; zato kromě \(\{P\}\) obstaja samo jeden klas \(\{\mathfrak K\}\), a to sut kvaterniony. Jednočasno dobyva se normalna forma, ako vzmemo \(-1\) kako normovany faktorny sistem v klasu asociovanyh sistemov. Zato že togda dobyva se \(\mathfrak K=\mathfrak Z+\mathfrak Zu\), s \(u^2=-1\). zato \(\mathfrak K=P+Pi+(P+Pi)u\) s \(i^2=-1\), \(u^2=-1\).

\noindent\textbf{Primětka k 1.} Takože dokaz znanogo fakta 2 možno na osnově konca § 30 voditi tako: Dobyva se \(\mathscr H=\mathfrak Z^{*S-1}\cong\mathfrak Z^*/P^*\), zato že vse i samo elementi od \(P^*\) stvarjajut identičny avtomorfizm glavnogo roda, ili takože pri priměnjenju od \(S-1\) prěhodet v jedinicu. zato pri konečnom \(\mathfrak Z\) čislo elementov od \(\mathscr H\) jest ravno indeksu od \(P^*\) v \(\mathfrak Z^*\). Zato čislo elementov od \(\mathfrak Z^*/\mathscr H\) jest ravno čislu elementov od \(P^*\); to samo zato imaje město za \(N(\mathfrak Z^*)=\mathfrak Z^*/\mathscr H\), i zato \(P^*=N(\mathfrak Z^*)\). (Sravni: čislo rodov = čislo ambigvnyh klas. Tomu čislovo-teoretičnomu načinu izražanja ale vsude trěba dodati „v minimalnom“.)

\noindent\textbf{3.} \emph{Ako \(P\) jest \(p\)-adično osnovno polje, a \(\mathfrak Z\) jest ciklično \(n\)-togo stupenja nad \(P\), to grupa \(\mathscr K_{\mathfrak Z}\) od \(\{\mathfrak K\}\) k \(\mathfrak Z\) jest izomorfna s Galoisovoj grupoju od \(\mathfrak Z\).}

To jest po koncu § 30 prěma poslědica „teorije klasovyh polj v malom“, ktora izreka
\[
P^*/N\mathfrak Z^*\cong\mathfrak G.
\]

\noindent\textbf{Primětka k 3.} Slěduje, že grupovy ręd elementov iz \(\mathscr K_{\mathfrak Z}\), to jest eksponent od \(\{\mathfrak K\}\), postavaje ravny jego indeksu. Zato že eksponent stvarjajučego elementa od \(\mathscr K_{\mathfrak Z}\) postavaje \(=n\), zato takože indeks kako dělitelj od \(n\) i mnogokratnik eksponenta; grupa \(\mathscr K_{\mathfrak Z}\) sodrživaje zato samo klasy, ktoryh indeks jest dělitelj od \(n\), i za ktore \(\mathfrak Z\) jest razpadno polje. Nadto H. Hasse pokazal, Math. Ann. 104, že \(\mathscr K_{\mathfrak Z}\) sostoji se iz vsih klas \(\{\mathfrak K\}\), za ktore indeks jest dělitelj od \(n\). Vsako ciklično razpadno polje \(\mathfrak Z\) \(n\)-togo stupenja nad \(p\)-adičnym \(P\) vodi zato k toj samoj grupě \(\mathscr K_{\mathfrak Z}\) klas těl \(\{\mathfrak K\}\).

Dva fakta:

\noindent\textbf{1.} grupa vsih \(\{\mathfrak K\}\) nad \(P\), ktoryh indeks jest dělitelj od \(n\), jest ciklična \(n\)-togo stupenja.

\noindent\textbf{2.} Vsako ciklično razpadno polje \(n\)-togo stupenja nad \(P\) stvarjaje polnu grupu od \(\{\mathfrak K\}\).

možno razgledati sućstvenym sodržanjem obratnogo teorema teorije klasovyh polj v malom.


\clearpage
% END BOOK_S31

\end{document}
